% !TEX program = lualatex
% ---------------------------------------------------------------------------------------------
%  Readback note 33: A fit family of twists exists: Ribet's method with tied primes, triangular Taylor–Wiles primes and the denominator of the projection to theta lifts
%
%  Build:  python3 tools/rb.py note <slug>
%
%  One unit of a Lean artifact, rewritten as ordinary mathematics. The prose is hand-written. Lean names,
%  line ranges and counts come from data.tex; quoted Lean statements are copied from the artifact by the build.
%
%  Rules of the text: standard terminology; every symbol is explained where it first appears and listed again
%  in the notation at the end; no letter has two meanings; a proof is complete or says what it takes for granted;
%  Lean names appear only in the LEAN lines and in the concordance.
%
%  The signature rule (the build enforces it): a function, relation or predicate that the artifact defines is
%  written only through its notation macro, which the build makes from the spec (`notation`) with one slot for
%  each explicit argument of the Lean definition. No argument is fixed and dropped, also not "for brevity" and
%  not for one section. Every argument and every hypothesis of a stated declaration stands inside its
%  statement, or in a Setting that the statement cites by number; only implicit and instance arguments (the
%  ambient objects) may stand in the conventions.
% ---------------------------------------------------------------------------------------------
\documentclass[11pt,a4paper]{article}
\ifdefined\pdfvariable\pdfvariable suppressoptionalinfo \numexpr 1+2+4+8\relax\fi
\usepackage[a4paper,hmargin=35mm,top=22mm,bottom=27mm,footskip=12mm,nohead]{geometry}

\usepackage{fontspec}
\directlua{
  luaotfload.add_fallback("gfallback", {"STIXTwoMath-Regular.otf:mode=node;", "DejaVuSans.ttf:mode=node;"})
  luaotfload.add_fallback("gmonofallback", {"DejaVuSansMono.ttf:mode=node;", "STIXTwoMath-Regular.otf:mode=node;"})
}
\setmainfont{NewCM10}[Extension=.otf,UprightFont=*-Book,ItalicFont=*-BookItalic,
  BoldFont=*-Bold,BoldItalicFont=*-BoldItalic,SmallCapsFont=*-Book,SmallCapsFeatures={Letters=SmallCaps},
  RawFeature={fallback=gfallback}]
\setsansfont{NewCMSans10}[Extension=.otf,UprightFont=*-Book,ItalicFont=*-BookOblique,
  BoldFont=*-Bold,BoldItalicFont=*-BoldOblique,RawFeature={fallback=gfallback}]
\setmonofont{JetBrainsMono}[Extension=.otf,UprightFont=*-Regular,ItalicFont=*-Italic,
  BoldFont=*-SemiBold,BoldItalicFont=*-SemiBoldItalic,Scale=MatchLowercase,
  RawFeature={-calt;-liga;fallback=gmonofallback}]
\usepackage{amsmath}
\usepackage{amsthm}
\usepackage[math-style=TeX]{unicode-math}
\setmathfont{NewCMMath-Book.otf}
\emergencystretch=2.5em\relax

\usepackage{xcolor}
\usepackage{tikz}
\usetikzlibrary{calc,arrows.meta,positioning,backgrounds}
\usepackage[protrusion=true,expansion=true]{microtype}
\usepackage{enumitem}
\usepackage{needspace}
\usepackage{booktabs}
\usepackage{array}
\usepackage{longtable}
\usepackage{fancyhdr}
\usepackage{lastpage}
\usepackage{multicol}
\usepackage[hidelinks,pdfusetitle]{hyperref}

\definecolor{ink}{HTML}{0F172A}
\definecolor{inksoft}{HTML}{475569}
\definecolor{inkfaint}{HTML}{94A3B8}
\definecolor{rulegrey}{HTML}{CBD5E1}
\definecolor{arith}{HTML}{C2410C}
\definecolor{est}{HTML}{7C3AED}
\definecolor{ana}{HTML}{0891B2}
\definecolor{ok}{HTML}{047857}
\definecolor{leancol}{HTML}{3730A3}

\title{A fit family of twists exists: Ribet's method with tied primes, triangular Taylor–Wiles primes and the denominator of the projection to theta lifts}
\author{grh-compact}
\date{}

% ------------------------------------------------------------------------------------- macros
\newcommand{\leanfont}{\ttfamily\fontsize{8.6}{10.5}\selectfont\color{leancol}}
\newcommand{\lean}[1]{\csname LN@#1\endcsname}
\newcommand{\fillin}[1]{{\color{arith}\itshape[#1]}}  % a place of the template that is not written yet; the build reports every one that is left
\newcommand{\leanquote}[1]{\csname LQ@#1\endcsname}
% generated by rb.py note; do not edit
\newcommand{\ConcDecls}{92}
\newcommand{\Pin}{72fead10}
\newcommand{\PinText}{that of commit {\leanfont 72fead10}}
\newcommand{\ArtifactName}{grh-compact}
\expandafter\newcommand\csname LN@IsConjSelfDual\endcsname{\href{https://github.com/lsernau/grh-lean/blob/72fead102f663d5c7a407964a13dc0cf9d977207/grh-compact/src/New/Selmer/AptTupleDefs.lean\#L21-L24}{{\leanfont IsConjSelfDual}}}
\expandafter\newcommand\csname LN@crudeUnit\endcsname{\href{https://github.com/lsernau/grh-lean/blob/72fead102f663d5c7a407964a13dc0cf9d977207/grh-compact/src/New/Selmer/AptTupleDefs.lean\#L27-L27}{{\leanfont crudeUnit}}}
\expandafter\newcommand\csname LN@SelmerLengthLe\endcsname{\href{https://github.com/lsernau/grh-lean/blob/72fead102f663d5c7a407964a13dc0cf9d977207/grh-compact/src/New/Selmer/AptTupleDefs.lean\#L34-L38}{{\leanfont SelmerLengthLe}}}
\expandafter\newcommand\csname LN@IsColumn\endcsname{\href{https://github.com/lsernau/grh-lean/blob/72fead102f663d5c7a407964a13dc0cf9d977207/grh-compact/src/New/Selmer/AptTupleDefs.lean\#L43-L45}{{\leanfont IsColumn}}}
\expandafter\newcommand\csname LN@IsRow\endcsname{\href{https://github.com/lsernau/grh-lean/blob/72fead102f663d5c7a407964a13dc0cf9d977207/grh-compact/src/New/Selmer/AptTupleDefs.lean\#L50-L59}{{\leanfont IsRow}}}
\expandafter\newcommand\csname LN@chiOneInv\endcsname{\href{https://github.com/lsernau/grh-lean/blob/72fead102f663d5c7a407964a13dc0cf9d977207/grh-compact/src/New/Selmer/AptTupleDefs.lean\#L68-L68}{{\leanfont chiOneInv}}}
\expandafter\newcommand\csname LN@normalisedRep\endcsname{\href{https://github.com/lsernau/grh-lean/blob/72fead102f663d5c7a407964a13dc0cf9d977207/grh-compact/src/New/Selmer/AptTupleDefs.lean\#L82-L82}{{\leanfont normalisedRep}}}
\expandafter\newcommand\csname LN@twistedLabel\endcsname{\href{https://github.com/lsernau/grh-lean/blob/72fead102f663d5c7a407964a13dc0cf9d977207/grh-compact/src/New/Selmer/AptTupleDefs.lean\#L86-L87}{{\leanfont twistedLabel}}}
\expandafter\newcommand\csname LN@IsApt\endcsname{\href{https://github.com/lsernau/grh-lean/blob/72fead102f663d5c7a407964a13dc0cf9d977207/grh-compact/src/New/Selmer/AptTupleDefs.lean\#L92-L100}{{\leanfont IsApt}}}
\expandafter\newcommand\csname LN@tupleProj\endcsname{\href{https://github.com/lsernau/grh-lean/blob/72fead102f663d5c7a407964a13dc0cf9d977207/grh-compact/src/New/Selmer/AptTupleDefs.lean\#L104-L104}{{\leanfont tupleProj}}}
\expandafter\newcommand\csname LN@HasRow\endcsname{\href{https://github.com/lsernau/grh-lean/blob/72fead102f663d5c7a407964a13dc0cf9d977207/grh-compact/src/New/Selmer/AptTupleDefs.lean\#L108-L110}{{\leanfont HasRow}}}
\expandafter\newcommand\csname LN@IsLabelledColumn\endcsname{\href{https://github.com/lsernau/grh-lean/blob/72fead102f663d5c7a407964a13dc0cf9d977207/grh-compact/src/New/Selmer/AptTupleDefs.lean\#L114-L116}{{\leanfont IsLabelledColumn}}}
\expandafter\newcommand\csname LN@obeyingCocycles\endcsname{\href{https://github.com/lsernau/grh-lean/blob/72fead102f663d5c7a407964a13dc0cf9d977207/grh-compact/src/New/Selmer/BlocksAndCocycles.lean\#L16-L17}{{\leanfont obeyingCocycles}}}
\expandafter\newcommand\csname LN@RowBlockIndex\endcsname{\href{https://github.com/lsernau/grh-lean/blob/72fead102f663d5c7a407964a13dc0cf9d977207/grh-compact/src/New/Selmer/BlocksAndCocycles.lean\#L41-L41}{{\leanfont RowBlockIndex}}}
\expandafter\newcommand\csname LN@blockRep\endcsname{\href{https://github.com/lsernau/grh-lean/blob/72fead102f663d5c7a407964a13dc0cf9d977207/grh-compact/src/New/Selmer/BlocksAndCocycles.lean\#L50-L51}{{\leanfont blockRep}}}
\expandafter\newcommand\csname LN@IsAuxiliary\endcsname{\href{https://github.com/lsernau/grh-lean/blob/72fead102f663d5c7a407964a13dc0cf9d977207/grh-compact/src/New/Selmer/Corepresentable.lean\#L55-L55}{{\leanfont IsAuxiliary}}}
\expandafter\newcommand\csname LN@evenClasses\endcsname{\href{https://github.com/lsernau/grh-lean/blob/72fead102f663d5c7a407964a13dc0cf9d977207/grh-compact/src/New/Selmer/Corepresentable.lean\#L136-L138}{{\leanfont evenClasses}}}
\expandafter\newcommand\csname LN@IsNewPlaces\endcsname{\href{https://github.com/lsernau/grh-lean/blob/72fead102f663d5c7a407964a13dc0cf9d977207/grh-compact/src/New/Selmer/Corepresentable.lean\#L142-L143}{{\leanfont IsNewPlaces}}}
\expandafter\newcommand\csname LN@withConjInv\endcsname{\href{https://github.com/lsernau/grh-lean/blob/72fead102f663d5c7a407964a13dc0cf9d977207/grh-compact/src/New/Selmer/Corepresentable.lean\#L147-L147}{{\leanfont withConjInv}}}
\expandafter\newcommand\csname LN@IsPrimeTriangle\endcsname{\href{https://github.com/lsernau/grh-lean/blob/72fead102f663d5c7a407964a13dc0cf9d977207/grh-compact/src/New/Selmer/Corepresentable.lean\#L152-L162}{{\leanfont IsPrimeTriangle}}}
\expandafter\newcommand\csname LN@IsAuxTwist\endcsname{\href{https://github.com/lsernau/grh-lean/blob/72fead102f663d5c7a407964a13dc0cf9d977207/grh-compact/src/New/Selmer/Corepresentable.lean\#L166-L167}{{\leanfont IsAuxTwist}}}
\expandafter\newcommand\csname LN@existscycCharpowsubonedvd\endcsname{\href{https://github.com/lsernau/grh-lean/blob/72fead102f663d5c7a407964a13dc0cf9d977207/grh-compact/src/New/Selmer/DenominatorBound.lean\#L26-L40}{{\leanfont exists\_cycChar\_pow\_sub\_one\_dvd}}}
\expandafter\newcommand\csname LN@denominatorLerowBudget\endcsname{\href{https://github.com/lsernau/grh-lean/blob/72fead102f663d5c7a407964a13dc0cf9d977207/grh-compact/src/New/Selmer/DenominatorBound.lean\#L47-L166}{{\leanfont denominatorLe\_rowBudget}}}
\expandafter\newcommand\csname LN@crysRels\endcsname{\href{https://github.com/lsernau/grh-lean/blob/72fead102f663d5c7a407964a13dc0cf9d977207/grh-compact/src/New/Selmer/DeterminantBound.lean\#L43-L43}{{\leanfont crysRels}}}
\expandafter\newcommand\csname LN@blockLoss\endcsname{\href{https://github.com/lsernau/grh-lean/blob/72fead102f663d5c7a407964a13dc0cf9d977207/grh-compact/src/New/Selmer/DeterminantBound.lean\#L47-L48}{{\leanfont blockLoss}}}
\expandafter\newcommand\csname LN@cmSign\endcsname{\href{https://github.com/lsernau/grh-lean/blob/72fead102f663d5c7a407964a13dc0cf9d977207/grh-compact/src/New/Selmer/DeterminantBound.lean\#L56-L56}{{\leanfont cmSign}}}
\expandafter\newcommand\csname LN@IsRowType\endcsname{\href{https://github.com/lsernau/grh-lean/blob/72fead102f663d5c7a407964a13dc0cf9d977207/grh-compact/src/New/Selmer/DeterminantBound.lean\#L64-L70}{{\leanfont IsRowType}}}
\expandafter\newcommand\csname LN@lambdaLoss\endcsname{\href{https://github.com/lsernau/grh-lean/blob/72fead102f663d5c7a407964a13dc0cf9d977207/grh-compact/src/New/Selmer/DeterminantBound.lean\#L74-L75}{{\leanfont lambdaLoss}}}
\expandafter\newcommand\csname LN@loss\endcsname{\href{https://github.com/lsernau/grh-lean/blob/72fead102f663d5c7a407964a13dc0cf9d977207/grh-compact/src/New/Selmer/DeterminantBound.lean\#L78-L79}{{\leanfont loss}}}
\expandafter\newcommand\csname LN@innerAtheckePolyOnForms\endcsname{\href{https://github.com/lsernau/grh-lean/blob/72fead102f663d5c7a407964a13dc0cf9d977207/grh-compact/src/New/Selmer/FitTwists.lean\#L19-L24}{{\leanfont LevelPlaces.\allowbreak{}LevelForm.\allowbreak{}innerAt\_heckePolyOnForms}}}
\expandafter\newcommand\csname LN@groupAlgAction\endcsname{\href{https://github.com/lsernau/grh-lean/blob/72fead102f663d5c7a407964a13dc0cf9d977207/grh-compact/src/New/Selmer/FitTwists.lean\#L35-L35}{{\leanfont AuxPair.\allowbreak{}groupAlgAction}}}
\expandafter\newcommand\csname LN@IsTested\endcsname{\href{https://github.com/lsernau/grh-lean/blob/72fead102f663d5c7a407964a13dc0cf9d977207/grh-compact/src/New/Selmer/FitTwists.lean\#L39-L40}{{\leanfont AuxPair.\allowbreak{}IsTested}}}
\expandafter\newcommand\csname LN@IsTestForm\endcsname{\href{https://github.com/lsernau/grh-lean/blob/72fead102f663d5c7a407964a13dc0cf9d977207/grh-compact/src/New/Selmer/FitTwists.lean\#L44-L45}{{\leanfont AuxPair.\allowbreak{}IsTestForm}}}
\expandafter\newcommand\csname LN@DenominatorLe\endcsname{\href{https://github.com/lsernau/grh-lean/blob/72fead102f663d5c7a407964a13dc0cf9d977207/grh-compact/src/New/Selmer/FitTwists.lean\#L50-L53}{{\leanfont DenominatorLe}}}
\expandafter\newcommand\csname LN@isOfFinOrderprodpsidet\endcsname{\href{https://github.com/lsernau/grh-lean/blob/72fead102f663d5c7a407964a13dc0cf9d977207/grh-compact/src/New/Selmer/FitTwists.lean\#L59-L61}{{\leanfont AuxPair.\allowbreak{}isOfFinOrder\_prod\_psi\_det}}}
\expandafter\newcommand\csname LN@innerAtformAction\endcsname{\href{https://github.com/lsernau/grh-lean/blob/72fead102f663d5c7a407964a13dc0cf9d977207/grh-compact/src/New/Selmer/FitTwists.lean\#L66-L70}{{\leanfont AuxPair.\allowbreak{}innerAt\_formAction}}}
\expandafter\newcommand\csname LN@rangegroupAlgActionmeminvtSubmodule\endcsname{\href{https://github.com/lsernau/grh-lean/blob/72fead102f663d5c7a407964a13dc0cf9d977207/grh-compact/src/New/Selmer/FitTwists.lean\#L72-L74}{{\leanfont AuxPair.\allowbreak{}range\_groupAlgAction\_mem\_invtSubmodule}}}
\expandafter\newcommand\csname LN@heckePolyOnFormscomm\endcsname{\href{https://github.com/lsernau/grh-lean/blob/72fead102f663d5c7a407964a13dc0cf9d977207/grh-compact/src/New/Selmer/FitTwists.lean\#L78-L82}{{\leanfont AuxPair.\allowbreak{}heckePolyOnForms\_comm}}}
\expandafter\newcommand\csname LN@innerAtgroupAlgAction\endcsname{\href{https://github.com/lsernau/grh-lean/blob/72fead102f663d5c7a407964a13dc0cf9d977207/grh-compact/src/New/Selmer/FitTwists.lean\#L86-L88}{{\leanfont AuxPair.\allowbreak{}innerAt\_groupAlgAction}}}
\expandafter\newcommand\csname LN@normSqAtstarProjectionle\endcsname{\href{https://github.com/lsernau/grh-lean/blob/72fead102f663d5c7a407964a13dc0cf9d977207/grh-compact/src/New/Selmer/FitTwists.lean\#L96-L112}{{\leanfont AuxPair.\allowbreak{}normSqAt\_starProjection\_le}}}
\expandafter\newcommand\csname LN@IsSplitOffLevel\endcsname{\href{https://github.com/lsernau/grh-lean/blob/72fead102f663d5c7a407964a13dc0cf9d977207/grh-compact/src/New/Selmer/FitTwists.lean\#L117-L117}{{\leanfont Data.\allowbreak{}IsSplitOffLevel}}}
\expandafter\newcommand\csname LN@heckeAlgebra\endcsname{\href{https://github.com/lsernau/grh-lean/blob/72fead102f663d5c7a407964a13dc0cf9d977207/grh-compact/src/New/Selmer/FitTwists.lean\#L122-L127}{{\leanfont Data.\allowbreak{}heckeAlgebra}}}
\expandafter\newcommand\csname LN@heckeOperator\endcsname{\href{https://github.com/lsernau/grh-lean/blob/72fead102f663d5c7a407964a13dc0cf9d977207/grh-compact/src/New/Selmer/FitTwists.lean\#L131-L136}{{\leanfont Data.\allowbreak{}heckeOperator}}}
\expandafter\newcommand\csname LN@eigenPart\endcsname{\href{https://github.com/lsernau/grh-lean/blob/72fead102f663d5c7a407964a13dc0cf9d977207/grh-compact/src/New/Selmer/FitTwists.lean\#L140-L140}{{\leanfont Data.\allowbreak{}eigenPart}}}
\expandafter\newcommand\csname LN@isSphericalWithofmemeigenPart\endcsname{\href{https://github.com/lsernau/grh-lean/blob/72fead102f663d5c7a407964a13dc0cf9d977207/grh-compact/src/New/Selmer/FitTwists.lean\#L145-L148}{{\leanfont Data.\allowbreak{}isSphericalWith\_of\_mem\_eigenPart}}}
\expandafter\newcommand\csname LN@occurringLabels\endcsname{\href{https://github.com/lsernau/grh-lean/blob/72fead102f663d5c7a407964a13dc0cf9d977207/grh-compact/src/New/Selmer/FitTwists.lean\#L152-L153}{{\leanfont Data.\allowbreak{}occurringLabels}}}
\expandafter\newcommand\csname LN@dvdcardresFieldsubone\endcsname{\href{https://github.com/lsernau/grh-lean/blob/72fead102f663d5c7a407964a13dc0cf9d977207/grh-compact/src/New/Selmer/FitTwists.lean\#L155-L158}{{\leanfont Data.\allowbreak{}dvd\_card\_resField\_sub\_one}}}
\expandafter\newcommand\csname LN@existsheckeprojeigenPart\endcsname{\href{https://github.com/lsernau/grh-lean/blob/72fead102f663d5c7a407964a13dc0cf9d977207/grh-compact/src/New/Selmer/FitTwists.lean\#L162-L167}{{\leanfont Data.\allowbreak{}exists\_hecke\_proj\_eigenPart}}}
\expandafter\newcommand\csname LN@commuteheckeAlgebraprincipalRep\endcsname{\href{https://github.com/lsernau/grh-lean/blob/72fead102f663d5c7a407964a13dc0cf9d977207/grh-compact/src/New/Selmer/FitTwists.lean\#L169-L173}{{\leanfont Data.\allowbreak{}commute\_heckeAlgebra\_principalRep}}}
\expandafter\newcommand\csname LN@heckeOnGood\endcsname{\href{https://github.com/lsernau/grh-lean/blob/72fead102f663d5c7a407964a13dc0cf9d977207/grh-compact/src/New/Selmer/FitTwists.lean\#L177-L179}{{\leanfont Data.\allowbreak{}heckeOnGood}}}
\expandafter\newcommand\csname LN@principalRepmemeigenPart\endcsname{\href{https://github.com/lsernau/grh-lean/blob/72fead102f663d5c7a407964a13dc0cf9d977207/grh-compact/src/New/Selmer/FitTwists.lean\#L181-L183}{{\leanfont Data.\allowbreak{}principalRep\_mem\_eigenPart}}}
\expandafter\newcommand\csname LN@finiteeigenPart\endcsname{\href{https://github.com/lsernau/grh-lean/blob/72fead102f663d5c7a407964a13dc0cf9d977207/grh-compact/src/New/Selmer/FitTwists.lean\#L187-L188}{{\leanfont Data.\allowbreak{}finite\_eigenPart}}}
\expandafter\newcommand\csname LN@iSupeigenPart\endcsname{\href{https://github.com/lsernau/grh-lean/blob/72fead102f663d5c7a407964a13dc0cf9d977207/grh-compact/src/New/Selmer/FitTwists.lean\#L192-L197}{{\leanfont Data.\allowbreak{}iSup\_eigenPart}}}
\expandafter\newcommand\csname LN@thetamuleqthetaformAction\endcsname{\href{https://github.com/lsernau/grh-lean/blob/72fead102f663d5c7a407964a13dc0cf9d977207/grh-compact/src/New/Selmer/FitTwists.lean\#L205-L212}{{\leanfont AuxPair.\allowbreak{}theta\_mul\_eq\_theta\_formAction}}}
\expandafter\newcommand\csname LN@existsthetatoprincipalForms\endcsname{\href{https://github.com/lsernau/grh-lean/blob/72fead102f663d5c7a407964a13dc0cf9d977207/grh-compact/src/New/Selmer/FitTwists.lean\#L216-L222}{{\leanfont Data.\allowbreak{}exists\_theta\_to\_principalForms}}}
\expandafter\newcommand\csname LN@auxLift\endcsname{\href{https://github.com/lsernau/grh-lean/blob/72fead102f663d5c7a407964a13dc0cf9d977207/grh-compact/src/New/Selmer/FitTwists.lean\#L226-L227}{{\leanfont Data.\allowbreak{}auxLift}}}
\expandafter\newcommand\csname LN@RankOneIndex\endcsname{\href{https://github.com/lsernau/grh-lean/blob/72fead102f663d5c7a407964a13dc0cf9d977207/grh-compact/src/New/Selmer/FitTwists.lean\#L231-L231}{{\leanfont Data.\allowbreak{}RankOneIndex}}}
\expandafter\newcommand\csname LN@goodLift\endcsname{\href{https://github.com/lsernau/grh-lean/blob/72fead102f663d5c7a407964a13dc0cf9d977207/grh-compact/src/New/Selmer/FitTwists.lean\#L237-L237}{{\leanfont Data.\allowbreak{}goodLift}}}
\expandafter\newcommand\csname LN@IsRamified\endcsname{\href{https://github.com/lsernau/grh-lean/blob/72fead102f663d5c7a407964a13dc0cf9d977207/grh-compact/src/New/Selmer/FitTwists.lean\#L241-L241}{{\leanfont Data.\allowbreak{}Pair.\allowbreak{}IsRamified}}}
\expandafter\newcommand\csname LN@IsRankOne\endcsname{\href{https://github.com/lsernau/grh-lean/blob/72fead102f663d5c7a407964a13dc0cf9d977207/grh-compact/src/New/Selmer/FitTwists.lean\#L245-L247}{{\leanfont Data.\allowbreak{}Pair.\allowbreak{}IsRankOne}}}
\expandafter\newcommand\csname LN@B\endcsname{\href{https://github.com/lsernau/grh-lean/blob/72fead102f663d5c7a407964a13dc0cf9d977207/grh-compact/src/New/Selmer/FitTwists.lean\#L250-L250}{{\leanfont Data.\allowbreak{}B₀}}}
\expandafter\newcommand\csname LN@IsFitFamily\endcsname{\href{https://github.com/lsernau/grh-lean/blob/72fead102f663d5c7a407964a13dc0cf9d977207/grh-compact/src/New/Selmer/FitTwists.lean\#L256-L263}{{\leanfont Data.\allowbreak{}IsFitFamily}}}
\expandafter\newcommand\csname LQ@IsFitFamily\endcsname{\leanquoteline{0}{structure Data.IsFitFamily (θ : PrimePadicIso → D.MFoldTwist) where}\leanquoteline{2}{coprime ι : Nat.Coprime ι.1.1 (θ ι).levelNorm}\leanquoteline{2}{denominator\_le j : ∃ γ, (∀ T : Finset ℕ, (∀ l ∈ T, l.Prime) → T.sum γ ≤ D.B₀) ∧ ∀ ι, DenominatorLe}\leanquoteline{4}{ι.2 \_ D.e₀ (D.goodLift ((θ ι).powTwist j) j).range (γ ι.1) ∧ ∀ p : D.Pair, ∀ H : p.IsRankOne j,}\leanquoteline{4}{DenominatorLe ι.2 \_ (tupleIdem \_ fun \_ ↦}\leanquoteline{4}{.X - 1) (ℂ ∙ (p.lValueData H.1 ((θ ι).powTwist j)).trivialLift) (γ ι.1)}\leanquoteline{2}{height\_le j (p : D.Pair) s : p.IsRamified → HeightLe (D.x \textasciicircum{} 999) fun ι ↦}\leanquoteline{4}{(p.ratio ((θ ι).powTwist j) s)⁻¹}}
\expandafter\newcommand\csname LN@existsisFitFamily\endcsname{\href{https://github.com/lsernau/grh-lean/blob/72fead102f663d5c7a407964a13dc0cf9d977207/grh-compact/src/New/Selmer/FitTwistsExist.lean\#L24-L65}{{\leanfont exists\_isFitFamily}}}
\expandafter\newcommand\csname LQ@existsisFitFamily\endcsname{\leanquoteline{0}{lemma exists\_isFitFamily \{P\} (D : Data P) : ∃ θ, D.IsFitFamily θ}}
\expandafter\newcommand\csname LN@AbsGal\endcsname{\href{https://github.com/lsernau/grh-lean/blob/72fead102f663d5c7a407964a13dc0cf9d977207/grh-compact/src/New/Selmer/LevelPlacesChosen.lean\#L25-L25}{{\leanfont AbsGal}}}
\expandafter\newcommand\csname LN@placesAboveG\endcsname{\href{https://github.com/lsernau/grh-lean/blob/72fead102f663d5c7a407964a13dc0cf9d977207/grh-compact/src/New/Selmer/LevelPlacesChosen.lean\#L28-L28}{{\leanfont placesAboveG}}}
\expandafter\newcommand\csname LN@badPlaces\endcsname{\href{https://github.com/lsernau/grh-lean/blob/72fead102f663d5c7a407964a13dc0cf9d977207/grh-compact/src/New/Selmer/LevelPlacesChosen.lean\#L32-L32}{{\leanfont badPlaces}}}
\expandafter\newcommand\csname LN@invAvatar\endcsname{\href{https://github.com/lsernau/grh-lean/blob/72fead102f663d5c7a407964a13dc0cf9d977207/grh-compact/src/New/Selmer/LevelPlacesChosen.lean\#L59-L60}{{\leanfont Data.\allowbreak{}TameChar.\allowbreak{}invAvatar}}}
\expandafter\newcommand\csname LN@centralRep\endcsname{\href{https://github.com/lsernau/grh-lean/blob/72fead102f663d5c7a407964a13dc0cf9d977207/grh-compact/src/New/Selmer/RankOneDenominator.lean\#L13-L14}{{\leanfont LValueData.\allowbreak{}centralRep}}}
\expandafter\newcommand\csname LN@twistedCentralOp\endcsname{\href{https://github.com/lsernau/grh-lean/blob/72fead102f663d5c7a407964a13dc0cf9d977207/grh-compact/src/New/Selmer/RankOneDenominator.lean\#L18-L18}{{\leanfont LValueData.\allowbreak{}twistedCentralOp}}}
\expandafter\newcommand\csname LN@innertwistedCentralOp\endcsname{\href{https://github.com/lsernau/grh-lean/blob/72fead102f663d5c7a407964a13dc0cf9d977207/grh-compact/src/New/Selmer/RankOneDenominator.lean\#L20-L24}{{\leanfont LValueData.\allowbreak{}inner\_twistedCentralOp}}}
\expandafter\newcommand\csname LN@twistedCentralOpkernelForm\endcsname{\href{https://github.com/lsernau/grh-lean/blob/72fead102f663d5c7a407964a13dc0cf9d977207/grh-compact/src/New/Selmer/RankOneDenominator.lean\#L28-L32}{{\leanfont LValueData.\allowbreak{}twistedCentralOp\_kernelForm}}}
\expandafter\newcommand\csname LN@twistedCentralOptrivialLift\endcsname{\href{https://github.com/lsernau/grh-lean/blob/72fead102f663d5c7a407964a13dc0cf9d977207/grh-compact/src/New/Selmer/RankOneDenominator.lean\#L34-L36}{{\leanfont LValueData.\allowbreak{}twistedCentralOp\_trivialLift}}}
\expandafter\newcommand\csname LN@rankonebound\endcsname{\href{https://github.com/lsernau/grh-lean/blob/72fead102f663d5c7a407964a13dc0cf9d977207/grh-compact/src/New/Selmer/RankOneDenominator.lean\#L43-L66}{{\leanfont DenominatorLe.\allowbreak{}rank\_one\_bound}}}
\expandafter\newcommand\csname LN@eqzeroofmemevenClasses\endcsname{\href{https://github.com/lsernau/grh-lean/blob/72fead102f663d5c7a407964a13dc0cf9d977207/grh-compact/src/New/Selmer/ReciprocityLaw.lean\#L12-L17}{{\leanfont eq\_zero\_of\_mem\_evenClasses}}}
\expandafter\newcommand\csname LN@IsSplitPlace\endcsname{\href{https://github.com/lsernau/grh-lean/blob/72fead102f663d5c7a407964a13dc0cf9d977207/grh-compact/src/New/Selmer/ReciprocityLaw.lean\#L24-L24}{{\leanfont IsSplitPlace}}}
\expandafter\newcommand\csname LN@primeBelowPlus\endcsname{\href{https://github.com/lsernau/grh-lean/blob/72fead102f663d5c7a407964a13dc0cf9d977207/grh-compact/src/New/Selmer/ReciprocityLaw.lean\#L27-L27}{{\leanfont primeBelowPlus}}}
\expandafter\newcommand\csname LN@IsLaw\endcsname{\href{https://github.com/lsernau/grh-lean/blob/72fead102f663d5c7a407964a13dc0cf9d977207/grh-compact/src/New/Selmer/ReciprocityLaw.lean\#L32-L40}{{\leanfont IsLaw}}}
\expandafter\newcommand\csname LN@primeBelowPlusstar\endcsname{\href{https://github.com/lsernau/grh-lean/blob/72fead102f663d5c7a407964a13dc0cf9d977207/grh-compact/src/New/Selmer/ReciprocityLaw.lean\#L42-L43}{{\leanfont primeBelowPlus\_star}}}
\expandafter\newcommand\csname LN@existsisLaw\endcsname{\href{https://github.com/lsernau/grh-lean/blob/72fead102f663d5c7a407964a13dc0cf9d977207/grh-compact/src/New/Selmer/ReciprocityLaw.lean\#L48-L195}{{\leanfont exists\_isLaw}}}
\expandafter\newcommand\csname LN@ribetNormaliser\endcsname{\href{https://github.com/lsernau/grh-lean/blob/72fead102f663d5c7a407964a13dc0cf9d977207/grh-compact/src/New/Selmer/RibetCramer.lean\#L79-L80}{{\leanfont ribetNormaliser}}}
\expandafter\newcommand\csname LN@rowBudget\endcsname{\href{https://github.com/lsernau/grh-lean/blob/72fead102f663d5c7a407964a13dc0cf9d977207/grh-compact/src/New/Selmer/RibetElement.lean\#L36-L37}{{\leanfont rowBudget}}}
\expandafter\newcommand\csname LN@memselmerModule\endcsname{\href{https://github.com/lsernau/grh-lean/blob/72fead102f663d5c7a407964a13dc0cf9d977207/grh-compact/src/New/Selmer/RibetElement.lean\#L42-L50}{{\leanfont obeyingCocycles.\allowbreak{}mem\_selmerModule}}}
\expandafter\newcommand\csname LN@isTraceIdealwithConjInv\endcsname{\href{https://github.com/lsernau/grh-lean/blob/72fead102f663d5c7a407964a13dc0cf9d977207/grh-compact/src/New/Selmer/RibetElement.lean\#L55-L70}{{\leanfont isTraceIdeal\_withConjInv}}}
\expandafter\newcommand\csname LN@sumblockLossle\endcsname{\href{https://github.com/lsernau/grh-lean/blob/72fead102f663d5c7a407964a13dc0cf9d977207/grh-compact/src/New/Selmer/RibetElement.lean\#L80-L114}{{\leanfont sum\_blockLoss\_le}}}
\expandafter\newcommand\csname LN@ribetcolumns\endcsname{\href{https://github.com/lsernau/grh-lean/blob/72fead102f663d5c7a407964a13dc0cf9d977207/grh-compact/src/New/Selmer/RibetElement.lean\#L122-L223}{{\leanfont ribet\_columns}}}
\expandafter\newcommand\csname LN@HasTriangleCost\endcsname{\href{https://github.com/lsernau/grh-lean/blob/72fead102f663d5c7a407964a13dc0cf9d977207/grh-compact/src/New/Selmer/RibetFirstOrder.lean\#L20-L24}{{\leanfont HasTriangleCost}}}
\expandafter\newcommand\csname LN@IsTraceIdeal\endcsname{\href{https://github.com/lsernau/grh-lean/blob/72fead102f663d5c7a407964a13dc0cf9d977207/grh-compact/src/New/Selmer/RibetFirstOrder.lean\#L67-L71}{{\leanfont IsTraceIdeal}}}
\expandafter\newcommand\csname LN@IsRibetElement\endcsname{\href{https://github.com/lsernau/grh-lean/blob/72fead102f663d5c7a407964a13dc0cf9d977207/grh-compact/src/New/Selmer/RibetFirstOrder.lean\#L75-L77}{{\leanfont IsRibetElement}}}
\expandafter\newcommand\csname LN@existsrowforalltwists\endcsname{\href{https://github.com/lsernau/grh-lean/blob/72fead102f663d5c7a407964a13dc0cf9d977207/grh-compact/src/New/Selmer/RowsExist.lean\#L29-L55}{{\leanfont exists\_row\_for\_all\_twists}}}
\expandafter\newcommand\csname LN@existshasRowauxPair\endcsname{\href{https://github.com/lsernau/grh-lean/blob/72fead102f663d5c7a407964a13dc0cf9d977207/grh-compact/src/New/Selmer/RowsExist.lean\#L61-L89}{{\leanfont exists\_hasRow\_auxPair}}}
\expandafter\newcommand\csname LN@existsisPrimeTriangle\endcsname{\href{https://github.com/lsernau/grh-lean/blob/72fead102f663d5c7a407964a13dc0cf9d977207/grh-compact/src/New/Selmer/TrianglePrimes.lean\#L24-L123}{{\leanfont exists\_isPrimeTriangle}}}
\expandafter\newcommand\csname LN@existsmFoldTwistselmerLengthLe\endcsname{\href{https://github.com/lsernau/grh-lean/blob/72fead102f663d5c7a407964a13dc0cf9d977207/grh-compact/src/New/Selmer/TwistForRows.lean\#L14-L59}{{\leanfont exists\_mFoldTwist\_selmerLengthLe}}}
\newcommand{\ConcordanceRows}{\crow{IsConjSelfDual}{\href{https://github.com/lsernau/grh-lean/blob/72fead102f663d5c7a407964a13dc0cf9d977207/grh-compact/src/New/Selmer/AptTupleDefs.lean\#L21-L24}{{\leanfont IsConjSelfDual}}}{AptTupleDefs.lean}{21--24}
\crow{crudeUnit}{\href{https://github.com/lsernau/grh-lean/blob/72fead102f663d5c7a407964a13dc0cf9d977207/grh-compact/src/New/Selmer/AptTupleDefs.lean\#L27-L27}{{\leanfont crudeUnit}}}{AptTupleDefs.lean}{27--27}
\crow{SelmerLengthLe}{\href{https://github.com/lsernau/grh-lean/blob/72fead102f663d5c7a407964a13dc0cf9d977207/grh-compact/src/New/Selmer/AptTupleDefs.lean\#L34-L38}{{\leanfont SelmerLengthLe}}}{AptTupleDefs.lean}{34--38}
\crow{IsColumn}{\href{https://github.com/lsernau/grh-lean/blob/72fead102f663d5c7a407964a13dc0cf9d977207/grh-compact/src/New/Selmer/AptTupleDefs.lean\#L43-L45}{{\leanfont IsColumn}}}{AptTupleDefs.lean}{43--45}
\crow{IsRow}{\href{https://github.com/lsernau/grh-lean/blob/72fead102f663d5c7a407964a13dc0cf9d977207/grh-compact/src/New/Selmer/AptTupleDefs.lean\#L50-L59}{{\leanfont IsRow}}}{AptTupleDefs.lean}{50--59}
\crow{chiOneInv}{\href{https://github.com/lsernau/grh-lean/blob/72fead102f663d5c7a407964a13dc0cf9d977207/grh-compact/src/New/Selmer/AptTupleDefs.lean\#L68-L68}{{\leanfont chiOneInv}}}{AptTupleDefs.lean}{68--68}
\crow{normalisedRep}{\href{https://github.com/lsernau/grh-lean/blob/72fead102f663d5c7a407964a13dc0cf9d977207/grh-compact/src/New/Selmer/AptTupleDefs.lean\#L82-L82}{{\leanfont normalisedRep}}}{AptTupleDefs.lean}{82--82}
\crow{twistedLabel}{\href{https://github.com/lsernau/grh-lean/blob/72fead102f663d5c7a407964a13dc0cf9d977207/grh-compact/src/New/Selmer/AptTupleDefs.lean\#L86-L87}{{\leanfont twistedLabel}}}{AptTupleDefs.lean}{86--87}
\crow{IsApt}{\href{https://github.com/lsernau/grh-lean/blob/72fead102f663d5c7a407964a13dc0cf9d977207/grh-compact/src/New/Selmer/AptTupleDefs.lean\#L92-L100}{{\leanfont IsApt}}}{AptTupleDefs.lean}{92--100}
\crow{tupleProj}{\href{https://github.com/lsernau/grh-lean/blob/72fead102f663d5c7a407964a13dc0cf9d977207/grh-compact/src/New/Selmer/AptTupleDefs.lean\#L104-L104}{{\leanfont tupleProj}}}{AptTupleDefs.lean}{104--104}
\crow{HasRow}{\href{https://github.com/lsernau/grh-lean/blob/72fead102f663d5c7a407964a13dc0cf9d977207/grh-compact/src/New/Selmer/AptTupleDefs.lean\#L108-L110}{{\leanfont HasRow}}}{AptTupleDefs.lean}{108--110}
\crow{IsLabelledColumn}{\href{https://github.com/lsernau/grh-lean/blob/72fead102f663d5c7a407964a13dc0cf9d977207/grh-compact/src/New/Selmer/AptTupleDefs.lean\#L114-L116}{{\leanfont IsLabelledColumn}}}{AptTupleDefs.lean}{114--116}
\crow{obeyingCocycles}{\href{https://github.com/lsernau/grh-lean/blob/72fead102f663d5c7a407964a13dc0cf9d977207/grh-compact/src/New/Selmer/BlocksAndCocycles.lean\#L16-L17}{{\leanfont obeyingCocycles}}}{BlocksAndCocycles.lean}{16--17}
\crow{RowBlockIndex}{\href{https://github.com/lsernau/grh-lean/blob/72fead102f663d5c7a407964a13dc0cf9d977207/grh-compact/src/New/Selmer/BlocksAndCocycles.lean\#L41-L41}{{\leanfont RowBlockIndex}}}{BlocksAndCocycles.lean}{41--41}
\crow{blockRep}{\href{https://github.com/lsernau/grh-lean/blob/72fead102f663d5c7a407964a13dc0cf9d977207/grh-compact/src/New/Selmer/BlocksAndCocycles.lean\#L50-L51}{{\leanfont blockRep}}}{BlocksAndCocycles.lean}{50--51}
\crow{IsAuxiliary}{\href{https://github.com/lsernau/grh-lean/blob/72fead102f663d5c7a407964a13dc0cf9d977207/grh-compact/src/New/Selmer/Corepresentable.lean\#L55-L55}{{\leanfont IsAuxiliary}}}{Corepresentable.lean}{55--55}
\crow{evenClasses}{\href{https://github.com/lsernau/grh-lean/blob/72fead102f663d5c7a407964a13dc0cf9d977207/grh-compact/src/New/Selmer/Corepresentable.lean\#L136-L138}{{\leanfont evenClasses}}}{Corepresentable.lean}{136--138}
\crow{IsNewPlaces}{\href{https://github.com/lsernau/grh-lean/blob/72fead102f663d5c7a407964a13dc0cf9d977207/grh-compact/src/New/Selmer/Corepresentable.lean\#L142-L143}{{\leanfont IsNewPlaces}}}{Corepresentable.lean}{142--143}
\crow{withConjInv}{\href{https://github.com/lsernau/grh-lean/blob/72fead102f663d5c7a407964a13dc0cf9d977207/grh-compact/src/New/Selmer/Corepresentable.lean\#L147-L147}{{\leanfont withConjInv}}}{Corepresentable.lean}{147--147}
\crow{IsPrimeTriangle}{\href{https://github.com/lsernau/grh-lean/blob/72fead102f663d5c7a407964a13dc0cf9d977207/grh-compact/src/New/Selmer/Corepresentable.lean\#L152-L162}{{\leanfont IsPrimeTriangle}}}{Corepresentable.lean}{152--162}
\crow{IsAuxTwist}{\href{https://github.com/lsernau/grh-lean/blob/72fead102f663d5c7a407964a13dc0cf9d977207/grh-compact/src/New/Selmer/Corepresentable.lean\#L166-L167}{{\leanfont IsAuxTwist}}}{Corepresentable.lean}{166--167}
\crow{existscycCharpowsubonedvd}{\href{https://github.com/lsernau/grh-lean/blob/72fead102f663d5c7a407964a13dc0cf9d977207/grh-compact/src/New/Selmer/DenominatorBound.lean\#L26-L40}{{\leanfont exists\_cycChar\_pow\_sub\_one\_dvd}}}{DenominatorBound.lean}{26--40}
\crow{denominatorLerowBudget}{\href{https://github.com/lsernau/grh-lean/blob/72fead102f663d5c7a407964a13dc0cf9d977207/grh-compact/src/New/Selmer/DenominatorBound.lean\#L47-L166}{{\leanfont denominatorLe\_rowBudget}}}{DenominatorBound.lean}{47--166}
\crow{crysRels}{\href{https://github.com/lsernau/grh-lean/blob/72fead102f663d5c7a407964a13dc0cf9d977207/grh-compact/src/New/Selmer/DeterminantBound.lean\#L43-L43}{{\leanfont crysRels}}}{DeterminantBound.lean}{43--43}
\crow{blockLoss}{\href{https://github.com/lsernau/grh-lean/blob/72fead102f663d5c7a407964a13dc0cf9d977207/grh-compact/src/New/Selmer/DeterminantBound.lean\#L47-L48}{{\leanfont blockLoss}}}{DeterminantBound.lean}{47--48}
\crow{cmSign}{\href{https://github.com/lsernau/grh-lean/blob/72fead102f663d5c7a407964a13dc0cf9d977207/grh-compact/src/New/Selmer/DeterminantBound.lean\#L56-L56}{{\leanfont cmSign}}}{DeterminantBound.lean}{56--56}
\crow{IsRowType}{\href{https://github.com/lsernau/grh-lean/blob/72fead102f663d5c7a407964a13dc0cf9d977207/grh-compact/src/New/Selmer/DeterminantBound.lean\#L64-L70}{{\leanfont IsRowType}}}{DeterminantBound.lean}{64--70}
\crow{lambdaLoss}{\href{https://github.com/lsernau/grh-lean/blob/72fead102f663d5c7a407964a13dc0cf9d977207/grh-compact/src/New/Selmer/DeterminantBound.lean\#L74-L75}{{\leanfont lambdaLoss}}}{DeterminantBound.lean}{74--75}
\crow{loss}{\href{https://github.com/lsernau/grh-lean/blob/72fead102f663d5c7a407964a13dc0cf9d977207/grh-compact/src/New/Selmer/DeterminantBound.lean\#L78-L79}{{\leanfont loss}}}{DeterminantBound.lean}{78--79}
\crow{innerAtheckePolyOnForms}{\href{https://github.com/lsernau/grh-lean/blob/72fead102f663d5c7a407964a13dc0cf9d977207/grh-compact/src/New/Selmer/FitTwists.lean\#L19-L24}{{\leanfont LevelPlaces.\allowbreak{}LevelForm.\allowbreak{}innerAt\_heckePolyOnForms}}}{FitTwists.lean}{19--24}
\crow{groupAlgAction}{\href{https://github.com/lsernau/grh-lean/blob/72fead102f663d5c7a407964a13dc0cf9d977207/grh-compact/src/New/Selmer/FitTwists.lean\#L35-L35}{{\leanfont AuxPair.\allowbreak{}groupAlgAction}}}{FitTwists.lean}{35--35}
\crow{IsTested}{\href{https://github.com/lsernau/grh-lean/blob/72fead102f663d5c7a407964a13dc0cf9d977207/grh-compact/src/New/Selmer/FitTwists.lean\#L39-L40}{{\leanfont AuxPair.\allowbreak{}IsTested}}}{FitTwists.lean}{39--40}
\crow{IsTestForm}{\href{https://github.com/lsernau/grh-lean/blob/72fead102f663d5c7a407964a13dc0cf9d977207/grh-compact/src/New/Selmer/FitTwists.lean\#L44-L45}{{\leanfont AuxPair.\allowbreak{}IsTestForm}}}{FitTwists.lean}{44--45}
\crow{DenominatorLe}{\href{https://github.com/lsernau/grh-lean/blob/72fead102f663d5c7a407964a13dc0cf9d977207/grh-compact/src/New/Selmer/FitTwists.lean\#L50-L53}{{\leanfont DenominatorLe}}}{FitTwists.lean}{50--53}
\crow{isOfFinOrderprodpsidet}{\href{https://github.com/lsernau/grh-lean/blob/72fead102f663d5c7a407964a13dc0cf9d977207/grh-compact/src/New/Selmer/FitTwists.lean\#L59-L61}{{\leanfont AuxPair.\allowbreak{}isOfFinOrder\_prod\_psi\_det}}}{FitTwists.lean}{59--61}
\crow{innerAtformAction}{\href{https://github.com/lsernau/grh-lean/blob/72fead102f663d5c7a407964a13dc0cf9d977207/grh-compact/src/New/Selmer/FitTwists.lean\#L66-L70}{{\leanfont AuxPair.\allowbreak{}innerAt\_formAction}}}{FitTwists.lean}{66--70}
\crow{rangegroupAlgActionmeminvtSubmodule}{\href{https://github.com/lsernau/grh-lean/blob/72fead102f663d5c7a407964a13dc0cf9d977207/grh-compact/src/New/Selmer/FitTwists.lean\#L72-L74}{{\leanfont AuxPair.\allowbreak{}range\_groupAlgAction\_mem\_invtSubmodule}}}{FitTwists.lean}{72--74}
\crow{heckePolyOnFormscomm}{\href{https://github.com/lsernau/grh-lean/blob/72fead102f663d5c7a407964a13dc0cf9d977207/grh-compact/src/New/Selmer/FitTwists.lean\#L78-L82}{{\leanfont AuxPair.\allowbreak{}heckePolyOnForms\_comm}}}{FitTwists.lean}{78--82}
\crow{innerAtgroupAlgAction}{\href{https://github.com/lsernau/grh-lean/blob/72fead102f663d5c7a407964a13dc0cf9d977207/grh-compact/src/New/Selmer/FitTwists.lean\#L86-L88}{{\leanfont AuxPair.\allowbreak{}innerAt\_groupAlgAction}}}{FitTwists.lean}{86--88}
\crow{normSqAtstarProjectionle}{\href{https://github.com/lsernau/grh-lean/blob/72fead102f663d5c7a407964a13dc0cf9d977207/grh-compact/src/New/Selmer/FitTwists.lean\#L96-L112}{{\leanfont AuxPair.\allowbreak{}normSqAt\_starProjection\_le}}}{FitTwists.lean}{96--112}
\crow{IsSplitOffLevel}{\href{https://github.com/lsernau/grh-lean/blob/72fead102f663d5c7a407964a13dc0cf9d977207/grh-compact/src/New/Selmer/FitTwists.lean\#L117-L117}{{\leanfont Data.\allowbreak{}IsSplitOffLevel}}}{FitTwists.lean}{117--117}
\crow{heckeAlgebra}{\href{https://github.com/lsernau/grh-lean/blob/72fead102f663d5c7a407964a13dc0cf9d977207/grh-compact/src/New/Selmer/FitTwists.lean\#L122-L127}{{\leanfont Data.\allowbreak{}heckeAlgebra}}}{FitTwists.lean}{122--127}
\crow{heckeOperator}{\href{https://github.com/lsernau/grh-lean/blob/72fead102f663d5c7a407964a13dc0cf9d977207/grh-compact/src/New/Selmer/FitTwists.lean\#L131-L136}{{\leanfont Data.\allowbreak{}heckeOperator}}}{FitTwists.lean}{131--136}
\crow{eigenPart}{\href{https://github.com/lsernau/grh-lean/blob/72fead102f663d5c7a407964a13dc0cf9d977207/grh-compact/src/New/Selmer/FitTwists.lean\#L140-L140}{{\leanfont Data.\allowbreak{}eigenPart}}}{FitTwists.lean}{140--140}
\crow{isSphericalWithofmemeigenPart}{\href{https://github.com/lsernau/grh-lean/blob/72fead102f663d5c7a407964a13dc0cf9d977207/grh-compact/src/New/Selmer/FitTwists.lean\#L145-L148}{{\leanfont Data.\allowbreak{}isSphericalWith\_of\_mem\_eigenPart}}}{FitTwists.lean}{145--148}
\crow{occurringLabels}{\href{https://github.com/lsernau/grh-lean/blob/72fead102f663d5c7a407964a13dc0cf9d977207/grh-compact/src/New/Selmer/FitTwists.lean\#L152-L153}{{\leanfont Data.\allowbreak{}occurringLabels}}}{FitTwists.lean}{152--153}
\crow{dvdcardresFieldsubone}{\href{https://github.com/lsernau/grh-lean/blob/72fead102f663d5c7a407964a13dc0cf9d977207/grh-compact/src/New/Selmer/FitTwists.lean\#L155-L158}{{\leanfont Data.\allowbreak{}dvd\_card\_resField\_sub\_one}}}{FitTwists.lean}{155--158}
\crow{existsheckeprojeigenPart}{\href{https://github.com/lsernau/grh-lean/blob/72fead102f663d5c7a407964a13dc0cf9d977207/grh-compact/src/New/Selmer/FitTwists.lean\#L162-L167}{{\leanfont Data.\allowbreak{}exists\_hecke\_proj\_eigenPart}}}{FitTwists.lean}{162--167}
\crow{commuteheckeAlgebraprincipalRep}{\href{https://github.com/lsernau/grh-lean/blob/72fead102f663d5c7a407964a13dc0cf9d977207/grh-compact/src/New/Selmer/FitTwists.lean\#L169-L173}{{\leanfont Data.\allowbreak{}commute\_heckeAlgebra\_principalRep}}}{FitTwists.lean}{169--173}
\crow{heckeOnGood}{\href{https://github.com/lsernau/grh-lean/blob/72fead102f663d5c7a407964a13dc0cf9d977207/grh-compact/src/New/Selmer/FitTwists.lean\#L177-L179}{{\leanfont Data.\allowbreak{}heckeOnGood}}}{FitTwists.lean}{177--179}
\crow{principalRepmemeigenPart}{\href{https://github.com/lsernau/grh-lean/blob/72fead102f663d5c7a407964a13dc0cf9d977207/grh-compact/src/New/Selmer/FitTwists.lean\#L181-L183}{{\leanfont Data.\allowbreak{}principalRep\_mem\_eigenPart}}}{FitTwists.lean}{181--183}
\crow{finiteeigenPart}{\href{https://github.com/lsernau/grh-lean/blob/72fead102f663d5c7a407964a13dc0cf9d977207/grh-compact/src/New/Selmer/FitTwists.lean\#L187-L188}{{\leanfont Data.\allowbreak{}finite\_eigenPart}}}{FitTwists.lean}{187--188}
\crow{iSupeigenPart}{\href{https://github.com/lsernau/grh-lean/blob/72fead102f663d5c7a407964a13dc0cf9d977207/grh-compact/src/New/Selmer/FitTwists.lean\#L192-L197}{{\leanfont Data.\allowbreak{}iSup\_eigenPart}}}{FitTwists.lean}{192--197}
\crow{thetamuleqthetaformAction}{\href{https://github.com/lsernau/grh-lean/blob/72fead102f663d5c7a407964a13dc0cf9d977207/grh-compact/src/New/Selmer/FitTwists.lean\#L205-L212}{{\leanfont AuxPair.\allowbreak{}theta\_mul\_eq\_theta\_formAction}}}{FitTwists.lean}{205--212}
\crow{existsthetatoprincipalForms}{\href{https://github.com/lsernau/grh-lean/blob/72fead102f663d5c7a407964a13dc0cf9d977207/grh-compact/src/New/Selmer/FitTwists.lean\#L216-L222}{{\leanfont Data.\allowbreak{}exists\_theta\_to\_principalForms}}}{FitTwists.lean}{216--222}
\crow{auxLift}{\href{https://github.com/lsernau/grh-lean/blob/72fead102f663d5c7a407964a13dc0cf9d977207/grh-compact/src/New/Selmer/FitTwists.lean\#L226-L227}{{\leanfont Data.\allowbreak{}auxLift}}}{FitTwists.lean}{226--227}
\crow{RankOneIndex}{\href{https://github.com/lsernau/grh-lean/blob/72fead102f663d5c7a407964a13dc0cf9d977207/grh-compact/src/New/Selmer/FitTwists.lean\#L231-L231}{{\leanfont Data.\allowbreak{}RankOneIndex}}}{FitTwists.lean}{231--231}
\crow{goodLift}{\href{https://github.com/lsernau/grh-lean/blob/72fead102f663d5c7a407964a13dc0cf9d977207/grh-compact/src/New/Selmer/FitTwists.lean\#L237-L237}{{\leanfont Data.\allowbreak{}goodLift}}}{FitTwists.lean}{237--237}
\crow{IsRamified}{\href{https://github.com/lsernau/grh-lean/blob/72fead102f663d5c7a407964a13dc0cf9d977207/grh-compact/src/New/Selmer/FitTwists.lean\#L241-L241}{{\leanfont Data.\allowbreak{}Pair.\allowbreak{}IsRamified}}}{FitTwists.lean}{241--241}
\crow{IsRankOne}{\href{https://github.com/lsernau/grh-lean/blob/72fead102f663d5c7a407964a13dc0cf9d977207/grh-compact/src/New/Selmer/FitTwists.lean\#L245-L247}{{\leanfont Data.\allowbreak{}Pair.\allowbreak{}IsRankOne}}}{FitTwists.lean}{245--247}
\crow{B}{\href{https://github.com/lsernau/grh-lean/blob/72fead102f663d5c7a407964a13dc0cf9d977207/grh-compact/src/New/Selmer/FitTwists.lean\#L250-L250}{{\leanfont Data.\allowbreak{}B₀}}}{FitTwists.lean}{250--250}
\crow{IsFitFamily}{\href{https://github.com/lsernau/grh-lean/blob/72fead102f663d5c7a407964a13dc0cf9d977207/grh-compact/src/New/Selmer/FitTwists.lean\#L256-L263}{{\leanfont Data.\allowbreak{}IsFitFamily}}}{FitTwists.lean}{256--263}
\crow{existsisFitFamily}{\href{https://github.com/lsernau/grh-lean/blob/72fead102f663d5c7a407964a13dc0cf9d977207/grh-compact/src/New/Selmer/FitTwistsExist.lean\#L24-L65}{{\leanfont exists\_isFitFamily}}}{FitTwistsExist.lean}{24--65}
\crow{AbsGal}{\href{https://github.com/lsernau/grh-lean/blob/72fead102f663d5c7a407964a13dc0cf9d977207/grh-compact/src/New/Selmer/LevelPlacesChosen.lean\#L25-L25}{{\leanfont AbsGal}}}{LevelPlacesChosen.lean}{25--25}
\crow{placesAboveG}{\href{https://github.com/lsernau/grh-lean/blob/72fead102f663d5c7a407964a13dc0cf9d977207/grh-compact/src/New/Selmer/LevelPlacesChosen.lean\#L28-L28}{{\leanfont placesAboveG}}}{LevelPlacesChosen.lean}{28--28}
\crow{badPlaces}{\href{https://github.com/lsernau/grh-lean/blob/72fead102f663d5c7a407964a13dc0cf9d977207/grh-compact/src/New/Selmer/LevelPlacesChosen.lean\#L32-L32}{{\leanfont badPlaces}}}{LevelPlacesChosen.lean}{32--32}
\crow{invAvatar}{\href{https://github.com/lsernau/grh-lean/blob/72fead102f663d5c7a407964a13dc0cf9d977207/grh-compact/src/New/Selmer/LevelPlacesChosen.lean\#L59-L60}{{\leanfont Data.\allowbreak{}TameChar.\allowbreak{}invAvatar}}}{LevelPlacesChosen.lean}{59--60}
\crow{centralRep}{\href{https://github.com/lsernau/grh-lean/blob/72fead102f663d5c7a407964a13dc0cf9d977207/grh-compact/src/New/Selmer/RankOneDenominator.lean\#L13-L14}{{\leanfont LValueData.\allowbreak{}centralRep}}}{RankOneDenominator.lean}{13--14}
\crow{twistedCentralOp}{\href{https://github.com/lsernau/grh-lean/blob/72fead102f663d5c7a407964a13dc0cf9d977207/grh-compact/src/New/Selmer/RankOneDenominator.lean\#L18-L18}{{\leanfont LValueData.\allowbreak{}twistedCentralOp}}}{RankOneDenominator.lean}{18--18}
\crow{innertwistedCentralOp}{\href{https://github.com/lsernau/grh-lean/blob/72fead102f663d5c7a407964a13dc0cf9d977207/grh-compact/src/New/Selmer/RankOneDenominator.lean\#L20-L24}{{\leanfont LValueData.\allowbreak{}inner\_twistedCentralOp}}}{RankOneDenominator.lean}{20--24}
\crow{twistedCentralOpkernelForm}{\href{https://github.com/lsernau/grh-lean/blob/72fead102f663d5c7a407964a13dc0cf9d977207/grh-compact/src/New/Selmer/RankOneDenominator.lean\#L28-L32}{{\leanfont LValueData.\allowbreak{}twistedCentralOp\_kernelForm}}}{RankOneDenominator.lean}{28--32}
\crow{twistedCentralOptrivialLift}{\href{https://github.com/lsernau/grh-lean/blob/72fead102f663d5c7a407964a13dc0cf9d977207/grh-compact/src/New/Selmer/RankOneDenominator.lean\#L34-L36}{{\leanfont LValueData.\allowbreak{}twistedCentralOp\_trivialLift}}}{RankOneDenominator.lean}{34--36}
\crow{rankonebound}{\href{https://github.com/lsernau/grh-lean/blob/72fead102f663d5c7a407964a13dc0cf9d977207/grh-compact/src/New/Selmer/RankOneDenominator.lean\#L43-L66}{{\leanfont DenominatorLe.\allowbreak{}rank\_one\_bound}}}{RankOneDenominator.lean}{43--66}
\crow{eqzeroofmemevenClasses}{\href{https://github.com/lsernau/grh-lean/blob/72fead102f663d5c7a407964a13dc0cf9d977207/grh-compact/src/New/Selmer/ReciprocityLaw.lean\#L12-L17}{{\leanfont eq\_zero\_of\_mem\_evenClasses}}}{ReciprocityLaw.lean}{12--17}
\crow{IsSplitPlace}{\href{https://github.com/lsernau/grh-lean/blob/72fead102f663d5c7a407964a13dc0cf9d977207/grh-compact/src/New/Selmer/ReciprocityLaw.lean\#L24-L24}{{\leanfont IsSplitPlace}}}{ReciprocityLaw.lean}{24--24}
\crow{primeBelowPlus}{\href{https://github.com/lsernau/grh-lean/blob/72fead102f663d5c7a407964a13dc0cf9d977207/grh-compact/src/New/Selmer/ReciprocityLaw.lean\#L27-L27}{{\leanfont primeBelowPlus}}}{ReciprocityLaw.lean}{27--27}
\crow{IsLaw}{\href{https://github.com/lsernau/grh-lean/blob/72fead102f663d5c7a407964a13dc0cf9d977207/grh-compact/src/New/Selmer/ReciprocityLaw.lean\#L32-L40}{{\leanfont IsLaw}}}{ReciprocityLaw.lean}{32--40}
\crow{primeBelowPlusstar}{\href{https://github.com/lsernau/grh-lean/blob/72fead102f663d5c7a407964a13dc0cf9d977207/grh-compact/src/New/Selmer/ReciprocityLaw.lean\#L42-L43}{{\leanfont primeBelowPlus\_star}}}{ReciprocityLaw.lean}{42--43}
\crow{existsisLaw}{\href{https://github.com/lsernau/grh-lean/blob/72fead102f663d5c7a407964a13dc0cf9d977207/grh-compact/src/New/Selmer/ReciprocityLaw.lean\#L48-L195}{{\leanfont exists\_isLaw}}}{ReciprocityLaw.lean}{48--195}
\crow{ribetNormaliser}{\href{https://github.com/lsernau/grh-lean/blob/72fead102f663d5c7a407964a13dc0cf9d977207/grh-compact/src/New/Selmer/RibetCramer.lean\#L79-L80}{{\leanfont ribetNormaliser}}}{RibetCramer.lean}{79--80}
\crow{rowBudget}{\href{https://github.com/lsernau/grh-lean/blob/72fead102f663d5c7a407964a13dc0cf9d977207/grh-compact/src/New/Selmer/RibetElement.lean\#L36-L37}{{\leanfont rowBudget}}}{RibetElement.lean}{36--37}
\crow{memselmerModule}{\href{https://github.com/lsernau/grh-lean/blob/72fead102f663d5c7a407964a13dc0cf9d977207/grh-compact/src/New/Selmer/RibetElement.lean\#L42-L50}{{\leanfont obeyingCocycles.\allowbreak{}mem\_selmerModule}}}{RibetElement.lean}{42--50}
\crow{isTraceIdealwithConjInv}{\href{https://github.com/lsernau/grh-lean/blob/72fead102f663d5c7a407964a13dc0cf9d977207/grh-compact/src/New/Selmer/RibetElement.lean\#L55-L70}{{\leanfont isTraceIdeal\_withConjInv}}}{RibetElement.lean}{55--70}
\crow{sumblockLossle}{\href{https://github.com/lsernau/grh-lean/blob/72fead102f663d5c7a407964a13dc0cf9d977207/grh-compact/src/New/Selmer/RibetElement.lean\#L80-L114}{{\leanfont sum\_blockLoss\_le}}}{RibetElement.lean}{80--114}
\crow{ribetcolumns}{\href{https://github.com/lsernau/grh-lean/blob/72fead102f663d5c7a407964a13dc0cf9d977207/grh-compact/src/New/Selmer/RibetElement.lean\#L122-L223}{{\leanfont ribet\_columns}}}{RibetElement.lean}{122--223}
\crow{HasTriangleCost}{\href{https://github.com/lsernau/grh-lean/blob/72fead102f663d5c7a407964a13dc0cf9d977207/grh-compact/src/New/Selmer/RibetFirstOrder.lean\#L20-L24}{{\leanfont HasTriangleCost}}}{RibetFirstOrder.lean}{20--24}
\crow{IsTraceIdeal}{\href{https://github.com/lsernau/grh-lean/blob/72fead102f663d5c7a407964a13dc0cf9d977207/grh-compact/src/New/Selmer/RibetFirstOrder.lean\#L67-L71}{{\leanfont IsTraceIdeal}}}{RibetFirstOrder.lean}{67--71}
\crow{IsRibetElement}{\href{https://github.com/lsernau/grh-lean/blob/72fead102f663d5c7a407964a13dc0cf9d977207/grh-compact/src/New/Selmer/RibetFirstOrder.lean\#L75-L77}{{\leanfont IsRibetElement}}}{RibetFirstOrder.lean}{75--77}
\crow{existsrowforalltwists}{\href{https://github.com/lsernau/grh-lean/blob/72fead102f663d5c7a407964a13dc0cf9d977207/grh-compact/src/New/Selmer/RowsExist.lean\#L29-L55}{{\leanfont exists\_row\_for\_all\_twists}}}{RowsExist.lean}{29--55}
\crow{existshasRowauxPair}{\href{https://github.com/lsernau/grh-lean/blob/72fead102f663d5c7a407964a13dc0cf9d977207/grh-compact/src/New/Selmer/RowsExist.lean\#L61-L89}{{\leanfont exists\_hasRow\_auxPair}}}{RowsExist.lean}{61--89}
\crow{existsisPrimeTriangle}{\href{https://github.com/lsernau/grh-lean/blob/72fead102f663d5c7a407964a13dc0cf9d977207/grh-compact/src/New/Selmer/TrianglePrimes.lean\#L24-L123}{{\leanfont exists\_isPrimeTriangle}}}{TrianglePrimes.lean}{24--123}
\crow{existsmFoldTwistselmerLengthLe}{\href{https://github.com/lsernau/grh-lean/blob/72fead102f663d5c7a407964a13dc0cf9d977207/grh-compact/src/New/Selmer/TwistForRows.lean\#L14-L59}{{\leanfont exists\_mFoldTwist\_selmerLengthLe}}}{TwistForRows.lean}{14--59}}
% the notations of the defined functions: one slot for each explicit argument of the Lean declaration
\newcommand{\IsConjSelfDual}[2]{\mathsf{CSD}_{#1}(#2)}
\newcommand{\crudeUnit}[2]{\mathrm{cr}_{#1}(#2)}
\newcommand{\SelmerLengthLe}[6]{\mathsf{Short}(#1,#2;\,#3,#4;\,#5,#6)}
\newcommand{\IsColumn}[5]{\mathsf{ColRep}_{#1}(#2,#3,#4;\,#5)}
\newcommand{\IsRow}[4]{\mathsf{RowRep}_{#1}(#2,#3;\,#4)}
\newcommand{\chiOneInv}[1]{\chi_{1,#1}^{-1}}
\newcommand{\normalisedRep}[2]{\operatorname{nm}_{#1}(#2)}
\newcommand{\twistedLabel}[3]{\operatorname{tl}_{#1}(#2,#3)}
\newcommand{\IsApt}[3]{\mathsf{Apt}(#1,#2;\,#3)}
\newcommand{\tupleProj}[2]{\operatorname{pr}_{#1}(#2)}
\newcommand{\HasRow}[3]{\mathsf{HasRow}(#1,#2;\,#3)}
\newcommand{\IsLabelledColumn}[6]{\mathsf{LabCol}(#1;\,#2,#3,#5;\,#4;\,#6)}
\newcommand{\obeyingCocycles}[3]{\mathcal{Z}_{#1}(#2;\,#3)}
\newcommand{\RowBlockIndex}[2]{\mathrm{I}_{#1}(#2)}
\newcommand{\blockRep}[3]{\operatorname{blk}_{#1}(#2)(#3)}
\newcommand{\IsAuxiliary}[2]{\mathsf{Auxil}(#1;\,#2)}
\newcommand{\evenClasses}[4]{\mathrm{Ev}(#1,#2;\,#3,#4)}
\newcommand{\IsNewPlaces}[4]{\mathsf{New}_{#1}(#2,#3,#4)}
\newcommand{\withConjInv}[2]{#1^{[#2]}}
\newcommand{\IsPrimeTriangle}[9]{\mathsf{Tri}_{#1}(#2,#3,#4;\,#5;\,#6,#7;\,#8,#9)}
\newcommand{\IsAuxTwist}[4]{\mathsf{AuxTw}(#1,#2;\,#3,#4)}
\newcommand{\crysRels}[1]{\mathcal{R}_{\mathrm{cris}}(#1)}
\newcommand{\blockLoss}[3]{\operatorname{bl}(#1;\,#2,#3)}
\newcommand{\cmSign}[1]{\operatorname{sgn}_{\mathrm{CM}}(#1)}
\newcommand{\IsRowType}[4]{\mathsf{RowType}_{#1}(#2,#3;\,#4)}
\newcommand{\lambdaLoss}[3]{\operatorname{lam}(#1,#2;\,#3)}
\newcommand{\loss}[2]{\operatorname{loss}_{#1}(#2)}
\newcommand{\groupAlgAction}[2]{\operatorname{act}_{#1}(#2)}
\newcommand{\IsTestForm}[3]{\mathsf{TestForm}(#1;\,#2;\,#3)}
\newcommand{\DenominatorLe}[5]{\mathsf{Den}_{#1}(#2,#3,#4;\,#5)}
\newcommand{\IsSplitOffLevel}[2]{\mathsf{SplitOff}_{#1}(#2)}
\newcommand{\heckeOperator}[2]{\operatorname{T}_{#1}(#2)}
\newcommand{\eigenPart}[3]{\operatorname{Eig}_{#1}(#2;\,#3)}
\newcommand{\occurringLabels}[4]{\operatorname{Lab}_{#1}(#2,#3;\,#4)}
\newcommand{\heckeOnGood}[3]{\operatorname{hg}_{#1,#2}(#3)}
\newcommand{\auxLift}[4]{\theta^{*}_{#1,#2,#3}(#4)}
\newcommand{\RankOneIndex}[1]{\mathrm{I}^{(1)}(#1)}
\newcommand{\goodLift}[3]{\theta^{*,\mathrm{good}}_{#1,#2,#3}}
\newcommand{\IsRamified}[1]{\mathsf{Ram}(#1)}
\newcommand{\IsRankOne}[2]{\mathsf{RkOne}(#1;\,#2)}
\newcommand{\Bzero}[1]{B_0(#1)}
\newcommand{\AbsGal}[1]{G_{E(#1)}}
\newcommand{\placesAboveG}[1]{\Sigma_{g}(#1)}
\newcommand{\badPlaces}[2]{\Sigma_{\mathrm{bad}}(#1,#2)}
\newcommand{\invAvatar}[2]{\operatorname{av}^{-}_{#1}(#2)}
\newcommand{\centralRep}[1]{\operatorname{cen}_{#1}}
\newcommand{\twistedCentralOp}[2]{\operatorname{tc}_{#1}(#2)}
\newcommand{\primeBelowPlus}[1]{#1^{+}}
\newcommand{\ribetNormaliser}[6]{\operatorname{rn}_{#1,#2}(#3,#4,#5)(#6)}
\newcommand{\rowBudget}[2]{\operatorname{rb}_{#2}(#1)}
\newcommand{\HasTriangleCost}[6]{\mathsf{TriCost}(#1;\,#2,#3;\,#4;\,#5,#6)}
\newcommand{\IsTraceIdeal}[7]{\mathsf{TrId}(#1;\,#2;\,#3;\,#4,#5,#6,#7)}
\newcommand{\IsRibetElement}[4]{\mathsf{Rib}(#1;\,#2;\,#3,#4)}
\newcommand{\UnitMembers}{205}
\newcommand{\UnitInNote}{92}
\newcommand{\UnitProofs}{96}
\newcommand{\UnitSteps}{975}
\newcommand{\UnitAllSteps}{1091}


\linespread{1.08}                       % 11 pt on 14.7 pt
\setlength{\parindent}{0pt}
\setlength{\parskip}{5.5pt plus 1.5pt minus 0.5pt}
\setlength{\abovedisplayskip}{7.5pt plus 2pt minus 3pt}
\setlength{\belowdisplayskip}{7.5pt plus 2pt minus 3pt}
\setlength{\abovedisplayshortskip}{3pt plus 2pt}
\setlength{\belowdisplayshortskip}{6pt plus 2pt minus 2pt}
\emergencystretch=1.5em
\clubpenalty=10000 \widowpenalty=10000 \displaywidowpenalty=10000
\relpenalty=7000 \binoppenalty=9000
\raggedbottom
\pagestyle{fancy}\fancyhf{}\renewcommand{\headrulewidth}{0pt}
\fancyfoot[C]{\sffamily\fontsize{7.6}{9}\selectfont\color{inkfaint}\thepage\,/\,\pageref*{LastPage}}
\setlist{topsep=2pt,partopsep=0pt,parsep=0pt,itemsep=4pt,leftmargin=2.1em,labelsep=0.55em}

\newcommand{\Q}{\mathbb{Q}}
\newcommand{\Z}{\mathbb{Z}}
\newcommand{\F}{\mathbb{F}}
\newcommand{\len}{\operatorname{len}}
\newcommand{\cok}{\operatorname{cok}}
\newcommand{\End}{\operatorname{End}}
\newcommand{\Hom}{\operatorname{Hom}}
\newcommand{\Gal}{\operatorname{Gal}}
\newcommand{\diag}{\operatorname{diag}}
\newcommand{\kw}[1]{\emph{#1}}
\newcommand{\m}[1]{\mbox{$#1$}}        % a short formula that must not be broken across lines
\newcommand{\head}[1]{\par\addvspace{17pt plus 4pt minus 2pt}\Needspace{6\baselineskip}%
  \noindent{\fontsize{13.2}{16}\selectfont\bfseries\color{ink}#1}\par\nopagebreak\vspace{-1pt}}
\newcommand{\sect}[2]{\refstepcounter{section}\label{#2}\par\addvspace{19pt plus 4pt minus 2pt}\Needspace{7\baselineskip}%
  \noindent\llap{{\fontsize{13.2}{16}\selectfont\bfseries\color{inkfaint}\thesection}\hspace{9pt}}%
  {\fontsize{13.2}{16}\selectfont\bfseries\color{ink}#1}\par\nopagebreak\vspace{-1pt}}
\newcommand{\subhead}[1]{\par\addvspace{9pt plus 2pt}\Needspace{4\baselineskip}\noindent\textbf{#1}\enspace\ignorespaces}
% where a statement lives in the Lean archive
\newcommand{\leantag}{{\sffamily\bfseries\fontsize{6.3}{8}\selectfont\color{leancol}\addfontfeatures{LetterSpace=7}LEAN}}
\newcommand{\sep}{\unskip\nobreak\hspace{5pt}{\color{inkfaint}\textperiodcentered}\hspace{5pt plus 2pt}\ignorespaces}
\newcommand{\leanline}[1]{\par\nopagebreak\vspace{-2.5pt}{\fontsize{8.9}{12.4}\selectfont\color{inksoft}\raggedright
  \noindent\llap{\leantag\hspace{9pt}}#1\par}\vspace{1pt}}
% a displayed statement in a tinted box: #1 colour, #2 text
\newcommand{\statement}[2]{\par\addvspace{7pt}\begingroup\setlength{\fboxsep}{8.5pt}\setlength{\fboxrule}{0.6pt}%
  \noindent\fcolorbox{#1!30!white}{#1!3!white}{\parbox{\dimexpr\linewidth-2\fboxsep-2\fboxrule\relax}{#2}}%
  \endgroup\par\addvspace{7pt}}
\newcommand{\nt}[2]{\par\noindent\hangindent=1.1em\hangafter=1 #1\enspace{\color{inksoft}#2}\par\vspace{1pt}}
\newcommand{\legdot}[1]{\tikz[baseline=-0.62ex]\fill[#1] (0,0) circle (2.5pt);\hspace{3pt}}
\newcommand{\crow}[4]{\csname CS@#1\endcsname & #2 & \textcolor{inksoft}{#3} & \textcolor{inksoft}{#4}\\}
\newcommand{\leanquoteline}[2]{\par\noindent\hangindent=\dimexpr#1\fontdimen2\font+2.2em\relax\hangafter=1
  \hspace*{#1\fontdimen2\font}#2}
\newcommand{\quoted}[1]{\par\vspace{3pt}{\ttfamily\fontsize{7.5}{10.4}\selectfont\color{leancol}\raggedright
  \setlength{\parskip}{0pt}\leanquote{#1}\par}\vspace{4pt}}
\newcommand{\bx}[3]{{\sffamily\bfseries\fontsize{7.3}{9.2}\selectfont\color{#1!82!black}#2}\\[1.5pt]#3}
\newcommand{\cs}[2]{\expandafter\newcommand\csname CS@#1\endcsname{#2}}

% theorems
\newtheoremstyle{plainx}{7pt plus 2pt}{7pt plus 2pt}{\itshape}{}{\bfseries}{.}{0.5em}{}
\newtheoremstyle{uprightx}{7pt plus 2pt}{7pt plus 2pt}{}{}{\bfseries}{.}{0.5em}{}
\theoremstyle{plainx}
\newtheorem{theorem}{Theorem}[section]
\newtheorem{lemma}[theorem]{Lemma}
\newtheorem{proposition}[theorem]{Proposition}
\newtheorem{corollary}[theorem]{Corollary}
\theoremstyle{uprightx}
\newtheorem{definition}[theorem]{Definition}
\newtheorem{setting}[theorem]{Setting}          % hypotheses shared by several statements; each of them cites it by number
\newtheorem{remark}[theorem]{Remark}
\newtheorem{example}[theorem]{Example}
\renewcommand{\qedsymbol}{{\setlength{\fboxsep}{0pt}\setlength{\fboxrule}{0.5pt}\fbox{\rule{0pt}{0.46em}\hspace{0.46em}}}}
% what holds throughout the note: the ambient objects only (implicit and instance arguments), never a hypothesis
\newenvironment{conventions}{\par\addvspace{5pt}\noindent\textbf{Conventions.}\enspace\ignorespaces}{\par\addvspace{3pt}}
\makeatletter
\renewenvironment{proof}[1][Proof]{\par\pushQED{\qed}\normalfont\topsep0pt\partopsep0pt
  \trivlist\item[\hskip\labelsep\itshape #1.]\ignorespaces}{\popQED\endtrivlist\@endpefalse}
\makeatother

% ---- macros of the three notes of this unit (background notions; the notions that the unit defines are in data.tex)
\newcommand{\CC}{\mathbb{C}}
\newcommand{\NN}{\mathbb{N}}
\newcommand{\Qb}{\bar{\mathbb{Q}}}
\newcommand{\Ql}{\mathbb{Q}_\ell}
\newcommand{\Zl}{\mathbb{Z}_\ell}
\newcommand{\Qlb}{\bar{\mathbb{Q}}_\ell}
\newcommand{\Zlb}{\bar{\mathbb{Z}}_\ell}
\newcommand{\Cl}{\mathbb{C}_\ell}
\newcommand{\Gl}{G_{\mathbb{Q}_\ell}}
\newcommand{\GE}[1]{G_{#1}}
\newcommand{\cyc}{\operatorname{cyc}}
\newcommand{\vl}{v_\ell}
\newcommand{\Sel}[4]{\operatorname{Sel}_{#2}^{#3}(#1;\,#4)}
\newcommand{\Ob}[2]{\operatorname{Ob}(#1;\,#2)}
\newcommand{\Coc}[1]{\operatorname{Coc}(#1)}
\newcommand{\Rel}[1]{\operatorname{Rel}(#1)}
\newcommand{\pd}[2]{\Sigma_{#2}(#1)}
\newcommand{\inab}[1]{I_{#1}}
\newcommand{\cj}[1]{\bar{#1}}
\newcommand{\Nm}{\operatorname{N}}
\newcommand{\nlen}{\operatorname{nlen}}
\newcommand{\cO}{\mathcal{O}}
\newcommand{\fP}{\mathfrak{P}}
\newcommand{\fq}{\mathfrak{q}}
\newcommand{\HT}{\mathcal{H}}
\newcommand{\GL}{\mathrm{GL}}
\newcommand{\tr}{\operatorname{tr}}
\newcommand{\rk}{\operatorname{rk}}
\newcommand{\Art}{\operatorname{Art}}
\newcommand{\av}{\operatorname{av}}
\newcommand{\Zf}{Z_{\mathrm{f}}}
% the data (as in the note on admissible data)
\newcommand{\Dm}[1]{m(#1)}
\newcommand{\Dr}[1]{r(#1)}
\newcommand{\kD}[1]{k(#1)}
\newcommand{\tD}[1]{t(#1)}
\newcommand{\fD}[1]{f(#1)}
\newcommand{\ED}[1]{E(#1)}
\newcommand{\Dg}[2]{g_{#1}(#2)}
\newcommand{\gD}[1]{g(#1)}
\newcommand{\lvl}[1]{N(#1)}
\newcommand{\hhD}[1]{h_2(#1)}
\newcommand{\hD}[1]{h(#1)}
\newcommand{\bone}[1]{b_1(#1)}
\newcommand{\btwo}[1]{b_2(#1)}
\newcommand{\wt}[3]{w_{#1,#2}(#3)}
\newcommand{\ordch}[1]{\operatorname{ord}_{\mathrm{ch}}(#1)}
\newcommand{\Dmu}[1]{\mu(#1)}
\newcommand{\ND}[1]{N_1(#1)}
\newcommand{\thick}[1]{\operatorname{thick}(#1)}
\newcommand{\inp}[2]{\statement{ana}{\textbf{#1.}\enspace #2}}

\newcommand{\Sha}{\mathrm{Sha}}

% where each declaration of the concordance is treated
\cs{IsConjSelfDual}{Definition~\ref{d:recall}}
\cs{crudeUnit}{Definition~\ref{d:recall}}
\cs{SelmerLengthLe}{Definition~\ref{d:recall}}
\cs{IsColumn}{Definition~\ref{d:recall}}
\cs{IsRow}{Definition~\ref{d:recall}}
\cs{chiOneInv}{Definition~\ref{d:recall}}
\cs{normalisedRep}{Definition~\ref{d:recall}}
\cs{twistedLabel}{Definition~\ref{d:recall}}
\cs{IsApt}{Definition~\ref{d:recall}}
\cs{tupleProj}{Definition~\ref{d:recall}}
\cs{HasRow}{Definition~\ref{d:recall}}
\cs{IsLabelledColumn}{Definition~\ref{d:recall}}
\cs{obeyingCocycles}{Definition~\ref{d:recall}}
\cs{RowBlockIndex}{Definition~\ref{d:recall}}
\cs{blockRep}{Definition~\ref{d:recall}}
\cs{IsAuxiliary}{Definition~\ref{d:recall}}
\cs{evenClasses}{Definition~\ref{d:recall}}
\cs{IsNewPlaces}{Definition~\ref{d:recall}}
\cs{withConjInv}{Definition~\ref{d:recall}}
\cs{IsPrimeTriangle}{Definition~\ref{d:recall}}
\cs{IsAuxTwist}{Definition~\ref{d:recall}}
\cs{existscycCharpowsubonedvd}{Lemma~\ref{l:cycdvd}}
\cs{denominatorLerowBudget}{Theorem B}
\cs{crysRels}{Definition~\ref{d:recall}}
\cs{blockLoss}{Definition~\ref{d:recall}}
\cs{cmSign}{Definition~\ref{d:recall}}
\cs{IsRowType}{Definition~\ref{d:recall}}
\cs{lambdaLoss}{Definition~\ref{d:recall}}
\cs{loss}{Definition~\ref{d:recall}}
\cs{innerAtheckePolyOnForms}{Lemma~\ref{l:heckefacts}}
\cs{groupAlgAction}{Definition~\ref{d:den}}
\cs{IsTested}{Definition~\ref{d:den}}
\cs{IsTestForm}{Definition~\ref{d:den}}
\cs{DenominatorLe}{Definition~\ref{d:den}}
\cs{isOfFinOrderprodpsidet}{Lemma~\ref{l:heckefacts}}
\cs{innerAtformAction}{Lemma~\ref{l:inneraction}}
\cs{rangegroupAlgActionmeminvtSubmodule}{Lemma~\ref{l:heckefacts}}
\cs{heckePolyOnFormscomm}{Lemma~\ref{l:heckecomm}}
\cs{innerAtgroupAlgAction}{Lemma~\ref{l:heckefacts}}
\cs{normSqAtstarProjectionle}{Lemma~\ref{l:normproj}}
\cs{IsSplitOffLevel}{Definition~\ref{d:hecke}}
\cs{heckeAlgebra}{Definition~\ref{d:hecke}}
\cs{heckeOperator}{Definition~\ref{d:hecke}}
\cs{eigenPart}{Definition~\ref{d:hecke}}
\cs{isSphericalWithofmemeigenPart}{Lemma~\ref{l:spherical}}
\cs{occurringLabels}{Definition~\ref{d:hecke}}
\cs{dvdcardresFieldsubone}{Lemma~\ref{l:resfield}}
\cs{existsheckeprojeigenPart}{Lemma~\ref{l:heckeproj}}
\cs{commuteheckeAlgebraprincipalRep}{Lemma~\ref{l:commute}}
\cs{heckeOnGood}{Definition~\ref{d:hecke}}
\cs{principalRepmemeigenPart}{Lemma~\ref{l:heckefacts}}
\cs{finiteeigenPart}{Lemma~\ref{l:heckefacts}}
\cs{iSupeigenPart}{Lemma~\ref{l:eigenparts}}
\cs{thetamuleqthetaformAction}{Lemma~\ref{l:thetaaction}}
\cs{existsthetatoprincipalForms}{Lemma~\ref{l:thetaprincipal}}
\cs{auxLift}{Definition~\ref{d:lifts}}
\cs{RankOneIndex}{Definition~\ref{d:lifts}}
\cs{goodLift}{Definition~\ref{d:lifts}}
\cs{IsRamified}{Definition~\ref{d:pairs}}
\newcommand{\Mom}[2]{M_\omega(#1,#2)}           % as in the note on admissible data
\newcommand{\tupleIdem}[2]{e^{\otimes}_{#1,#2}}   % as in the note on primitive labels
\cs{IsRankOne}{Definition~\ref{d:pairs}}
\cs{B}{Definition~\ref{d:pairs}}
\cs{IsFitFamily}{Theorem A}
\cs{existsisFitFamily}{Theorem A}
\cs{AbsGal}{Definition~\ref{d:recall}}
\cs{placesAboveG}{Definition~\ref{d:recall}}
\cs{badPlaces}{Definition~\ref{d:recall}}
\cs{invAvatar}{Definition~\ref{d:recall}}
\cs{centralRep}{Definition~\ref{d:central}}
\cs{twistedCentralOp}{Definition~\ref{d:central}}
\cs{innertwistedCentralOp}{Lemma~\ref{l:central}}
\cs{twistedCentralOpkernelForm}{Lemma~\ref{l:central}}
\cs{twistedCentralOptrivialLift}{Lemma~\ref{l:central}}
\cs{rankonebound}{Lemma~\ref{l:rankone}}
\cs{eqzeroofmemevenClasses}{Lemma~\ref{l:evenzero}}
\cs{IsSplitPlace}{Definition~\ref{d:law}}
\cs{primeBelowPlus}{Definition~\ref{d:law}}
\cs{IsLaw}{Definition~\ref{d:law}}
\cs{primeBelowPlusstar}{Lemma~\ref{l:heckefacts}}
\cs{existsisLaw}{Lemma~\ref{l:law}}
\cs{ribetNormaliser}{Definition~\ref{d:recall}}
\cs{rowBudget}{Definition~\ref{d:rowbudget}}
\cs{memselmerModule}{Lemma~\ref{l:obeysel}}
\cs{isTraceIdealwithConjInv}{Lemma~\ref{l:traceconj}}
\cs{sumblockLossle}{Lemma~\ref{l:blockloss}}
\cs{ribetcolumns}{Theorem C}
\cs{HasTriangleCost}{Definition~\ref{d:recall}}
\cs{IsTraceIdeal}{Definition~\ref{d:recall}}
\cs{IsRibetElement}{Definition~\ref{d:recall}}
\cs{existsrowforalltwists}{Lemma~\ref{l:rowtwists}}
\cs{existshasRowauxPair}{Theorem E}
\cs{existsisPrimeTriangle}{Theorem D}
\cs{existsmFoldTwistselmerLengthLe}{Theorem F}

\begin{document}
\color{ink}

{\raggedright
{\sffamily\bfseries\fontsize{7.2}{9}\selectfont\color{leancol}\addfontfeatures{LetterSpace=8}READBACK NOTE 33\par}
\vspace{5pt}
{\fontsize{21.5}{25.5}\selectfont\bfseries A fit family of twists exists: Ribet's method with tied primes, triangular Taylor–Wiles primes and the denominator of the projection to theta lifts\par}
\vspace{5pt}
{\fontsize{12.4}{16}\selectfont\itshape\color{inksoft} One unit of the Lean artifact \ArtifactName, read back into
mathematics\par}}
\vspace{8pt}
{\color{rulegrey}\hrule height 0.6pt}
\vspace{7pt}
\thispagestyle{empty}

\statement{arith}{\phantomsection\label{t:A}\textbf{Theorem A} (a fit family of twists exists).\enspace
\emph{Let $P$ be parameters and $D$ admissible data for~$P$. Then there is a family
\m{\theta=(\theta_{\ell,\iota})} of $m$-fold twists of~$D$, one for each pair \m{(\ell,\iota)} of a prime number $\ell$
and a field isomorphism \m{\iota\colon\CC\to\Qlb}, with the following three properties (the family is \emph{fit
for~$D$}).}
\begin{enumerate}[label=\textup{(F\arabic*)},leftmargin=3em,itemsep=2pt]
\item \emph{For every \m{(\ell,\iota)}, the prime $\ell$ is prime to the product \m{\Nm(Q(\theta_{\ell,\iota}))} of the norms
of the places of \m{Q(\theta_{\ell,\iota})}.}
\item \emph{For every lift index \m{1\le j\le\Dm{D}} there are real numbers \m{(\eta_l)_{l\in\NN}} such that
\m{\sum_{l\in T}\eta_l\le\Bzero{D}} for every finite set $T$ of prime numbers, and such that for every
\m{(\ell,\iota)}, with \m{\vartheta=\theta_{\ell,\iota}^{\,j}}:}
\[
\DenominatorLe{\iota}{O_D(\vartheta,j)}{e_0(D)}{\operatorname{im}\goodLift{D}{\vartheta}{j}}{\eta_\ell},
\]
\emph{and, for every pair $p$ of~$D$ that is of rank one for~$j$, with $r_p$ the index set of the rows
of~\m{O_p(\vartheta)},}
\[
\DenominatorLe{\iota}{O_p(\vartheta)}{\tupleIdem{r_p}{X-1}}{\CC\cdot\Xi_p(\vartheta)}{\eta_\ell}.
\]
\item \emph{For every lift index \m{1\le j\le\Dm{D}}, every pair $p$ of~$D$ and every sign \m{s\in\{1,-1\}}: if $p$ is
ramified, then the family \m{(\ell,\iota)\mapsto\varrho_p(\theta_{\ell,\iota}^{\,j},s)^{-1}} of complex numbers has height at
most \m{x(D)^{999}}.}
\end{enumerate}}
\leanline{\lean{existsisFitFamily}\sep \lean{IsFitFamily}}

Here \m{\Bzero{D}=\operatorname{budget}(D)/(\Dm{D}\kD{D})}; \m{O_D(\vartheta,j)} is the $j$-th auxiliary pair of
the twist~$\vartheta$, \m{e_0(D)} the good idempotent and \m{\goodLift{D}{\vartheta}{j}} the adjoint theta lift on the
good forms (Definition~\ref{d:lifts}); \m{O_p(\vartheta)} is the auxiliary pair of the $L$-value data of $p$
and~$\vartheta$, of type \m{(r_p,n)} with \m{\#r_p=1}, \m{\tupleIdem{r_p}{X-1}} the idempotent of the tuple with all
labels \m{X-1} (written \m{\tupleIdem{r}{c}} for an index set $r$ and a tuple~$c$, as in the note \emph{Primitive
labels}) and \m{\Xi_p(\vartheta)} the lift of the trivial character; \m{\DenominatorLe{\iota}{O}{e_0}{Z}{\beta}} is the bound on the denominator of the orthogonal projection
to~$Z$ of Definition~\ref{d:den}; \m{\varrho_p(\vartheta,s)} is the ratio of $L$-values
\[
\varrho_p(\vartheta,s)=\frac{\bigl|L^{(\mathrm{Ch}\cup Q(\vartheta))}\bigl(z_p,(\nu_p\vartheta^{s})^{-1}\bigr)\bigr|}{\bigl|L^{(\mathrm{Ch})}\bigl(z_p,\nu_p^{-1}\bigr)\bigr|},
\]
where the upper index names a set of finite places at which, and at whose conjugates, the Euler factors are
removed, and \m{\mathrm{Ch}} is the set of the chosen places; and a family \m{y=(y_{\ell,\iota})} of complex numbers
\emph{has height at most~$c$} if there are \m{H\in\CC^{\times}}, a real number $\alpha$ and real numbers
\m{(\eta_l)_{l\in\NN}} with \m{|\sigma(H)|\le e^{\alpha}} for all field automorphisms $\sigma$ of~$\CC$,
\m{|\iota(H\,y_{\ell,\iota})|\le e^{\eta_\ell}} for all \m{(\ell,\iota)}, and \m{\alpha+\sum_{l<M\text{ prime}}\eta_l\le c} for
all large~$M$ (note \emph{Heights}). The letter \m{\eta_l} is that of the overview of this handout, for the
budget of (F2) and for the local bounds of a height; the lift index is written~$j$ here and in the formal text,
and~$i$ in the overview.

Theorem~A is the statement \emph{Twists} of the overview of this handout. The following points of its reading
were checked against the formal text by the second reader of this write-up, whose reading this note follows.
The twist \m{\theta_{\ell,\iota}} may depend on $\iota$ at a fixed~$\ell$. The whole family is chosen before the lift index,
the pair and the sign are quantified. In (F2) the numbers \m{\eta_l} depend on~$j$ only (the proof even gives the
same for all~$j$), the same \m{\eta_\ell} serves every $\iota$ over~$\ell$, both the image of the good lift and
every rank-one line, and the twist in the bound is the $j$-th power \m{\theta_{\ell,\iota}^{\,j}}, whose auxiliary support
is that of~\m{\theta_{\ell,\iota}}. In (F3) it is the \emph{inverse} ratio whose height is bounded, with one multiplier
$H$ for all embeddings.

\textbf{Summary.} A denominator of the orthogonal projection of integral forms to the image of a theta lift is
a congruence between a lift and a form that is not a lift. This note proves that for each pair \m{(\ell,\iota)}
there is a twist, with its first $m$ powers, at which these denominators are within a budget whose sum over all
primes is at most~\m{\Bzero{D}} (Theorem~A). The chain of the proof: the Galois representations of the lifts are
a row plus a string of powers of the cyclotomic character, the row serving all twists (Theorem~E); one
anticyclotomic twist makes the Selmer modules of all rows short and keeps the $L$-values in check (Theorem~F,
from the second note); Ribet's lemma, fed with a prime triangle (Theorem~D, by reciprocity laws) and the bound
for the odd part of Tate--Shafarevich groups, gives a Hecke element of bounded valuation (Theorem~C); and that
element bounds the denominator (Theorem~B).

\statement{leancol}{\fontsize{9.6}{13.2}\selectfont\textbf{What this note is.} The artifact \ArtifactName{} is a
body of Lean code: the compact edition of a formal proof, due to L.~Alpöge with the AI system Claude, of the
Riemann hypothesis and of its generalisation to Dirichlet $L$-functions. That headline is a claim. The authors
report that every file compiles and that Lean's kernel accepted the proofs; an independent replay is recorded in
the handout's record of checks when it reports. No mathematician has read the argument through, and nothing here
says that the Riemann hypothesis is proved. This note is the third and last on the unit \emph{The choice of the
twists}: it treats the declarations of its last nine files and recalls, from the notes \emph{Columns and rows as
Galois representations} and \emph{Short Selmer groups after one anticyclotomic twist}, the definitions it needs;
together \UnitInNote{} of the unit's \UnitMembers{} declarations appear here. All statements are complete. The
proofs of the short lemmas are complete, and the proofs of Theorems A, E and~F follow the formal proofs step by
step; the proofs of Theorems B, C, D and of Lemma~\ref{l:law} are outlines of formal proofs of a hundred to a
hundred and fifty lines, and say so. Results of the background library and of other units are quoted as Inputs with their owners. Every
statement was compared with its Lean text, argument by argument. The Lean text read is \PinText; Appendix~A is
the concordance of \ConcDecls{} declarations.}

\sect{The other main statements}{s:intro}
The conventions and the recalled notions are in §\ref{s:back}, the definitions of this note in §\ref{s:defs}.
Throughout, $\ell$ is a prime number and \m{\iota\colon\CC\to\Qlb} a fixed isomorphism of fields.

\statement{arith}{\phantomsection\label{t:B}\textbf{Theorem B} (the denominator of the projection to lifts).\enspace
\emph{Let $P$ be parameters, $D$ admissible data for~$P$, and $\vartheta$ a twist of~$D$ with
\m{\IsAuxTwist{\badPlaces{D}{\ell}}{Q(\vartheta)}{\invAvatar{\vartheta}{\cdot}}{\crudeUnit{\ell}{D}}}. Let $r$, $n$ be finite sets, $O$
an auxiliary pair of type \m{(r,n)} over~\m{\ED{D}}, $e_0$ an element of the complex group algebra of the group
\m{\prod_i\GL_r(\F_{w_i})} of the chosen places of~$O$, and $Z$ a subspace of~\m{\mathcal{F}(O)}. Suppose that $Z$ is
contained in the span of the set of the \m{y\in Z} with the following two properties:}
\begin{enumerate}[label=\textup{(\alph*)},leftmargin=2.3em,itemsep=2pt]
\item \emph{there is a family $c$ of polynomials with \m{y\in\operatorname{im}\tupleProj{O}{c}},
\m{\IsApt{O}{c}{\vartheta}} and \m{\operatorname{im}\tupleProj{O}{c}\subseteq\operatorname{im}\groupAlgAction{O}{e_0}};}
\item \emph{there are Hecke data $z$ such that $y$ lifts~$z$ on~$O$, a homomorphism
\m{R\colon\AbsGal{D}\to\GL_r(\Qlb)} and an auxiliary pair $V$ of type \m{(r,n)} over~\m{\ED{D}} with \m{\HasRow{V}{z}{R}}
and \m{\mu_O=\mu_V\cdot\vartheta}, and elements \m{\gamma\in\AbsGal{D}}, \m{s'\in\Zlb} with
\m{\SelmerLengthLe{\invAvatar{\vartheta}{\cdot}}{R}{Q(\vartheta)}{\lfloor\Dr{D}\cdot\rowBudget{D}{\ell}\rfloor+1}{\gamma}{s'}}.}
\end{enumerate}
\emph{Then \m{\DenominatorLe{\iota}{O}{e_0}{Z}{\rowBudget{D}{\ell}\cdot\log\ell}}.}}
\leanline{\lean{denominatorLerowBudget}}

\statement{arith}{\phantomsection\label{t:C}\textbf{Theorem C} (Ribet's method for columns).\enspace
\emph{Let $P$, $D$ be as in Theorem~B, $K$ a finite totally ordered set, $Y$, $n$, $r$ finite sets, \m{\nu\in\NN},
\m{w=(w_k)_{k\in K}} finite places of~\m{\ED{D}}, \m{\tau=(\tau_k)_k}, \m{\phi=(\phi_k)_k} families in~\m{\AbsGal{D}},
\m{u=(u_x)_{x\in Y}} homomorphisms \m{u_x\colon\AbsGal{D}\to\GL_n(\Qlb)}, $V$ an auxiliary pair of type \m{(r,n)}
over~\m{\ED{D}} with rank difference \m{t_V=\#n-\#r}, $\vartheta$ a twist and $c$ a family of polynomials with
\m{\IsApt{V}{c}{\vartheta}}. Let $\HT'$ be a type, \m{\psi\colon\AbsGal{D}\to\Qlb} a multiplicative homomorphism and $Q$
a finite set of finite places with \m{\IsAuxTwist{\badPlaces{D}{\ell}}{Q}{\psi}{\crudeUnit{\ell}{D}}}. Let
\m{\mathrm{ty}\colon K\to\{e:0<e<t_V\}} be monotone, \m{a\in\AbsGal{D}} and \m{h\in\NN} with
\m{h\le v_\ell(\ordch{D})} and
\m{\IsPrimeTriangle{\nu}{w}{\tau}{\phi}{a}{t_V}{\mathrm{ty}}{\hhD{D}h}{\hhD{D}^{22}h}}. Let $\mathcal{S}$ be a
\m{\Zlb}-subalgebra of the functions \m{Y\to\Qlb[X_k:k\in K]}. Suppose that for every \m{x\in Y}}
\[
\IsLabelledColumn{V}{t_V}{\HT'}{(\twistedLabel{V}{c}{i}(X-1)^{t_V})_i}{Q\cup\cj{Q}\cup w(K)\cup\cj{w(K)}}{u_x}
\]
\emph{and \m{(u_x(g)-\psi(g))(u_x(g)-1)=0} for all \m{g\in\inab{Q\cup\cj{Q}}}. Let \m{\mathcal{X}\subseteq Y}, and for
\m{x\in\mathcal{X}} let $\HT_x$ be a type, \m{R_x\colon\AbsGal{D}\to\GL_r(\Qlb)} a homomorphism,
\m{\gamma_x\in\AbsGal{D}} and \m{y_x\in\Zlb}; let \m{m_0\in\NN} with \m{\lfloor\Dr{D}\cdot\rowBudget{D}{\ell}\rfloor<m_0}.
Suppose that for every \m{x\in\mathcal{X}}: \m{\IsRow{D}{t_V}{\HT_x}{R_x}},
\m{\SelmerLengthLe{\psi}{R_x}{Q}{m_0}{\gamma_x}{y_x}}, and
\m{\tr u_x(g)=\psi(g)\tr R_x(g)+\sum_{e<t_V}\cyc(g)^{-e}} for all~$g$. Let $\mathcal{J}$ be an
$\mathcal{S}$-submodule of the functions \m{Y\to\Qlb[X_k:k\in K]} and \m{\varepsilon=(\varepsilon_k)_{k\in K}}
polynomials over~\m{\Zlb} with \m{\IsTraceIdeal{\mathcal{J}}{u}{\ell^{\nu}}{\tau}{\phi}{\mathrm{id}}{\varepsilon}} and
\m{\varepsilon_k(1)=1} for all~$k$. Let \m{\pi\in\mathcal{X}} and $\theta'$ a finitely supported function
\m{\AbsGal{D}\to\NN}, and let \m{\beta(x)=\det\bigl(\sum_g\ribetNormaliser{\gamma_\pi}{a}{y_\pi}{\ell^{m_0+\hhD{D}h}}{\theta'}{g}\,(u_x(g)-1)\bigr)}
for \m{x\in Y}. Then there is \m{d'\in\Qlb} with}
\[
\IsRibetElement{\mathcal{J}}{\beta}{\pi}{d'}\qquad\text{\emph{and}}\qquad\vl\bigl(\ordch{D}^{3}\cdot d'\bigr)\le\#r\cdot\rowBudget{D}{\ell}.
\]}
\leanline{\lean{ribetcolumns}}

\statement{arith}{\phantomsection\label{t:D}\textbf{Theorem D} (triangular Taylor--Wiles primes).\enspace
\emph{Let $E$ be a CM field as in the conventions, \m{t,N,y,h} natural numbers with \m{N\ge1}, \m{\gamma\in\GE{E}}, and
$W$ a set of Frobenius data of~$E$. Suppose that the image of~$E$ in~\m{\Qb} lies in the cyclotomic field of
conductor~$N$; that \m{N+t<y}; that \m{\cyc(\gamma)^{a}-1} divides \m{\ell^{h}} in~\m{\Zl} for all \m{0<a<t}; that
\m{\ell\le y^2} or \m{h=0}; that $E$ has at least four infinite places; and that for all \m{\nu\in\NN} and
\m{g\in\GE{E}} there is \m{\delta\in\GE{E}} with \m{\cyc(\delta)=1} such that $W$ contains most Frobenius data
near \m{\delta^{\ell^{\nu}}g}. Then there are a finite totally ordered set $J$ and a monotone map
\m{\mathrm{ty}\colon J\to\{e:0<e<t\}} such that for every \m{\nu>0} there is a system
\m{(w_j,\fP_j,\phi_j)_{j\in J}} of Taylor--Wiles places of~$E$ for the modulus~\m{\ell^{\nu}} with}
\[
\IsPrimeTriangle{\nu}{(w_j)_j}{(\tau_{\fP_j})_j}{(\phi_j)_j}{\gamma^{-1}}{t}{\mathrm{ty}}{yh}{y^{22}h}
\]
\emph{and \m{(w_j,\fP_j,\phi_j)\in W} for every~$j$.}}
\leanline{\lean{existsisPrimeTriangle}}

\statement{arith}{\phantomsection\label{t:E}\textbf{Theorem E} (the good eigensystems have rows).\enspace
\emph{Let $P$, $D$ be as in Theorem~B, $p$ a function from the pairs of a finite place of~\m{\ED{D}} and a natural
number to~$\CC$, and $j$ a lift index \m{(1\le j\le\Dm{D})}. Suppose that the space of the good forms on
\m{W_{\Dm{D}}(D)} meets \m{\eigenPart{D}{W_{\Dm{D}}(D)}{p}} non-trivially. Then there is a homomorphism
\m{T\colon\AbsGal{D}\to\GL_r(\Qlb)}, with $r$ the index set of the rows of~$D$, such that
\m{\HasRow{O_D(1,j)}{p}{T}}.}}
\leanline{\lean{existshasRowauxPair}}

\statement{arith}{\phantomsection\label{t:F}\textbf{Theorem F} (one $m$-fold twist for finitely many rows).\enspace
\emph{Let $P$, $D$ be as in Theorem~B, $\ell$ a prime number and \m{\iota\colon\CC\to\Qlb} an isomorphism of
fields. Let $I$ be a finite set, \m{(r_i)_{i\in I}} finite sets, \m{R_i\colon\AbsGal{D}\to\GL_{r_i}(\Qlb)}
homomorphisms such that for every~$i$ there are $t_i$ and a type $\HT_i$ with \m{\IsRow{D}{t_i}{\HT_i}{R_i}},
and \m{M\in\NN}. Then there is an $m$-fold twist $\tau$ of~$D$ such that for every lift index~$j$
\m{(1\le j\le\Dm{D})}, with \m{\vartheta=\tau^{j}}:}
\begin{enumerate}[label=\textup{(\arabic*)},leftmargin=2.3em,itemsep=2pt]
\item \emph{\m{\IsAuxTwist{\badPlaces{D}{\ell}}{Q(\vartheta)}{\invAvatar{\vartheta}{\cdot}}{\crudeUnit{\ell}{D}}};}
\item \emph{for every $i$ there are \m{\gamma\in\AbsGal{D}} and \m{y\in\Zlb} with
\m{\SelmerLengthLe{\invAvatar{\vartheta}{\cdot}}{R_i}{Q(\vartheta)}{M}{\gamma}{y}};}
\item \emph{for every pair $p$ of~$D$ and sign~$s$: if $p$ is ramified, then
\m{\bigl|\iota\bigl(\varrho_p(\vartheta,s)\bigr)^{-1}\bigr|\le\exp\bigl(9\,\crudeUnit{\ell}{D}\log\ell\bigr)}.}
\end{enumerate}}
\leanline{\lean{existsmFoldTwistselmerLengthLe}}

\sect{Conventions and recalled notions}{s:back}
\begin{conventions}\phantomsection\label{c:conventions}
The conventions are those of the first note of the unit. The natural numbers are \m{0,1,2,\dots}. $\ell$ is a
prime number (also $l$, when several primes occur); \m{\Ql}, \m{\Zl}, \m{\Qlb}, \m{\Zlb}, \m{\Cl} and
\m{\Gl=\Gal(\Qlb/\Ql)} are as usual, with \m{|\ell|=1/\ell}; \m{\iota\colon\CC\to\Qlb} is a fixed isomorphism of fields;
families indexed by the pairs \m{(\ell,\iota)} carry both as indices. \m{\Qb} is a fixed algebraic closure
of~$\mathbb{Q}$ with a fixed complex conjugation~$\mathsf{c}$. A \emph{CM field as in the conventions} is a CM
number field $E$ with a fixed embedding in~\m{\Qb} that contains a primitive root of unity of order~$5040$, with
its CM type and its Weil data (Chapter~VII); \m{\GE{E}=\Gal(\Qb/E)}; $E^{+}$ is its maximal totally real subfield.
$P$ denotes parameters and $D$ admissible data for~$P$; \m{\ED{D}} is a CM field as in the conventions. The
letters $n$, $r$, $d$, $\sigma$, $I$, $J$, $K$, $Y$, $X$, $T$, when they index matrices or families, denote finite sets;
\m{\#n} is the number of elements. A \emph{coefficient ring} is the ring of integers $\cO$ of a finite Galois
extension of~\m{\Ql} inside~\m{\Qlb}; ``for all large~$\cO$'' means: for all coefficient rings containing a
suitable one. $V$ also denotes a hermitian space with $s$ a parameter of forms on it, as in Chapter~VII, and
then \m{\mathcal{M}_D(V)} is the space of the forms of principal level at the chosen places on~$V$, which is
assumed finite-dimensional where this is said.
\end{conventions}

\subhead{2.1. Background.} The background notions are those of §2 of the first note and of §2.1 of the second:
finite places, norms \m{\Nm(w)}, conjugates \m{\cj{w}}, \emph{split} places (\m{\cj{w}\ne w}), the sets
\m{\pd{E}{M}}; places $\fP$ of~\m{\Qb}, the chosen place \m{\fP_w} over~$w$, groups \m{G_{\fP}}, \m{I_{\fP}},
\m{\inab{S}}, Frobenius elements and tame generators \m{\tau_{\fP}}; \m{g^{\mathsf c}} and \m{g^{[\pm]}};
\m{\cyc} and \m{\cyc_M}; characters from $\mathrm{U}(1)$; relations, \m{\Ob{L}{B}}, Selmer modules
\m{\Sel{\varrho}{S}{L}{B}}; lengths \m{\len} and \m{\nlen_{\cO}}; the valuation \m{\vl}; Tate twists \m{\cO(-e)};
lattices of representations; Galois lattices and their classes; characters of~$\mathrm{U}(1)$ with weight and
tame part, avatars; tame characters, twists and pairs \m{p=(\nu_p,z_p)} of~$D$, with the weight \m{-c_p} of~\m{\nu_p}
and its tame part~$\psi_p$; chosen places; auxiliary pairs $O$ of type \m{(r,n)} with their chosen places \m{w_i},
characters \m{\mu_O}, \m{\psi_O}, \m{a_O}, forms \m{\mathcal{F}(O)} and Hecke operators \m{H_w(X)}; labels and the
idempotents \m{\tupleIdem{r}{c}} of the tuples $c$ of labels, in the group algebra of \m{\prod_i\GL_r(\F_{w_i})}; eigenforms, Hecke polynomials and attached representations; base levels. Further:
\begin{itemize}[leftmargin=1.6em,itemsep=1pt]
\item A \emph{Frobenius datum} of~$E$ is a triple \m{(w,\fP,\phi)} of a finite place of prime norm, a place
of~\m{\Qb} over it and an element $\phi$ whose inverse is an \m{\Nm(w)}-Frobenius at~$\fP$ (a geometric Frobenius
element). A set $W$ of Frobenius data \emph{contains most Frobenius data} if every non-empty open subset
of~\m{\GE{E}} contains a non-empty open subset $U$ such that, for all but finitely many finite places~$w$, every
Frobenius datum \m{(w,\fP,\phi)} with \m{\phi\in U} belongs to~$W$; such sets form a filter. The set $W$
\emph{contains most Frobenius data near~$g$}, for \m{g\in\GE{E}}, if it contains the intersection of a set that
contains most Frobenius data with the set of the data whose element $\phi$ lies in some neighbourhood of~$g$. A
statement holds \emph{for most Frobenius data} (near~$g$) if the set of the data for which it holds has this
property. A system of
\emph{Taylor--Wiles places} of~$E$ for the modulus~$M$, indexed by~$J$, is a family of Frobenius data
\m{(w_j,\fP_j,\phi_j)} with the $w_j$ pairwise distinct, no \m{\cj{w_i}} among them, and \m{M\mid\Nm(w_j)-1}.
\item An \emph{$m$-fold twist} of~$D$ is a twist $\vartheta$ such that \m{\psi_{\vartheta,w}^{\,j}\ne1} for all
\m{w\in Q(\vartheta)} and \m{1\le j\le\Dm{D}}; its powers \m{\vartheta^{j}}, \m{1\le j\le\Dm{D}}, are twists with the
same set \m{Q(\vartheta)}; \m{\Nm(Q(\vartheta))} is the product of the norms of its places.
\item For a form on~$O$ the \emph{integrality at a ring homomorphism \m{\CC\to\Cl}} is that of the overview of
this handout (the values lie in the weight lattice at each point); the \emph{integrality at~$\iota$} is the
integrality at $\iota$ followed by \m{\Qlb\subseteq\Cl}. \m{\langle y,y'\rangle} is the Petersson product
of~\m{\mathcal{F}(O)}, \m{\langle y,y'\rangle_\sigma} its form at an automorphism $\sigma$ of~$\CC$ and
\m{\|y\|_\sigma^2=\langle y,y\rangle_\sigma}; \m{\mathrm{pr}_Z} is the orthogonal projection to a subspace~$Z$.
\item \m{\theta_O} is the theta lift from the forms of~$O$ to forms on its hermitian space of rank~\m{\#r}
(Chapter~VII); a form \emph{lifts} Hecke data~$z$ on~$O$ if at the split places off the places of~$O$ it is
spherical with the Hecke polynomial of the string times that of~$z$ (Chapter~VII; the notion is used in
Definition~3.21 of the first note).
\item For admissible data~$D$: \m{W_j(D)} is the hermitian space of the $j$-th single lift, with index set of
rank \m{\Dr{D}}; \m{e_0(D)} the good idempotent in the group algebra of \m{\mathsf{G}(D)=\prod_w\GL_r(\F_w)} over the
chosen places, the sum of the \m{\tupleIdem{r}{c}} over the good tuple families~$c$ ($r$ the index set of the rows); the \emph{good forms} on a space are the image
of~\m{e_0(D)} in its forms of principal level; \m{O_D(\vartheta,j)} the $j$-th auxiliary pair of a twist;
\m{O_p(\vartheta)}, \m{\Xi_p(\vartheta)} the auxiliary pair and the lift of the trivial character of the $L$-value data
of a pair $p$ and a twist; \m{x(D)} the first scale and \m{\operatorname{budget}(D)} the budget (Chapters III, V, VII).
\item \m{\Sha^{\mathrm{odd}}(\ell,E,\cO,e,\nu)} is the module of §2.1 of the second note.
\end{itemize}

\begin{definition}[notions recalled from the first two notes]\label{d:recall}
The following are defined in the notes \emph{Columns and rows as Galois representations} (numbers 3.x) and
\emph{Short Selmer groups after one anticyclotomic twist} (numbers 3.x of that note, marked~II).
\begin{enumerate}[label=\textup{(\arabic*)},leftmargin=2.3em,itemsep=2pt]
\item (3.1, 3.2, 3.5, 3.6) \m{\obeyingCocycles{u}{\varrho}{a}}: the cocycles for~$\varrho$ modulo~$a$ that obey the
relations among the \m{u_x(g)-1}. \m{\RowBlockIndex{r}{i}} is $r$ for \m{i=\ast} and a point for \m{i=e<t};
\m{\blockRep{T}{i}{\cdot}} is $T$ for \m{i=\ast} and \m{\cO(-e)} for \m{i=e}. \m{\crysRels{u}} is the set of the
crystalline relations of~$u$, and \m{\blockLoss{u}{w}{M}} the length of
\m{\obeyingCocycles{w}{u}{\ell^{M}}} modulo its elements that obey \m{\crysRels{u}}.
\item (3.7--3.9) \m{\cmSign{\sigma}=\pm1} according as the complex embedding $\sigma$ belongs to the CM type or not.
\m{\IsRowType{D}{r}{t}{\HT}}: the conditions (R1)--(R4) on a type. \m{\lambdaLoss{h}{b}{l}} is the exponent of $l$
in the integer \m{(2b)!\,(\prod_{b-30h^3<k\le b+30h^3}k)^{h^3}}, and
\[
\loss{D}{l}=10^{6}\,\frac{\Dm{D}}{\kD{D}}\,\bigl(\bone{D}\,v_l(\gD{D})+\lambdaLoss{\hhD{D}}{\bone{D}}{l}\bigr).
\]
\item (3.10, 3.11, 3.13) \m{\AbsGal{D}=\Gal(\Qb/\ED{D})}. The set \m{\placesAboveG{D}} consists of the places
above~\m{\gD{D}}, and \m{\badPlaces{D}{\ell}} of the places containing \m{\ell\gD{D}\lvl{D}}. For a tame character
$\vartheta$, \m{\invAvatar{\vartheta}{g}} is the value at~$g$ of the avatar of~\m{\vartheta^{-1}}. Finally
\m{\crudeUnit{\ell}{D}=v_\ell(\ordch{D})+2^{-\ell}}.
\item (3.12, 3.14--3.16) \m{\IsConjSelfDual{t}{R}}: $R$ is continuous with non-degenerate trace form and
\m{\tr R(g^{\mathsf c})=\cyc(g)^{-(t-1)}\tr R(g^{-1})}. \m{\SelmerLengthLe{\psi}{R}{Q}{m}{\gamma}{z}}: for all large~$\cO$ some
lattice $u$ of~$R$ and \m{\psi'=\psi} make \m{\varrho=\psi'u} scalar \m{=z} at~$\gamma$ with \m{z-1} a unit,
\m{\cyc(\gamma)=1}, and the cocycles vanishing at~$\gamma$ in
\m{\Sel{\varrho}{\pd{\ED{D}}{\ell}\cup\placesAboveG{D}\cup Q\cup\cj{Q}}{\crysRels{\varrho}}{\cO/\ell^{m}}} have normalised length
at most \m{(10\fD{D}\#n+1)\crudeUnit{\ell}{D}}. \m{\IsColumn{D}{t}{\HT}{S}{R}} and \m{\IsRow{D}{t}{\HT}{R}}: columns of
level~$S$ and rows.
\item (3.17--3.22) \m{\chiOneInv{O}} the normalising character and \m{\normalisedRep{O}{R}=\chiOneInv{O}\cdot R};
\m{\twistedLabel{O}{c}{i}} the twisted label; \m{\IsApt{O}{c}{\vartheta}}: the clauses (A1)--(A6);
\m{\tupleProj{O}{c}} the action of~\m{\tupleIdem{r}{c}} on~\m{\mathcal{F}(O)}; \m{\HasRow{O}{z}{R}}: $R$ is a row and the normalised
representations of the eigenforms lifting~$z$ with character \m{\mu_O\vartheta} have the traces
\m{\invAvatar{\vartheta}{g}\tr R(g)+\sum_{e<\#n-\#r}\cyc(g)^{-e}}; \m{\IsLabelledColumn{O}{k}{\HT}{c}{S}{R}}: a column labelled
by the $c_i$ at the chosen places.
\item (3.25, 3.27--3.30) \m{\IsAuxiliary{S}{I}}; \m{\evenClasses{\cO}{S}{\nu}{e}}; \m{\IsNewPlaces{\nu}{w}{\tau}{\phi}};
\m{\withConjInv{\tau}{k}=\tau_j^{[b]}} for \m{k=(j,b)}; \m{\IsPrimeTriangle{\nu}{w}{\tau}{\phi}{a}{t}{\mathrm{ty}}{h}{c}}: the
clauses (T1)--(T5); \m{\IsAuxTwist{S}{Q}{\psi}{c}}: $Q$ is auxiliary for~$S$, $\psi$ is from $\mathrm{U}(1)$ relative to
\m{\pd{E}{\ell}} with modulus~$Q$, and \m{\sum_{\fq\in Q}\vl(\psi(\sigma_{\fq})-1)\le c} at tame generators.
\item (II: 3.3--3.6) The normaliser \m{\ribetNormaliser{\gamma}{a}{y}{p}{\theta}{\cdot}} is the function
\m{\delta_{\gamma a}-y\delta_a+p\theta}. The triangle cost \m{\HasTriangleCost{Z}{\tau}{\phi}{I}{C}{b}} consists of the
clauses (C1), (C2).
\m{\IsTraceIdeal{\mathcal{J}}{u}{p}{\tau}{\phi}{k}{\varepsilon}}: $\mathcal{S}$ contains the variables and the coefficients
of the characteristic polynomials, $\mathcal{J}$ the families \m{p\tr u_x(g)} and the tie relations of traces;
\m{\IsRibetElement{\mathcal{J}}{\beta}{\pi}{c}}: \m{\beta\in\mathcal{S}} vanishes on~$\mathcal{X}$, and some
\m{f\in\mathcal{S}} with \m{f\beta\in\mathcal{J}} has on~$\mathcal{X}$ no monomial properly dividing \m{\prod X_j} and the
coefficient $c$ of \m{\prod X_j} at~$\pi$.
\end{enumerate}
\end{definition}
\leanline{\lean{obeyingCocycles}\sep \lean{RowBlockIndex}\sep \lean{blockRep}\sep \lean{crysRels}\sep \lean{blockLoss}\sep
\lean{cmSign}\sep \lean{IsRowType}\sep \lean{lambdaLoss}\sep \lean{loss}\sep \lean{AbsGal}\sep \lean{placesAboveG}\sep
\lean{badPlaces}\sep \lean{invAvatar}\sep \lean{crudeUnit}\sep \lean{IsConjSelfDual}\sep \lean{SelmerLengthLe}\sep
\lean{IsColumn}\sep \lean{IsRow}\sep \lean{chiOneInv}\sep \lean{normalisedRep}\sep \lean{twistedLabel}\sep \lean{IsApt}\sep
\lean{tupleProj}\sep \lean{HasRow}\sep \lean{IsLabelledColumn}\sep \lean{IsAuxiliary}\sep \lean{evenClasses}\sep
\lean{IsNewPlaces}\sep \lean{withConjInv}\sep \lean{IsPrimeTriangle}\sep \lean{IsAuxTwist}\sep \lean{ribetNormaliser}\sep
\lean{HasTriangleCost}\sep \lean{IsTraceIdeal}\sep \lean{IsRibetElement}\sep defined in the first two notes of the unit}

\sect{The definitions of this note}{s:defs}

\begin{definition}[tests and denominators]\label{d:den}
Let $E$ be a CM field as in the conventions, $r$, $n$ finite sets, $O$ an auxiliary pair of type \m{(r,n)}
over~$E$, $e_0$ an element of the complex group algebra of \m{\prod_i\GL_r(\F_{w_i})}, and $Z$ a subspace
of~\m{\mathcal{F}(O)}.
\begin{enumerate}[label=\textup{(\alph*)},leftmargin=2.3em,itemsep=2pt]
\item \m{\groupAlgAction{O}{e_0}} is the endomorphism of~\m{\mathcal{F}(O)} by which $e_0$ acts.
\item A form \m{y\in\mathcal{F}(O)} is \emph{tested by \m{(O,e_0,Z)}} if for every finite place $w$ of~$E$ with
\m{\cj{w}\ne w} that lies outside the places of~$O$ and their conjugates, $y$ is orthogonal to
\[
\operatorname{im}\groupAlgAction{O}{e_0}\ \cap\ \ker H_w(X_w)\ \cap\ Z,\qquad X_w=\bigl(\nu_w\,\Nm(w)^{(\#n-\#r)-1}\bigr)^{-1},
\]
with \m{\nu_w} the Frobenius value of~$a_O$ in rank~\m{\#r} at~$w$ and the exponent truncated at~$0$.
\item For a prime number $l$ and \m{\Xi\in\mathcal{F}(O)}, \m{\IsTestForm{O}{l}{\Xi}} holds if $\Xi$ is integral at
every ring homomorphism \m{\CC\to\mathbb{C}_l} and is tested by \m{(O,\tupleIdem{r}{X-1},\ker\theta_O)}, with
\m{\tupleIdem{r}{X-1}} the idempotent of the tuple of labels all equal to \m{X-1}.
\item For an isomorphism of fields \m{\iota\colon\CC\to\Qlb} and a real number~$\beta$,
\m{\DenominatorLe{\iota}{O}{e_0}{Z}{\beta}} holds if, with \m{d=\dim Z}, for all families \m{(y_a)_{a<d}} and
\m{(y'_b)_{b<d}} of forms in~\m{\mathcal{F}(O)} such that every $y_a$ is integral at every ring homomorphism
\m{\CC\to\Cl}, and every $y'_b$ is integral at~$\iota$ and tested by \m{(O,e_0,Z^{\perp})}:
\[
\bigl|\iota\bigl(\det(\langle y_a,\mathrm{pr}_Z\,y'_b\rangle)_{a,b}\bigr)\bigr|\ \le\ \exp\bigl(d\cdot\#r\cdot\beta\bigr).
\]
\end{enumerate}
\end{definition}
\leanline{\lean{groupAlgAction}\sep \lean{IsTested}\sep \lean{IsTestForm}\sep \lean{DenominatorLe}}

\begin{definition}[the Hecke algebra away from $g$ and its eigenparts]\label{d:hecke}
Let $D$ be admissible data.
\begin{enumerate}[label=\textup{(\alph*)},leftmargin=2.3em,itemsep=2pt]
\item For a finite place $w$ of~\m{\ED{D}}, \m{\IsSplitOffLevel{D}{w}} holds if \m{\cj{w}\ne w} and neither $w$ nor
\m{\cj{w}} is a chosen place above~\m{\gD{D}}.
\item The \emph{Hecke algebra away from~$g$}, written \m{\mathbb{T}(D)}, is the algebra of the families
\m{(T_V)_V} of endomorphisms of the spaces \m{\mathcal{M}_D(V)}, indexed by the hermitian spaces $V$ of rank
\m{\Dr{D}} with a parameter~$s$, that are given, on all $V$ at once, by one element of the group algebra of the
group of finite adelic points trivial at the chosen places.
\item For a pair \m{x=(w,i)} of a finite place and a natural number, \m{\heckeOperator{D}{x}\in\mathbb{T}(D)} acts
on every \m{\mathcal{M}_D(V)} by the Hecke operator \m{T_{w,i}} if \m{\IsSplitOffLevel{D}{w}}, and is $0$ otherwise.
\item For a space $V$ and a function $p$ from the pairs \m{(w,i)} to~$\CC$, the \emph{$p$-part}
\m{\eigenPart{D}{V}{p}} is the space of the \m{y\in\mathcal{M}_D(V)} on which each \m{\heckeOperator{D}{x}} acts by the
scalar~\m{p(x)}.
\item For $V$, $p$ and a chosen place $x$ above~\m{\gD{D}}, \m{\occurringLabels{D}{V}{p}{x}} is the set of the
primitive labels $c$ of degree~\m{\Dr{D}} over the residue field of~$x$ whose idempotent, acting at~$x$, does not
kill \m{\eigenPart{D}{V}{p}}.
\item For $V$ and \m{T\in\mathbb{T}(D)}, \m{\heckeOnGood{D}{V}{T}} is the restriction of~$T_V$ to the good forms
on~$V$ (it preserves them by Lemma~\ref{l:commute}: $T_V$ commutes with the action of~\m{e_0(D)}); \m{T\mapsto\heckeOnGood{D}{V}{T}} is multiplicative.
\end{enumerate}
\end{definition}
\leanline{\lean{IsSplitOffLevel}\sep \lean{heckeAlgebra}\sep \lean{heckeOperator}\sep \lean{eigenPart}\sep
\lean{occurringLabels}\sep \lean{heckeOnGood}}

\begin{definition}[the adjoint lifts]\label{d:lifts}
Let $D$ be admissible data, $\vartheta$ a tame character of~$D$ and $j$ a natural number. By
Lemma~\ref{l:thetaprincipal} there is a linear map \m{L_{\vartheta,j}} from the forms of the auxiliary pair
\m{O_D(\vartheta,j)} to \m{\mathcal{M}_D(W_j(D))} which, for \m{1\le j\le\Dm{D}}, is the theta lift
\m{\theta_{O_D(\vartheta,j)}}; the formal text chooses one. For \m{F\in\mathcal{M}_D(W_j(D))},
\m{\auxLift{D}{\vartheta}{j}{F}} is the image of~$F$ under the adjoint of~\m{L_{\vartheta,j}}, a form
of~\m{O_D(\vartheta,j)}; and \m{\goodLift{D}{\vartheta}{j}} is the restriction of \m{\auxLift{D}{\vartheta}{j}{\cdot}} to the good
forms on~\m{W_j(D)}. For a natural number~$z$, \m{\RankOneIndex{z}} is the index set with
\m{12+4+4+(2z-20)} elements used for the group \m{\mathrm{U}(2z)} (\m{2z} elements when \m{z\ge10}).
\end{definition}
\leanline{\lean{auxLift}\sep \lean{goodLift}\sep \lean{RankOneIndex}}

\begin{definition}[ramified pairs, pairs of rank one, the budget for congruences]\label{d:pairs}
Let $D$ be admissible data and $p$ a pair of~$D$. Then \m{\IsRamified{p}} (\emph{$p$ is ramified}) holds if, for
every finite place~$w$, \m{\psi_{p,w}\ne1} exactly when $w$ is chosen. For a natural number~$j$,
\m{\IsRankOne{p}{j}} (\emph{$p$ is of rank one for~$j$}) holds if $p$ has lift rank \m{\RankOneIndex{z_p}}
(Chapter~VII), \m{\Dm{D}<z_p\le\Dr{D}}, \m{1\le j\le\Dm{D}}, \m{\psi_p=\omega_D^{\,j}}, and for every infinite
place~$v$: \m{c_p(v)\le\wt{D}{j}{v}} and \m{\wt{D}{j}{v}+\Dr{D}\le c_p(v)+\tD{D}}. The \emph{budget for congruences} is
\m{\Bzero{D}=\operatorname{budget}(D)/(\Dm{D}\,\kD{D})}.
\end{definition}
\leanline{\lean{IsRamified}\sep \lean{IsRankOne}\sep \lean{B}}

The predicate ``fit for~$D$'' of Theorem~A is the definition of the formal text: a family
\m{\theta=(\theta_{\ell,\iota})} of $m$-fold twists of~$D$, indexed by the pairs \m{(\ell,\iota)}, is fit if (F1), (F2)
and (F3) hold; the three clauses are its three fields, and the box of Theorem~A is its statement.

\begin{definition}[the budget of a prime]\label{d:rowbudget}
Let $D$ be admissible data and $\ell$ a natural number. Let \m{X_D\ne0} be the integer that Input~4 provides for
the conductor \m{\ND{D}d}, the number \m{\tD{D}} and the field \m{\ED{D}} (the formal text chooses one). Then
\[
\rowBudget{D}{\ell}=\hhD{D}^{27}\bigl(\crudeUnit{\ell}{D}+v_\ell(X_D)\bigr)+4\,\fD{D}\,\loss{D}{\ell}.
\]
\end{definition}
\leanline{\lean{rowBudget}}

\begin{definition}[reciprocity laws]\label{d:law}
Let $E$ be a CM field as in the conventions and \m{r'\in\NN}.
\begin{enumerate}[label=\textup{(\alph*)},leftmargin=2.3em,itemsep=2pt]
\item A Frobenius datum \m{x=(w_x,\fP_x,\phi_x)} of~$E$ is \emph{split for~$r'$} if $w_x$ does not contain~$\ell$
and \m{\cyc_{\ell^{r'}}(\phi_x)=1}.
\item For a finite place $w$ of~$E$, \m{\primeBelowPlus{w}} is the place of~$E^{+}$ below~$w$.
\item Let \m{e\in\NN}, $\Phi$ a subset of~\m{\GE{E}}, $T$ a finite set of finite places of~$E^{+}$ and
\m{z\colon\GE{E}\to\mathbb{Z}/\ell^{r'}} a function. Then $z$ is a \emph{law for} \m{(r',e,\Phi,T)} if: $z$ is
locally constant; for all $g$ and all $a$, $b$ with \m{\cyc_{\ell^{r'}}(a)=\cyc_{\ell^{r'}}(b)=1},
\[
z(gag^{-1}b)=\cyc_{\ell^{r'}}(g)^{e+1}z(a)+z(b);
\]
\m{z(b^{\mathsf c})=-z(b)} for all such~$b$; and for every finite family
\m{(x_j)_{j\in J}} of Frobenius data that are split for~$r'$, with the places \m{\primeBelowPlus{w_{x_j}}} pairwise
distinct and outside~$T$, every index $j$ with \m{z(\phi_{x_i})=0} for all \m{i\ne j}, every coefficient
ring~$\cO$, every \m{c\in\evenClasses{\cO}{\{w_{x_i}\}_i}{r'}{e}} that vanishes on~$\Phi$ and every natural number
$n'$ divisible by \m{2z(\phi_{x_j})} in~\m{\mathbb{Z}/\ell^{r'}}: \m{n'\cdot c(\tau_{\fP_{x_j}})=0}.
\end{enumerate}
\end{definition}
\leanline{\lean{IsSplitPlace}\sep \lean{primeBelowPlus}\sep \lean{IsLaw}}

\begin{definition}[central operators on forms of $L$-value data]\label{d:central}
Let $D$ be admissible data, $n$ a finite set and $Y'$ $L$-value data over~\m{\ED{D}} for~$n$ (Chapter~VII), with
auxiliary pair~$O_{Y'}$. For $z$ in the finite adelic points of~$\mathrm{U}(1)$, \m{\centralRep{Y'}(z)} is the
translation of the forms of~$O_{Y'}$ by the central element~$z$, and
\m{\twistedCentralOp{Y'}{z}=b_{Y'}(\det z)\cdot\centralRep{Y'}(z)}, with \m{b_{Y'}} the second character of the
auxiliary pair.
\end{definition}
\leanline{\lean{centralRep}\sep \lean{twistedCentralOp}}

\sect{Forms, Hecke operators and the test}{s:forms}

\begin{lemma}[small facts]\label{l:heckefacts}
Let $O$ be an auxiliary pair of type \m{(r,n)} over a CM field $E$ as in the conventions, and $D$ admissible data.
\textup{(a)} For an automorphism $\sigma$ of~$\CC$, forms \m{\varphi,\varphi'}, a split place $w$ outside the places
of~$O$ and their conjugates and \m{X\in\CC}: \m{\langle H_w(X)\varphi,\varphi'\rangle_\sigma=\langle\varphi,H_{\cj w}(X^{*})\varphi'\rangle_\sigma},
with \m{X^{*}=\sigma^{-1}(\overline{\sigma(X)})}.
\textup{(b)} For \m{g\in\prod_i\GL_r(\F_{w_i})}, the scalar \m{\prod_i\psi_{O,w_i}(\det g_i)} is a root of unity. The
image of the action of a group algebra element is stable under every endomorphism that commutes with the
action of the group; and a central idempotent acts self-adjointly for every \m{\langle\cdot,\cdot\rangle_\sigma}.
\textup{(c)} For a space $V$ as in the conventions, a function~$p$ and \m{y\in\eigenPart{D}{V}{p}}, the translates
of~$y$ by the group algebra of~\m{\mathsf{G}(D)} lie in \m{\eigenPart{D}{V}{p}}. If \m{\mathcal{M}_D(V)} is
finite-dimensional, then for every subspace~$N'$ only finitely many $p$ have \m{N'\cap\eigenPart{D}{V}{p}\ne0}.
\textup{(d)} For a finite place $w$ of~$E$: \m{\primeBelowPlus{\cj{w}}=\primeBelowPlus{w}}.
\end{lemma}
\begin{proof}
(a) is the adjoint formula for the Hecke operators at a split place (Chapter~VII). (b): the characters
\m{\psi_{O,w}} have finite order; the second assertion was used in Lemma~5.8 of the first note; the third follows
from Lemma~\ref{l:inneraction}. (c) follows from Lemma~\ref{l:commute}, and from the independence of the
simultaneous eigenspaces. (d): an element of~$E^{+}$ equals its conjugate.
\end{proof}
\leanline{\lean{innerAtheckePolyOnForms}\sep \lean{isOfFinOrderprodpsidet}\sep \lean{rangegroupAlgActionmeminvtSubmodule}\sep
\lean{innerAtgroupAlgAction}\sep \lean{principalRepmemeigenPart}\sep \lean{finiteeigenPart}\sep \lean{primeBelowPlusstar}}

\begin{lemma}\label{l:inneraction}
Let $E$ be a CM field as in the conventions, $r$, $n$ finite sets and $O$ an auxiliary pair of type \m{(r,n)}
over~$E$. For every automorphism $\sigma$ of~$\CC$, every \m{g\in\prod_i\GL_r(\F_{w_i})} and all forms
\m{x,y\in\mathcal{F}(O)}: \m{\langle g\cdot x,g\cdot y\rangle_\sigma=\langle x,y\rangle_\sigma}.
\end{lemma}
\begin{proof}
The action of~$g$ is translation by a lift of~$g$ times the scalar \m{\prod_i\psi_{O,w_i}(\det g_i)}. Translation
preserves the pairing, and the scalar contributes its absolute value squared after~$\sigma$, which is~$1$ for a
root of unity.
\end{proof}
\leanline{\lean{innerAtformAction}}

\begin{lemma}\label{l:heckecomm}
Let $E$, $r$, $n$, $O$ be as in Lemma~\ref{l:inneraction}. Let $w$ and $u$ be finite places of~$E$ with
\m{\cj{w}\ne w} and \m{\cj{u}\ne u}, both outside the places of~$O$ and their conjugates, let \m{X,X'\in\CC} and
\m{y\in\mathcal{F}(O)}. Then \m{H_w(X)H_u(X')y=H_u(X')H_w(X)y}.
\end{lemma}
\begin{proof}
The Hecke operators \m{T_{w,i}} and \m{T_{u,i'}} commute on the vectors spherical at both places, and the two sides
are the same double sum.
\end{proof}
\leanline{\lean{heckePolyOnFormscomm}}

\begin{lemma}[projection to a tested space does not increase norms]\label{l:normproj}
Let $E$, $r$, $n$, $O$ be as in Lemma~\ref{l:inneraction}, with its Weil data. Let $w$ be a finite place with
\m{\cj{w}\ne w} outside the places of~$O$ and their conjugates, $e_0'$ an element of the group algebra and $Z'$ a
subspace of~\m{\mathcal{F}(O)}. Suppose that $e_0'$ is a central idempotent, that
\m{Z'\subseteq\operatorname{im}\groupAlgAction{O}{e_0'}\cap\ker H_w(X_w)}, and that $y$ is tested by
\m{(O,e_0',Z'^{\perp})}. Then \m{\|\mathrm{pr}_{Z'}y\|_\sigma^2\le\|y\|_\sigma^2} for every automorphism $\sigma$ of~$\CC$.
\end{lemma}
\begin{proof}
The test at~$w$ says that $y$ is orthogonal to
\m{\operatorname{im}\groupAlgAction{O}{e_0'}\cap\ker H_w(X_w)\cap Z'^{\perp}}. By Lemma~\ref{l:heckefacts}(a),(b) the orthogonal
complement of \m{\operatorname{im}e_0'\cap\ker H_w(X_w)} is \m{\ker e_0'+\operatorname{im}H_{\cj w}(X_w^{*})}; hence
\m{y=z+x} with \m{z\in Z'} and \m{x\in\ker e_0'+\operatorname{im}H_{\cj w}(X_w^{*})}. For every~$\sigma$, the form
\m{\langle\cdot,\cdot\rangle_\sigma} pairs $x$ to zero with~$Z'$ (again by the adjoint formulas, as $Z'$ lies in the image
of~$e_0'$ and is killed by \m{H_w(X_w)}; Lemma~\ref{l:heckecomm} is used to move the operators). So
\m{\mathrm{pr}_{Z'}y=z} and \m{\|y\|_\sigma^2=\|z\|_\sigma^2+\|x\|_\sigma^2\ge\|z\|_\sigma^2}, the norm \m{\|x\|_\sigma^2} being a sum of
non-negative terms over the classes.
\end{proof}
\leanline{\lean{normSqAtstarProjectionle}}

\begin{lemma}\label{l:spherical}
Let $D$ be admissible data, $V$ a space as in the conventions, $p$ a function on the pairs \m{(w,i)},
\m{y\in\eigenPart{D}{V}{p}} and $w$ a finite place with \m{\IsSplitOffLevel{D}{w}}. Then $y$ is spherical at~$w$, and
the Hecke polynomial at~$w$ acts on~$y$ through \m{X\mapsto\sum_i(-X)^i\,\Nm(w)^{i(i-1)/2}\,p(w,i)}.
\end{lemma}
\begin{proof}
$y$ has principal level at the chosen places only, so it is fixed by the maximal compact subgroup at~$w$, and
\m{T_{w,i}} acts on it by \m{p(w,i)} by the definition of the $p$-part.
\end{proof}
\leanline{\lean{isSphericalWithofmemeigenPart}}

\begin{lemma}\label{l:resfield}
Let $P$ be parameters, $D$ admissible data for~$P$ and $x$ a chosen place above~\m{\gD{D}}, with residue field
of~$q$ elements. Then \m{9\mid q-1}, \m{\Mom{P}{\Dmu{D}}^2\mid q-1} and \m{\Mom{P}{\Dmu{D}}+1<q}, where
\m{\Mom{P}{\Dmu{D}}=(N_0d)^{\Dmu{D}+6}} is the integer of Chapter~III.
\end{lemma}
\begin{proof}
\m{q=\Dg{i}{D}} for the prime below~$x$, and the three facts are properties of the primes \m{\Dg{i}{D}} of
admissible data (Chapter~III).
\end{proof}
\leanline{\lean{dvdcardresFieldsubone}}

\begin{lemma}[separating eigensystems]\label{l:heckeproj}
Let $D$ be admissible data, $q'$ a function on the pairs \m{(w,i)} and $F$ a finite set of such functions. Then
there is \m{b\in\mathbb{T}(D)} such that for every \m{d'\in F\cup\{q'\}}, every space $V$ as in the conventions and
every \m{y\in\eigenPart{D}{V}{d'}}: \m{b_V(y)=y} if \m{d'=q'}, and \m{b_V(y)=0} otherwise.
\end{lemma}
\begin{proof}
By induction on~$F$. For a further system \m{e'\ne q'} choose $x$ with \m{e'(x)\ne q'(x)} and replace $b$ by
\m{b-(e'(x)-q'(x))^{-1}\,b\,(\heckeOperator{D}{x}-q'(x))}: it acts as before on the parts where $b$ acts by $0$ or on the
\m{q'}-part, and kills the $e'$-part.
\end{proof}
\leanline{\lean{existsheckeprojeigenPart}}

\begin{lemma}\label{l:commute}
Let $D$ be admissible data, $V$ a space as in the conventions, \m{T\in\mathbb{T}(D)} and $x$ an element of the group
algebra of~\m{\mathsf{G}(D)}. Then $T_V$ commutes with the action of~$x$ on~\m{\mathcal{M}_D(V)}.
\end{lemma}
\begin{proof}
$T_V$ is given by translations by elements trivial at the chosen places, $x$ by translations by elements
supported at the chosen places; such translations commute.
\end{proof}
\leanline{\lean{commuteheckeAlgebraprincipalRep}}

\begin{lemma}\label{l:eigenparts}
Let $D$ be admissible data and $V$ a space as in the conventions such that \m{\mathcal{M}_D(V)} is
finite-dimensional. Then \m{\mathcal{M}_D(V)} is the sum of the spaces \m{\eigenPart{D}{V}{p}} over all functions~$p$.
\end{lemma}
\begin{proof}
Right translations are unitary for the Petersson product, and the commutation and adjoint formulas for the
spherical Hecke operators at the places $w$ with \m{\IsSplitOffLevel{D}{w}} make them a commuting family of normal
operators on the finite-dimensional space \m{\mathcal{M}_D(V)} (the adjoint of an operator at~$w$ is one
at~\m{\cj{w}}). Such a family has a basis of joint eigenvectors (an input of the unit on the Weil
representation), and a joint eigenvector lies in the $p$-part of its eigenvalues.
\end{proof}
\leanline{\lean{iSupeigenPart}}

\begin{lemma}[the theta lift and the action of the chosen group]\label{l:thetaaction}
Let $E$, $r$, $n$, $O$ be as in Lemma~\ref{l:inneraction}, with its Weil data. Let $k$ be an element of the
stabiliser of the standard lattice of the hermitian space of~$O$ and $s'$ an element of \m{\prod_i\GL_r(\F_{w_i})}
such that $k$ has the reduction~$s'$ at the chosen places. Then for every form \m{\Xi\in\mathcal{F}(O)} and every
adelic point~$g$: \m{(\theta_O\Xi)(g\,k)=\theta_O(s'\cdot\Xi)(g)}.
\end{lemma}
\begin{proof}
Write the lift as the pairing of~$\Xi$ with the theta kernel. The pairing is equivariant, and the Weil
representation of~$k$ acts on the kernel, in the variable of~$\Xi$, through the reduction~$s'$ and the character
that enters the action of the chosen group (an input of the unit on forms, Chapter~VII).
\end{proof}
\leanline{\lean{thetamuleqthetaformAction}}

\begin{lemma}\label{l:thetaprincipal}
Let $D$ be admissible data, $\vartheta$ a tame character of~$D$ and $j$ a natural number. Then there is a linear map
$L'$ from \m{\mathcal{F}(O_D(\vartheta,j))} to \m{\mathcal{M}_D(W_j(D))} such that, if \m{1\le j\le\Dm{D}}, then
\m{L'(\Xi)=\theta_{O_D(\vartheta,j)}\Xi} for every~$\Xi$.
\end{lemma}
\begin{proof}
If $j$ is not a lift index, take \m{L'=0}. Otherwise the theta lift of an algebraic form of the auxiliary pair
is a form of the right type on~\m{W_j(D)} (Chapter~VII), and it has principal level at the chosen places by
Lemma~\ref{l:thetaaction} with \m{s'=1}.
\end{proof}
\leanline{\lean{existsthetatoprincipalForms}}

\sect{Rows exist, and the twist for rows: proofs of Theorems E and F}{s:rows}

\begin{lemma}[one row serves all twists]\label{l:rowtwists}
Let $D$ be admissible data, $r$, $n$ finite sets, $O$ an auxiliary pair of type \m{(r,n)} over~\m{\ED{D}}, and $c$
a family of polynomials with \m{\IsApt{O}{c}{1}}, $1$ the trivial twist. Let $b$ be Hecke data (a function of a
finite place and a natural number), $\HT$ a type and \m{R\colon\AbsGal{D}\to\GL_r(\Qlb)} a homomorphism. Suppose that
\m{\IsLabelledColumn{O}{\#r}{\HT}{c}{\emptyset}{R}}; that for every \m{g\in\AbsGal{D}}, \m{\iota(b(w,1))} tends to
\m{\tr R(g)} along the Frobenius data \m{(w,\fP,\phi)} with $\phi$ near~$g$ (in the filter of most Frobenius data
near~$g$); and that \m{\IsRowType{D}{\#r}{\#n-\#r}{\HT^{\flat}}}, where \m{\HT^{\flat}(\sigma)} is \m{\HT(\sigma)} shifted
by \m{-\bigl((\#r-(\#n-\#r))/2-\cmSign{\sigma}\,e_O([\sigma])\bigr)}, with $e_O$ the weight of~\m{\mu_O}. Then there is
a homomorphism \m{T\colon\AbsGal{D}\to\GL_r(\Qlb)} with \m{\HasRow{O}{b}{T}}.
\end{lemma}
\begin{proof}
By (A3), \m{\#n} is even; put \m{a'=(\#r-(\#n-\#r))/2}. Let $T$ be the twist of~$R$ by the avatar of
\m{\mu_O^{-1}[0,a']}. By Lemma~9.4 of the first note it is a labelled column with exponent \m{\#n-\#r}, the
type~\m{\HT^{\flat}} and the labels of~$c$ scaled by the name of the tame part of~\m{\mu_O^{-1}}, that is the
twisted labels \m{\twistedLabel{O}{c}{\cdot}}; by Theorem~B of the first note and the hypothesis on \m{\HT^{\flat}} it is a
row. For the trace identity of \m{\HasRow{O}{b}{T}}, let $\vartheta$, $V$, $y$, $H$, $S$ be as there. Both sides are
continuous, so it suffices to compare limits along Frobenius data near a given~$g$. At a datum \m{(w,\fP,\phi)}
with $w$ split, prime to~$\ell$ and off the places, the trace of~\m{S(\phi)} is close to minus the second
coefficient of the Hecke polynomial of~$y$, which is that of the string \m{\prod_{i<\#n-\#r}(1-q^{i}\nu Y)} times the
polynomial of \m{b(w,\cdot)} at \m{q^{\#n-\#r}\nu Y/\nu'} (the definition of lifting, with $Y$ the variable; \m{q=\Nm(w)}, and \m{\nu,\nu'} the
Frobenius values of \m{a_V} in the two ranks). Multiplying by \m{\chiOneInv{V}(\phi)=\iota(\nu)^{-1}} (Lemma~5.4 of the
first note), the string contributes \m{\sum_{e<\#n-\#r}q^{e}=\sum_e\cyc(\phi)^{-e}}, and the other factor contributes
\m{\iota(b(w,1))} times the values at~$\phi$ of the avatars of \m{\vartheta^{-1}} and of \m{\mu_O^{-1}[0,a']}; this is a
computation with the values of the characters at the uniformiser. Passing to the limit gives
\m{\invAvatar{\vartheta}{g}\tr T(g)+\sum_e\cyc(g)^{-e}}.
\end{proof}
\leanline{\lean{existsrowforalltwists}}

\begin{proof}[Proof of Theorem~E]
Let $z$ be a non-zero good form on \m{W_{\Dm{D}}(D)} in the $p$-part. Its image under the passage to algebraic
forms (Chapter~VII) is a form of some weight~$\kappa$ on the standard definite space of rank \m{\Dr{D}}, of principal
level at the chosen places, and an eigenform whose Hecke polynomial at \m{w} with \m{\IsSplitOffLevel{D}{w}} is
that of \m{p(w,\cdot)} (Lemma~\ref{l:spherical}). As \m{e_0(D)z=z\ne0} and \m{e_0(D)} is the sum of the \m{\tupleIdem{r}{c}} over the
good tuple families, some good family $c$ has \m{\tupleIdem{r}{c}z\ne0}; by Lemma~6.1(a) of the first note,
\m{\IsApt{O_D(1,j)}{c}{1}}. Let $R$ be a representation attached to~$z$ (Input~3 of the first note). At each chosen
place the form $z$ has the parahoric type cut out by the corresponding factor of~\m{\tupleIdem{r}{c}}, hence the product label
of~$c_i$ with no further blocks; by Lemma~9.5 of the first note,
\m{\IsLabelledColumn{O_D(1,j)}{\Dr{D}}{\HT_{\kappa,\Dr{D}}}{c}{\emptyset}{R}}. The traces of~$R$ are the limits of the Hecke values
\m{\iota(p(w,1))} (attachment), and the shifted type satisfies the conditions of a row type: this is the comparison
of the Harish-Chandra parameter of~$\kappa$ with the weights of the data, an input of the unit on the rows
(Chapter~VII). Lemma~\ref{l:rowtwists} gives~$T$.
\end{proof}

\begin{proof}[Proof of Theorem~F]
By Lemma~5.5(c) of the first note there is a chosen place $w$ in \m{\placesAboveG{D}\setminus\pd{\ED{D}}{\ell}} with
\m{\Nm(w)=\Dg{i}{D}}. For each ramified pair $p$ and sign~$s$, Theorem~D of the first note gives a family
\m{(\nu^{p,s},\mathrm{val}^{p,s})} for \m{\badPlaces{D}{\ell}}; there are finitely many pairs. By Theorem~C of the second note
there are a coefficient ring~$\cO$, lattices $u_i$ of the~$R_i$ and \m{\gamma_i}, \m{\beta_i} such that $u_i$ is
balanced, in the sense of Definition~3.33 of the first note, for the prime~$\ell$, the set
\m{\Sigma=\pd{\ED{D}}{\ell}\cup\placesAboveG{D}}, the relations \m{\crysRels{u_i}}, the pair \m{(\gamma_i,\beta_i)} and
the precision~$M$, and with the normalised length of
\m{\beta_i+\sum_bv_{\cO}(\det(u_i(\tau_w^{[b]})-1))} at most \m{10\fD{D}\#r_i\,\crudeUnit{\ell}{D}}. Apply Theorem~A of the
second note with \m{S=\Sigma}, \m{S_0=\badPlaces{D}{\ell}}, \m{z=\tau_w}, the bounds \m{B_i=10\fD{D}\#r_i\crudeUnit{\ell}{D}},
\m{c=\crudeUnit{\ell}{D}} (by Theorem~D(1) of the first note and Lemma~5.7(a) there, \m{\vl(\nu_i(\tau_w)-1)\le\vl(\Nm(w)-1)\le\crudeUnit{\ell}{D}}),
\m{\eta=\crudeUnit{\ell}{D}>0} and \m{m=\Dm{D}}. It gives a finite set $Q$ and~$\psi$ such that for every power
\m{1\le d\le\Dm{D}} (the lift index~$j$ of the statement):
\m{\IsAuxTwist{\badPlaces{D}{\ell}}{Q}{\psi^{d}}{\crudeUnit{\ell}{D}}}, the values of the families at \m{(Q,\psi^{d})} have
valuations at most \m{\vl(\mathrm{val}(\emptyset,1))+9\crudeUnit{\ell}{D}}, and the Selmer modules of the \m{\psi^{d}u_i}
modulo~\m{\ell^{M}} have normalised length at most \m{(10\fD{D}\#r_i+1)\crudeUnit{\ell}{D}} over all large coefficient
rings. By Lemma~8.3 of the first note, \m{\psi=\invAvatar{\tau}{\cdot}} for a tame character $\tau$ with \m{Q(\tau)=Q}, and
then \m{\psi^{d}=\invAvatar{\tau^{d}}{\cdot}}. The character $\tau$ is an $m$-fold twist: if \m{\psi_{\tau,w'}^{\,d}=1} at
some \m{w'\in Q} and some \m{1\le d\le\Dm{D}}, then \m{\psi^{d}} would be trivial at the tame generator at~$w'$, so
that \m{\vl(\psi^{d}(\sigma_{w'})-1)=\infty}, against the bound \m{\crudeUnit{\ell}{D}} in the auxiliary twist for that
power~$d$. This uses conclusion~(2) of Theorem~A of the second note for every power \m{d\le\Dm{D}}, not only for
\m{d=1}; its proof there chooses the roots of unity so that all \m{\zeta_v^{d}-1} have small valuation. Now (1) holds by
construction; (3) follows from Theorem~D(2) of the first note with \m{a=9\crudeUnit{\ell}{D}}; and (2) is the Selmer bound,
read as \m{\SelmerLengthLe{\invAvatar{\vartheta}{\cdot}}{R_i}{Q(\vartheta)}{M}{\gamma_i}{y_i}} with $y_i$ the scalar by which
\m{u_i(\gamma_i)} acts: the Selmer module with the relations \m{\crysRels{u_i}\cup\{\delta_{\gamma_i}\}} carried to a
larger coefficient ring contains the cocycles vanishing at~$\gamma_i$ in the Selmer module with the relations
\m{\crysRels{\psi'u_i}}, because the crystalline cocycles are compatible with the extension of the coefficients and
\m{\psi'} is trivial on the decomposition groups above~$\ell$; \m{\psi'(\gamma_i)=1} because $\gamma_i$ is inertial
at a place of~$\Sigma$ prime to~$\ell$, where a character from $\mathrm{U}(1)$ with modulus~$Q$ is trivial.
\end{proof}

\sect{Ribet's element for columns: proof of Theorem C}{s:ribet}

\begin{lemma}[obeying cocycles are Selmer cocycles]\label{l:obeysel}
Let $D$ be admissible data, $\cO$ a coefficient ring, $Y$, $n$, $r$ finite sets, \m{u=(u_x)_{x\in Y}}
homomorphisms \m{u_x\colon\AbsGal{D}\to M_n(\cO)}, $S$ a set of finite places of~\m{\ED{D}},
\m{\varrho\colon\AbsGal{D}\to M_r(\cO)} a homomorphism, \m{b\in\cO} non-zero, $L$ a set of relations and
\m{c\colon\AbsGal{D}\to(\cO/b)^r} a function. Suppose that every $u_x$ is continuous, that \m{u_x(g)=1} for all
\m{g\in\inab{S^{\complement}}} and all~$x$, that \m{c\in\obeyingCocycles{u}{\varrho}{b}} and that \m{c\in\Ob{L}{\cO/b}}.
Then \m{c\in\Sel{\varrho}{S}{L}{\cO/b}}.
\end{lemma}
\begin{proof}
By Lemma~4.1 of the first note there are $k$ and a module $Y_0$, with injective first projection, containing all
\m{((u_x(g)-1)_x+b^kv,c(g))}. If \m{u_x(g)=1} for all~$x$, then \m{(0,c(g))\in Y_0}, so \m{c(g)=0}: $c$ is unramified
outside~$S$. If \m{u_x(h)\equiv u_x(g)} modulo~$b^k$ for all~$x$, the two points have the same first coordinate
for a suitable~$v$, so \m{c(h)=c(g)}; this condition is open because the $u_x$ are continuous: $c$ is locally
constant.
\end{proof}
\leanline{\lean{memselmerModule}}

\begin{lemma}[doubling the tied primes]\label{l:traceconj}
Let $D$ be admissible data, $\cO$ a coefficient ring, $K$, $Y$ sets, $n$ a finite set, \m{t,\nu\in\NN},
\m{w=(w_k)_{k\in K}} finite places of~\m{\ED{D}}, \m{\tau,\phi\colon K\to\AbsGal{D}}, \m{u=(u_x)_{x\in Y}}
homomorphisms \m{\AbsGal{D}\to\GL_n(\Qlb)}, $\mathcal{S}$ a subalgebra and $\mathcal{J}$ an $\mathcal{S}$-submodule of
the functions \m{Y\to\Qlb[X_k:k\in K]}, and \m{\varepsilon=(\varepsilon_k)_k} polynomials. Suppose
\m{\IsNewPlaces{\nu}{w}{\tau}{\phi}}, \m{\IsConjSelfDual{t}{u_x}} for all~$x$, and
\m{\IsTraceIdeal{\mathcal{J}}{u}{\ell^{\nu}}{\tau}{\phi}{\mathrm{id}}{\varepsilon}}. Then
\m{\IsTraceIdeal{\mathcal{J}}{u}{\ell^{\nu}}{\withConjInv{\tau}{\cdot}}{\withConjInv{\phi}{\cdot}}{\mathrm{pr}_1}{\varepsilon}}, where the
families are indexed by the pairs \m{(k,b)} of an index and a sign and \m{\mathrm{pr}_1(k,b)=k}.
\end{lemma}
\begin{proof}
Only the relations at the indices \m{(k,-)} are new. By conjugate self-duality, the trace of
\m{u_x((h^{[-]})^{i}g)} is \m{\cyc(\cdot)^{-(t-1)}} of that element times the trace of \m{u_x(h^{i}(g^{\mathsf c})^{-1})}.
By Lemma~7.2 of the first note, \m{\cyc(\tau_k^{[-]})=1} and \m{\cyc(\phi_k^{[-]})^{-e}\equiv1} modulo~\m{\ell^{\nu}}. So
the new relation at~$g$ is \m{\cyc(g)^{-(t-1)}} times the old relation at \m{(g^{\mathsf c})^{-1}}, up to $X_k$ times
a combination of the families \m{\ell^{\nu}\tr u_x(\cdot)}; all of these lie in~$\mathcal{J}$.
\end{proof}
\leanline{\lean{isTraceIdealwithConjInv}}

\begin{lemma}[crystalline losses above $\ell$ are affordable]\label{l:blockloss}
Let $D$ be admissible data, $Y$, $n$, $r$ finite sets, \m{t\in\NN}, $V$ an auxiliary pair of type \m{(r,n)}
over~\m{\ED{D}}, $\vartheta$ a twist and $c$ a family of polynomials with \m{\IsApt{V}{c}{\vartheta}}. Let $\HT$, $\HT'$ be
types, $S$ a set of finite places, \m{R\colon\AbsGal{D}\to\GL_r(\Qlb)} a homomorphism with \m{\IsRow{D}{t}{\HT}{R}}, and
\m{R'=(R'_x)_{x\in Y}} homomorphisms \m{\AbsGal{D}\to\GL_n(\Qlb)} with
\m{\IsLabelledColumn{V}{t}{\HT'}{(\twistedLabel{V}{c}{i}(X-1)^{t})_i}{S}{R'_x}} for all~$x$. Let $\epsilon$ be a bijection
from the disjoint union of the sets \m{\RowBlockIndex{r}{i}}, \m{i\in\{\ast\}\sqcup\{0,\dots,t-1\}}, to~$n$, and
\m{\pi\in Y}. Then for all large~$\cO$: for all homomorphisms \m{T,T_0\colon\AbsGal{D}\to M_r(\cO)}, every family
\m{u=(u_x)_{x\in Y}} of homomorphisms \m{\AbsGal{D}\to M_n(\cO)} and every family \m{(M_i)_i} of natural numbers such
that \m{u_\pi(g)} is the block-diagonal matrix of the \m{\blockRep{T}{i}{g}} transported by~$\epsilon$, every $u_x$ is
a lattice of~$R'_x$, $T_0$ is a lattice of~$R$, and $T$ and $T_0$ agree on \m{G_{\fP_v}} for every
\m{v\in\pd{\ED{D}}{\ell}}:
\[
\sum_{i}\nlen_{\cO}\bigl(\blockLoss{\blockRep{T}{i}{\cdot}}{u}{M_i}\bigr)\ \le\ 4\,\#r\cdot\fD{D}\cdot\loss{D}{\ell}.
\]
\end{lemma}
\begin{proof}
Let $\cO$ be large enough for Theorem~E of the first note at every place $v$ above~$\ell$, with the chart of~$v$.
A function obeys the crystalline relations of a block if and only if it obeys, for each~$v$, the relations of
the crystalline cocycles at~$v$ (Lemma~4.2(a) of the first note); so each block loss is at most the sum over~$v$
of the lengths of the obeying cocycles modulo those obeying the relations at~$v$. By Lemma~7.3 of the second
note, applied to the restrictions to \m{G_{\fP_v}} and to a class that contains the local lattices of all $u_x$
and of the block, this is at most the length of the crystalline cocycles of the twisted dual of the block
modulo the cup-orthogonal of the class. The class is that of Lemma~7.2 of the second note: by that lemma the
local lattices of the labelled columns $u_x$ belong to it, with the jumps of~$\HT'$, the tame type at the places
above~\m{\gD{D}} coming from Lemma~9.2 of the first note and Lemma~7.1 of the second (the eigenvalues of inertia
are roots of unity of bounded order); the block~$T$ agrees above~$\ell$ with a lattice of the row, and the string
blocks have the single jump~$e$. Theorem~E of the first note bounds each term by
\m{[\Gl:j_v(G_{\fP_v})]\cdot\rk\cdot\loss{D}{\ell}/x}, with \m{x=1} for the row block of rank~\m{\#r} and
\m{x=\tD{D}} for each of the \m{t\le\tD{D}} string blocks. Summing over the blocks gives at most
\m{2\#r\,\loss{D}{\ell}} times the index, and the sum of the indices over~$v$ is at most the degree
\m{[\ED{D}:\mathbb{Q}]=2\fD{D}} (Lemma~4.4 of the first note).
\end{proof}
\leanline{\lean{sumblockLossle}}

\begin{proof}[Proof of Theorem~C]
The formal proof has a hundred lines; this is its outline. \emph{Lattices.} Choose one coefficient ring $\cO'$
over which: each row $R_x$, \m{x\in\mathcal{X}}, has a lattice with the Selmer bound of
\m{\SelmerLengthLe{\psi}{R_x}{Q}{m_0}{\gamma_x}{y_x}}; each column $u_x$ has a lattice~$u'_x$; and for
\m{x\in\mathcal{X}} this lattice is block diagonal, a lattice $T_x$ of \m{\psi R_x} and the string
\m{\cO'(0),\dots,\cO'(-(t_V-1))}. The last point is the theorem of Brauer and Nesbitt applied to the trace
identity of the hypothesis: the row is absolutely irreducible because its lattices span all matrices (W3), and
the characters \m{\cyc^{-e}} are distinct. By Lemma~\ref{l:blockloss} the block losses of \m{T_\pi} and of the
string, at any moduli, total at most \m{4\#r\,\fD{D}\,\loss{D}{\ell}}.

\emph{The application of Ribet's lemma.} Apply Theorem~D of the second note over \m{A=\cO'}, with the points~$Y$,
the subset~$\mathcal{X}$, the block \m{\varrho_x=T_x}, the characters \m{\xi_e=\cyc^{-e}} (\m{\xi_0=1}),
\m{p=\ell^{\nu}}, the doubled families \m{\withConjInv{\tau}{\cdot}}, \m{\withConjInv{\phi}{\cdot}} with the trace ideal of
Lemma~\ref{l:traceconj}, the element $a$ of the triangle with \m{C=\ell^{\hhD{D}h}} (clause (T3)),
\m{p_0=\ell^{m_0}}, the elements \m{\gamma_x} and scalars \m{y_x} of the Selmer bounds (\m{y_x-1} a unit,
\m{\cyc(\gamma_x)=1}), and, as elements \m{\delta_v} with polynomials $o_v$, two kinds of inertial elements, each
with both signs. The first are the tame generators at the places of~$Q$: they are killed by
\m{(X-\psi)(X-1)}, and \m{o_v=X-\psi(\delta_v)}; the values \m{o_v(1)} have total valuation at most
\m{2\crudeUnit{\ell}{D}} by the auxiliary twist. The second are the generators that Lemma~9.2 of the first note
provides for the inertia above~\m{\gD{D}} off~$\ell$. Their images under every~$u_x$ are killed by
\m{X^{\ordch{D}}-1}: the eigenvalues are roots of unity of orders dividing some \m{\Dg{i}{D}^{d}-1} with
\m{d\le\Dr{D}\tD{D}} (Lemma~9.2 of the first note), and these numbers divide \m{\ordch{D}} (Lemma~5.5(a) there). So
\[
o_v=1+X+\dots+X^{\ordch{D}-1}\qquad\text{satisfies}\qquad(o_v\cdot(X-1))(u_x(\delta_v))=0,
\]
and \m{o_v(1)=\ordch{D}} has valuation at most \m{\crudeUnit{\ell}{D}}. (The killing polynomial of Lemma~9.2 divided
by \m{X-1} would not do: for a labelled column it contains the root~$1$ with the multiplicity~$t_V$ of the
string, so the quotient vanishes at~$1$ as soon as \m{t_V\ge2}.) Hypothesis~(e) holds with $b_0$ the Selmer length
bound plus the loss of the row block: by Lemma~\ref{l:obeysel} the obeying cocycles modulo~\m{\ell^{m_0}} that obey
the crystalline relations and vanish at the tame generators are Selmer cocycles for \m{T_\pi} unramified outside
\m{\pd{\ED{D}}{\ell}\cup\placesAboveG{D}\cup Q\cup\cj{Q}} (Lemma~7.3 of the first note), and
\m{m_0>\Dr{D}\,\rowBudget{D}{\ell}} exceeds that bound. Hypothesis~(g) holds by Theorem~E of the second note for the
prime triangle, with $b_e$ the loss of the $e$-th string block plus the length of
\m{\Sha^{\mathrm{odd}}(\ell,\ED{D},\cO',e,\nu)}, which is at most the valuation of the integer~$X_D$ (Input~4), and with
the finite set
\[
s'=s\cup\{\tau_k^{[b]}:k\in K,\ b=\pm\},\qquad s\text{ the set of clause (T2)}:
\]
a cocycle that vanishes at the \m{\tau_k^{[b]}} is unramified at the~$w_k$ and their conjugates (Lemma~7.3 of the
first note), and then (T2) detects it on~$s$.

\emph{The bound.} Theorem~D of the second note gives $c$ with \m{\IsRibetElement{\mathcal{J}}{\beta}{\pi}{c}} and
\[
\vl(c)\le b_0+\sum_eb_e+t_V\Bigl(\sum_v\vl(o_v(1))+2\#s'\cdot\hhD{D}h\Bigr)+(\#K+1)\hhD{D}h,
\]
where \m{\#s'\le\#s+2\#K}.
Each term is at most a multiple of \m{\crudeUnit{\ell}{D}}, of \m{v_\ell(X_D)} or of \m{\loss{D}{\ell}}. The Selmer
bound is \m{(10\fD{D}\#r+1)\crudeUnit{\ell}{D}}. The losses total \m{4\#r\fD{D}\loss{D}{\ell}}. The number $h$ is at
most \m{v_\ell(\ordch{D})}, hence at most \m{\crudeUnit{\ell}{D}}, and \m{(\#K+\#s)\hhD{D}h\le\hhD{D}^{22}h} by (T2). The sum of the \m{\vl(o_v(1))} is at most
\m{2\crudeUnit{\ell}{D}} for the places of~$Q$ plus \m{\crudeUnit{\ell}{D}} for each of the at most \m{2\cdot4\fD{D}}
generators above~\m{\gD{D}} (Lemma~5.5(b) of the first note). Finally \m{t_V\le\tD{D}} and \m{\fD{D}} are bounded by
powers of \m{\hhD{D}} (Chapter~III). With \m{3v_\ell(\ordch{D})} for the factor
\m{\ordch{D}^3}, the total is at most \m{\#r\,(\hhD{D}^{27}(\crudeUnit{\ell}{D}+v_\ell(X_D))+4\fD{D}\loss{D}{\ell})}. The comparison
of the constants was not re-derived here.
\end{proof}

\inp{Input 4 (a bound for the odd part of Tate--Shafarevich groups of abelian CM fields; background library, unit of the cyclotomic fields)}{Let
\m{N\ge1} and $t$ be natural numbers ($t$ may be~$0$) and $E$ a CM field whose fixed embedding lands in the
cyclotomic field of conductor~$N$. Then there is a natural number \m{X\ne0} with \m{X\le2^{4(N+t+3)^{27}}\cdot h_E}, $h_E$
the class number, such that for every prime~$l$, every discrete valuation ring $A$ of characteristic~$0$ that
is finite and free over~\m{\mathbb{Z}_l}, has finite quotients and contains the roots of unity of order the exponent
of \m{(\mathbb{Z}/N)^{\times}}, every \m{e<t} and every~$n'$: \m{\len_A\Sha^{\mathrm{odd}}(l,E,A,e,n')\le v_A(X)}. For
\m{e=0} the module is cut out by vanishing on the inertia above~$l$, with no oddness condition; for \m{e>0} by
oddness under conjugation and vanishing on the kernel of the cyclotomic character in the decomposition groups
above~$l$. \emph{References} (second reading): C.~Greither, \emph{Class groups of abelian fields, and the main
conjecture}, Theorem~3.2, p.~469; B.~Mazur, A.~Wiles, \emph{Class fields of abelian extensions of~$\mathbb{Q}$},
pp.~179--184. \emph{Relation:} an explicit finite-level estimate. Greither's theorem is an identity of
characteristic ideals in an Iwasawa algebra; the single integer, uniform in the prime, the twist and the
precision, and its explicit size need further estimates, which the library carries out (it treats \m{e=0}
separately and uses Stickelberger bounds and finite detecting sets). No printed statement with this uniform
bound has been identified.}

\sect{Reciprocity laws and prime triangles: proof of Theorem D}{s:laws}

\begin{lemma}\label{l:evenzero}
Let $E$ be a CM field as in the conventions, $\cO$ a coefficient ring, $S$ a set of finite places, \m{r',e\in\NN},
\m{c\in\evenClasses{\cO}{S}{r'}{e}} and \m{b\in\cO/\ell^{r'}}. Suppose that \m{u'(b)=0} for every additive map
\m{u'\colon\cO/\ell^{r'}\to\mathbb{Z}/\ell^{r'}} such that \m{u'\circ c} is a continuous cocycle of~\m{\GE{E}} for the
character \m{\cyc_{\ell^{r'}}^{-e}}. Then \m{b=0}.
\end{lemma}
\begin{proof}
The module \m{\cO/\ell^{r'}} is killed by~\m{\ell^{r'}}, so the additive maps to \m{\mathbb{Z}/\ell^{r'}} separate its
points; and for every such map \m{u'\circ c} is a continuous cocycle for \m{\cyc_{\ell^{r'}}^{-e}}, because $c$ is a
locally constant cocycle for~\m{\cO(-e)} and $u'$ is \m{\mathbb{Z}}-linear.
\end{proof}
\leanline{\lean{eqzeroofmemevenClasses}}

\begin{lemma}[a reciprocity law with prescribed zeros]\label{l:law}
Let $E$ be a CM field as in the conventions, \m{r',e,k} natural numbers, $I$ a finite set, \m{(x_i)_{i\in I}}
Frobenius data of~$E$ and \m{(y_i)_{i<k}} elements of~\m{\GE{E}}. Suppose that \m{0<r'}, that $E$ has at least four
infinite places, that \m{k\le\#I}, that every $x_i$ is split for~$r'$, and that the places
\m{\primeBelowPlus{w_{x_i}}} are pairwise distinct. Then there are a function \m{z\colon\GE{E}\to\mathbb{Z}/\ell^{r'}} and
a finite set $T$ of finite places of~$E^{+}$ such that $z$ is a law for \m{(r',e,\{\phi_{x_i}\}_i,T)},
\m{z(y_i)=0} for all \m{i<k}, and \m{z(g)} is a unit for some $g$ with \m{\cyc_{\ell^{r'}}(g)=1}.
\end{lemma}
\begin{proof}
Outline of a formal proof of 140 lines. Pass to~$E^{+}$, a totally real field: \m{\GE{E}} has index~$2$ in
\m{\GE{E^{+}}} and contains no complex conjugation; choose a section of the restriction and let
\m{M'=\mathbb{Z}/\ell^{r'}} with \m{\GE{E^{+}}} acting through the corresponding lift of \m{\cyc^{-e}}, which is
trivial on complex conjugations. Apply the relations at split primes (Input~1 of the second note, with
\m{K=E^{+}}, \m{m=\ell^{r'}} and \m{X=M'}; its Selmer groups carry the condition of vanishing on the complex
conjugations, which is not empty here). The set $S$ consists of the places of~$E^{+}$ above~$\ell$ and the places
\m{\primeBelowPlus{w_{x_i}}}; the local condition at a place above~$\ell$ is the vanishing on the kernel of the
character, and at the places \m{\primeBelowPlus{w_{x_i}}} it is the strict one: vanishing on the whole
decomposition group (not ``no condition''). The subgroup $U'$ is an open normal subgroup inside the kernel of the
character and of \m{\cyc_{\ell^{r'}}}. This gives a group of relation families \m{(n_i)} among the values of Selmer
cocycles at the elements~\m{\gamma_i}, each of which extends, by Input~1(b), to an equivariant family on~$U'$
that satisfies the reciprocity relations at further split places outside the finite set~$T'$. A count of the $\ell$-torsion of
the Selmer group (the same input for~\m{M'[\ell]} with two further primes from Chebotarev's theorem, a procyclic
generator of the character, the hypotheses \m{k\le\#I} and on the four infinite places) shows that the relation
families have rank greater than~$k$; so there is a primitive one, with a unit entry, that vanishes at the
$y_i$. Let \m{z(g)} be its component at the class of the restriction of~\m{g^{-1}}. Local constancy, the rule for
\m{gag^{-1}b} and the oddness come from the equivariance of the extended family; the reciprocity clause is the
relation of the family at new split places, applied to the cocycle obtained from an even class through
Lemma~\ref{l:evenzero} (the factor~$2$ comes from the passage from~$E$ to~$E^{+}$).
\end{proof}
\leanline{\lean{existsisLaw}}

\begin{proof}[Proof of Theorem~D]
Outline of a formal proof of a hundred lines. By Input~5 there is a finite set \m{\mathfrak{p}\subseteq\GE{E}} with at
most \m{y^{20}} elements that detects the cocycles of all \m{\cO(-e)} unramified outside~$\ell$. Let $J$ be the set
of the pairs \m{(e,i)} with \m{0<e<t} and \m{i\in\mathfrak{p}}, in lexicographic order, and \m{\mathrm{ty}} the first
projection. Fix \m{\nu>0}. By recursion along~$J$ one chooses, for \m{j=(e,i_0)}, a Frobenius datum
\m{(w_j,\fP_j,\phi_j)\in W}, a law $z_j$ of twist~$e$ (Lemma~\ref{l:law}) and its exceptional set, as follows. The
law $z_j$ is chosen to vanish at points that realise a basis of the values of the earlier laws of the same
twist, and to be a unit at some point. An element \m{u'} with \m{\cyc_{\ell^{\nu}}(u')=1} is found at which every
earlier law vanishes (the laws of another twist~\m{e'} by the factor \m{m^{e}-m^{e'}}, with $m$ the class of
\m{\cyc(\gamma)}), while \m{z_j(u')} is a unit times \m{\prod_{e'\ne e}(m^{e}-m^{e'})}. By the hypothesis on~$W$ there
is a Frobenius datum in~$W$ whose Frobenius element is close to
\m{\delta^{\ell^{\nu}}\cdot g_{i_0}g_{i_0}^{\mathsf c}\cdot u'(u'^{\mathsf c})^{-1}}, with \m{g_{i}=i^{-1}\gamma\,i\,\gamma^{-1}}, outside the earlier
primes and exceptional sets; there \m{z_j(\phi_j)=2z_j(u')}, the earlier laws vanish, and every even class $c$
of twist~$e$ satisfies \m{2c(g_{i_0})=c(\phi_j)}. The places are Taylor--Wiles places for~\m{\ell^{\nu}} because the
Frobenius elements fix the \m{\ell^{\nu}}-th roots of unity. The clauses: (T1) by construction; (T2) with
\m{s=\mathfrak{p}}, the cost \m{(\#J+\#\mathfrak{p})yh\le y^{22}h} following from \m{\#\mathfrak{p}\le y^{20}} and \m{t<y};
(T3) from the hypothesis that \m{\cyc(\gamma)^{a}-1} divides~\m{\ell^{h}}, with room for the factor \m{\ell^{3h}}; (T4)
from the reciprocity clauses of the laws, the divisor \m{2z_j(\phi_j)} dividing \m{\ell^{yh}}; (T5) from the detection
by~$\mathfrak{p}$ and Sah's lemma: an even class that vanishes at the $\phi_j$ of type~$e$ vanishes, up to the
factor~$2$, at the \m{g_i}, \m{i\in\mathfrak{p}}, which gives the coboundary identity
\m{\delta(a)\cdot2z(g)=\delta(g)\cdot2z(a)} with \m{a=\gamma^{-1}}.
\end{proof}

\inp{Input 5 (finitely many points detect the cocycles unramified outside $\ell$; background library, unit of the cyclotomic fields)}{Let
\m{N\ge1}, $E$ a number field whose fixed embedding lands in the cyclotomic field of conductor~$N$, and $l$ a
prime. Then there is a finite set \m{\mathfrak{p}\subseteq\GE{E}}, with at most \m{y^{20}} elements for every $y$ with
\m{N\le y} and \m{l\le y^2}, such that for every discrete valuation ring $A$ of characteristic~$0$ that is finite
and free over~\m{\mathbb{Z}_l} with finite quotients, and all \m{n',e\in\NN}, a cocycle in \m{\Sel{A(-e)}{\pd{E}{l}}{\emptyset}{A/l^{n'}}} that vanishes on~$\mathfrak{p}$ is zero.
\emph{Reference:} none; the catalogue of this write-up records that the library proves it by generators of ray
class groups and a count through Minkowski's bound, and that no published source is given.}

\sect{The denominator bound: proof of Theorem B}{s:den}

\begin{lemma}\label{l:cycdvd}
Let $P$ be parameters and $D$ admissible data for~$P$. Then there are a natural number
\m{h\le v_\ell(\ordch{D})}, with \m{\ell\le\hhD{D}^2} or \m{h=0}, and \m{\gamma\in\AbsGal{D}} such that \m{\cyc(\gamma)^{a}-1}
divides \m{\ell^{h}} in~\m{\Zl} for every \m{0<a<\tD{D}}.
\end{lemma}
\begin{proof}
If \m{\hhD{D}^2<\ell}, take \m{h=0} and for $\gamma$ an element of an inertia group above~$\ell$ whose cyclotomic
character reduces to a generator of \m{(\mathbb{Z}/\ell)^{\times}}; it fixes \m{\ED{D}}, which lies in a cyclotomic field
of conductor less than~$\ell$ and prime to it, and \m{\cyc(\gamma)^{a}-1} is a unit because \m{a<\tD{D}<\ell-1}.
Otherwise take \m{h=v_\ell(\ordch{D})} and for $\gamma$ a Frobenius element at a chosen place of norm \m{\Dg{i}{D}}
prime to~$\ell$ (Lemma~5.5(c) of the first note): then \m{\cyc(\gamma)^{a}-1=\Dg{i}{D}^{a}-1} divides \m{\ordch{D}}
(Lemma~5.5(a) there), which divides \m{\ell^{h}} in~\m{\Zl}.
\end{proof}
\leanline{\lean{existscycCharpowsubonedvd}}

\begin{proof}[Proof of Theorem~B]
Outline of a formal proof of 120 lines. Let $\mathcal{A}$ be the set of the $y$ of the hypothesis; its span
is~$Z$. Let \m{d=\dim Z} and \m{(y_a)}, \m{(y'_b)} be as in Definition~\ref{d:den}(d). By Lemma~4.5 of the first
note, applied with \m{V=\mathcal{F}(O)}, the map $j$ to smooth forms, the linear forms
\m{\eta_a=\langle y_a,\cdot\rangle} and \m{m=\ordch{D}^{3}}, it suffices to find, for each non-zero
\m{p_1\in\mathcal{A}} with its family $c$ of labels, a linear map $G$ on~\m{\mathcal{F}(O)} that acts by a scalar
(times~$j$) on every element of~$\mathcal{A}$, whose scalar $\lambda$ at~$p_1$ satisfies
\m{\vl(\ordch{D}^{3}\lambda)\le\#r\,\rowBudget{D}{\ell}}, and such that \m{\ordch{D}^{3}\langle y_a,w'\rangle} is
$\ell$-integral whenever \m{j(w')=G(\mathrm{pr}_Zy'_b)}. The exponent of the conclusion is then
\m{d\cdot\#r\cdot\rowBudget{D}{\ell}\log\ell}.

\emph{The base level and the eigensystems.} Theorem~C of the first note, for $O$, $c$ and~$\vartheta$, gives a base
level~$B$ containing the image of \m{\tupleProj{O}{c}}, such that \m{\ordch{D}\cdot\tupleProj{O}{c}} preserves
$\ell$-integrality, and whose eigensystems $x$ at levels off \m{\badPlaces{D}{\ell}} have normalised representations
\m{u_x} that are labelled columns with the quadratic relation on the inertia above \m{Q(\vartheta)\cup\cj{Q(\vartheta)}}. The
elements of~$\mathcal{A}$ with the label~$c$ are eigenvectors (they lift Hecke data); their eigensystems form the
set~$\mathcal{X}$, and by hypothesis~(b) each has a row $R_x$ with the trace identity of
\m{\HasRow{V}{z}{R_x}} and the Selmer bound. An eigensystem on which every Galois element has the eigenvalue~$1$
without being in~$\mathcal{X}$ is a \emph{weak lift}.

\emph{The primes.} Let $h$ and $\gamma$ be as in Lemma~\ref{l:cycdvd}. Theorem~D is applied to \m{E=\ED{D}} in the
cyclotomic field of conductor \m{\ND{D}d}, with \m{t=\#n-\#r} and \m{y=\hhD{D}}, and to the set $W$ of the Frobenius
data \m{(w,\fP,\phi)} with three properties. First, the twisted row \m{\psi R_{p_1}}, with
\m{\psi=\invAvatar{\vartheta}{\cdot}} and \m{R_{p_1}} the row of~$p_1$, satisfies
\m{|\det(1-\psi(\phi)R_{p_1}(\phi))|\ge1}: no eigenvalue of \m{\psi(\phi)R_{p_1}(\phi)} is congruent to~$1$ modulo the
maximal ideal. Secondly, $w$ lies off \m{\badPlaces{D}{\ell}}, off its conjugates and off the level of the base
level. Thirdly, the weak lifts that are congruent to~$p_1$ (their characteristic polynomials are congruent to
those of~$p_1$ coefficient by coefficient) have at~$w$ their Hecke string (Lemma~10.4 of the first note). The
hypothesis of Theorem~D on~$W$ holds: by Lemma~10.2 of the first note, applied to \m{T=\psi R_{p_1}} and to the
powers \m{\delta=\gamma_0^{\,i}} of the scalar inertial element $\gamma_0$ of~(W4), on which $\psi$ is trivial and
\m{\cyc=1}, the first property holds for most Frobenius data near every \m{\delta^{\ell^{\nu}}g}; the other two hold
for most Frobenius data. Theorem~D gives for every~$\nu$ Taylor--Wiles places with
\m{\IsPrimeTriangle{\nu}{w}{\tau}{\phi}{\gamma^{-1}}{t}{\mathrm{ty}}{\hhD{D}h}{\hhD{D}^{22}h}}.

\emph{Ribet's element and the operator.} Let $\mathcal{S}$ be the algebra generated by the Hecke operators of the
base level at the level of the Taylor--Wiles places, with the tie parameters as variables, acting on the
eigensystems, and $\mathcal{J}$ the module of the families that kill the forms of that level; the trace ideal
condition holds by the theory of the note on patching (unit of the ties at auxiliary primes). Theorem~C, with a
normaliser~$\theta'$ chosen so that Ribet's determinant~$\beta$ does not vanish on the eigensystems that are
neither lifts nor weak lifts, and with $\nu$ large with respect to that determinant, gives \m{f\in\mathcal{S}} with
\m{f\beta\in\mathcal{J}} and the coefficient~$d'$. An input of the note on patching turns $f$ into an operator \m{E_s}
on the forms of the base level, diagonal on the eigensystems, with eigenvalue the coefficient~$d'$ on the lift
\m{p_1}, such that \m{\ordch{D}^{2}E_s} preserves $\ell$-integrality on the forms orthogonal to the weak lifts that
are not lifts. Put \m{G=E_s\circ\tupleProj{O}{c}}.

\emph{The test.} It remains to see that \m{\tupleProj{O}{c}(y'_b-\mathrm{pr}_Zy'_b)} is killed by~\m{E_s}, that is, has
no component on the lifts of~$\mathcal{X}$ and on the weak lifts congruent to them. This is where the test of
Definition~\ref{d:den}(b) is used: a weak lift is killed by the Hecke polynomial \m{H_w(X_w)} at a suitable split
place~$w$ (its Hecke polynomial has the string of roots of Lemma~10.3 of the first note, and $X_w$ is the
inverse of the last root), it lies in the image of~$e_0$, and \m{y'_b} is orthogonal to such forms inside
\m{Z^{\perp}}. In detail, Lemma~6.4 of the first note reduces the claim to the isotypic vectors of the space
\[
\operatorname{im}\tupleProj{O}{c}\ \cap\ Z^{\perp}\ \cap\ \bigl(\operatorname{im}\groupAlgAction{O}{e_0}\cap\ker H_w(X_w)\bigr)^{\perp},
\]
which contains \m{\tupleProj{O}{c}(y'_b-\mathrm{pr}_Zy'_b)} by the test and is stable under the Hecke operators
(Lemmas 5.7(c) and~5.8 of the first note, Lemma~\ref{l:heckecomm}). An isotypic vector of this space that \m{E_s}
does not kill is a weak lift; hence it is killed by \m{H_w(X_w)} (Lemma~10.3 of the first note) and lies in the
image of~$e_0$, so it is orthogonal to itself, and is zero. (Lemma~\ref{l:normproj} is not used here.) The orders of the stabilisers divide \m{\ordch{D}}, which accounts for the third factor.
\end{proof}

\sect{The fit family: proof of Theorem A}{s:fit}

\begin{proof}[Proof of Theorem~A]
\emph{Step 1: rows.} For every \m{(\ell,\iota)} and every pair $p$ of rank one for an index~$j$, Lemma~\ref{l:rowtwists}
gives a row \m{R'_{p,j}} for the auxiliary pair \m{O_p(1)} and the Hecke data of the trivial lift: it is applied to
the trivial character of rank one, a labelled column with labels \m{X-1} (Lemma~6.1(b) of the first note) and
the type \m{\{0\}}, whose shift is a row type by the inequalities of a pair of rank one. For every lift index $j$
and every function $p'$ whose part meets the good forms on \m{W_{\Dm{D}}(D)} non-trivially, Theorem~E gives a row
\m{R_{j,p'}} for \m{O_D(1,j)}; by Lemma~\ref{l:heckefacts}(c) there are finitely many such~$p'$, and there are
finitely many pairs. Note that the rows \m{R_{j,p'}} are built from eigensystems found on the one space
\m{W_{\Dm{D}}(D)}, for every~$j$; Step~4 meets eigensystems on \m{W_j(D)}, and Input~6 below connects the two.

\emph{Step 2: the twists.} For each pair \m{(\ell,\iota)}, Theorem~F, applied to the finitely many rows of Step~1
and to \m{M=\lfloor\Dr{D}\,\rowBudget{D}{\ell}\rfloor+1}, gives an $m$-fold twist \m{\theta_{\ell,\iota}}. This is the
family.

\emph{Step 3: two budgets.} For every finite set $T$ of prime numbers,
\[
\sum_{l\in T}\rowBudget{D}{l}\log l\le\Bzero{D}\qquad\text{and}\qquad\sum_{l\in T}9\,\crudeUnit{l}{D}\log l\le x(D)^{999}.
\]
Indeed \m{\sum_{l\in T}v_l(a)\log l\le\log a} for an integer \m{a\ge1}, which bounds the contributions of
\m{\ordch{D}}, of~$X_D$ (by Input~4 and the bound of the class number by the discriminant) and of
\m{\gD{D}} and the factorials in \m{\loss{D}{l}}; \m{\sum_l2^{-l}\log l} is bounded; and the resulting sums of
logarithms are at most \m{\Bzero{D}/2} each, resp.\ \m{x(D)^{999}}, by the inequalities of the data (Chapter~III).

\emph{Step 4: the clauses.} (F1): by Theorem~F(1) for \m{j=1}, \m{Q(\theta_{\ell,\iota})} is auxiliary for
\m{\badPlaces{D}{\ell}}, so $\ell$ is prime to its norms (Lemma~8.2 of the first note). (F2): take
\m{\eta_l=\rowBudget{D}{l}\log l} for all~$j$; the first budget of Step~3 is the condition on the sums. Fix~$j$
and \m{(\ell,\iota)}, and let \m{\vartheta=\theta_{\ell,\iota}^{\,j}}; by Theorem~F(1) it satisfies the hypothesis of
Theorem~B.

\emph{The good lift.} The image of \m{\goodLift{D}{\vartheta}{j}} is spanned by the images of the eigenvectors $F$ in
the parts \m{\eigenPart{D}{W_j(D)}{p'}} of the good forms (Lemma~\ref{l:eigenparts}), cut by the idempotents
\m{\tupleIdem{r}{c}} of the good families (inputs of Chapter~VII). Consider such an image that is not zero. It lies
in \m{\operatorname{im}\tupleProj{O_D(\vartheta,j)}{c}}, with \m{\IsApt{O_D(\vartheta,j)}{c}{\vartheta}} (Lemma~6.1(a) of the
first note) and \m{\operatorname{im}\tupleIdem{r}{c}\subseteq\operatorname{im}e_0(D)}; this is hypothesis~(a) of Theorem~B.
It lifts the Hecke data~$p'$ (the adjoint lift of an eigenform, Chapter~VII). The eigensystem~$p'$ occurs in the
good forms on \m{W_j(D)}, since $F$ is a non-zero good form in its part. By Input~6, applied with the lift
indices $j$ and~\m{\Dm{D}}, it occurs in the good forms on \m{W_{\Dm{D}}(D)} as well; so the row \m{R_{j,p'}} of
Step~1 is defined. It is a row for \m{V=O_D(1,j)}, whose character differs from that of \m{O_D(\vartheta,j)}
by~$\vartheta$, and the Selmer bound is Theorem~F(2), the row being among those of Step~1; this is hypothesis~(b).
So Theorem~B gives
\[
\DenominatorLe{\iota}{O_D(\vartheta,j)}{e_0(D)}{\operatorname{im}\goodLift{D}{\vartheta}{j}}{\eta_\ell}.
\]

\emph{The lines of rank one.} For a pair $p$ of rank one for~$j$, the line \m{\CC\,\Xi_p(\vartheta)} is spanned by one
vector, which lies in the image of the idempotent with labels \m{X-1} (Lemma~6.1(c) of the first note, with (b)
for aptness), lifts the Hecke data of the trivial lift, and has the row \m{R'_{p,j}}; Theorem~B applies in the
same way, and no transfer is needed.

(F3): by Theorem~F(3), for every \m{(\ell,\iota)},
\m{|\iota(\varrho_p(\theta_{\ell,\iota}^{\,j},s)^{-1})|\le\exp(9\crudeUnit{\ell}{D}\log\ell)}; with \m{H=1}, \m{\alpha=0} and
the local bounds \m{9\crudeUnit{l}{D}\log l} for the primes~$l$, the second budget of Step~3 gives the height
bound (note \emph{Heights}).
\end{proof}

\inp{Input 6 (eigensystems of good forms on the spaces of two single lifts; unit \emph{Good idempotents, labels and the transfer of traces}, Chapter~VII: note \emph{Good forms and rows}, Theorems 2(b) and~3)}{Let
$P$ be parameters, $D$ admissible data for~$P$, $j$ and $j'$ lift indices \m{(1\le j,j'\le\Dm{D})}, and $p$ a
function from the pairs of a finite place of~\m{\ED{D}} and a natural number to~$\CC$.
\textup{(a)} \emph{Rigid labels.} For every such~$p$ and every chosen place $x$ above~\m{\gD{D}}, the set
\m{\occurringLabels{D}{W_j(D)}{p}{x}} has at most one element.
\textup{(b)} \emph{Transfer.} If the space of the good forms on \m{W_{j'}(D)} meets
\m{\eigenPart{D}{W_{j'}(D)}{p}} only in~$0$, then the space of the good forms on \m{W_j(D)} meets
\m{\eigenPart{D}{W_j(D)}{p}} only in~$0$. Equivalently: an eigensystem that occurs in the good forms on \m{W_j(D)}
occurs in the good forms on \m{W_{j'}(D)}.
Part~(b) is proved in that note from part~(a) for both spaces and from the equality of the traces of the Hecke
operators on the good forms of two unitary groups that differ at the real places; that equality rests on a
comparison of trace formulas, an input of the background library there. The spaces \m{W_j(D)} and \m{W_{j'}(D)}
have different signatures, so (b) is not a formal consequence of the definitions.}

\sect{A bound for rank one}{s:rankone}

\begin{lemma}[the central operators]\label{l:central}
Let $D$ be admissible data, $n$ a finite set, $Y'$ $L$-value data over~\m{\ED{D}} for~$n$ and $z$ a finite adelic
point of~$\mathrm{U}(1)$.
\textup{(a)} For forms \m{x,y} of~$O_{Y'}$: \m{\langle x,\twistedCentralOp{Y'}{z}y\rangle=\langle\twistedCentralOp{Y'}{z^{-1}}x,y\rangle}.
\textup{(b)} For every adelic point~$g$, \m{\twistedCentralOp{Y'}{z^{-1}}} carries the kernel form of~$Y'$ at~$g$ to
the kernel form at~\m{gz}.
\textup{(c)} \m{\twistedCentralOp{Y'}{z}\,\Xi_{Y'}=\Xi_{Y'}}, where \m{\Xi_{Y'}} is the lift of the trivial character.
\end{lemma}
\begin{proof}
(a) The scalar \m{b_{Y'}(\det z)} has absolute value~$1$, and translation by a central element preserves the
pairing. (b) Compare the values at each point: the Weil representation of the scalar pair attached to
\m{\det z^{-1}} is the value of the character times the translation (an input of Chapter~VII). (c) The lift of
the trivial character is the class sum of the kernel forms, also of the translated ones, which (b) maps back.
\end{proof}
\leanline{\lean{innertwistedCentralOp}\sep \lean{twistedCentralOpkernelForm}\sep \lean{twistedCentralOptrivialLift}}

\begin{lemma}[the denominator of a rank-one line, in terms of two test forms]\label{l:rankone}
Let $D$ be admissible data, $n$ a finite set, $Y'$ $L$-value data over~\m{\ED{D}} for~$n$, \m{(\ell,\iota)} a pair of a
prime and an isomorphism \m{\CC\to\Qlb}, $\beta$ a real number and \m{\Xi,\Xi'} forms of~$O_{Y'}$. Suppose that
\m{\DenominatorLe{\iota}{O_{Y'}}{\tupleIdem{r}{X-1}}{\CC\cdot\Xi_{Y'}}{\beta}} ($r$ the index set of the rows of~$O_{Y'}$), that the stabilisers in the rational points
of~$\mathrm{U}(1)$ of all classes at the level of the chosen places of~$Y'$ are trivial, and that
\m{\IsTestForm{O_{Y'}}{\ell}{\Xi}} and \m{\IsTestForm{O_{Y'}}{\ell}{\Xi'}}. Then, with $h'$ the number of classes
of~$\mathrm{U}(1)$ at that level,
\[
\bigl|\iota\bigl(h'\cdot\langle\Xi,\mathrm{pr}_{\CC\Xi_{Y'}}\Xi'\rangle\bigr)\bigr|\le e^{\beta}.
\]
\end{lemma}
\begin{proof}
If \m{\Xi_{Y'}=0} the left side is~$0$. Otherwise the line has dimension~$1$, and the type of~$O_{Y'}$ has
\m{\#r=1}. Choose representatives \m{s_1,\dots,s_{h'}} of the classes that are integral above~$\ell$. Apply the
definition of the denominator bound with \m{d=1}, \m{y_0=\Xi} and \m{y'_0=\sum_c\twistedCentralOp{Y'}{s_c^{-1}}\Xi'}. By
Lemma~\ref{l:central}(a),(c) each term contributes \m{\langle\Xi,\mathrm{pr}\,\Xi'\rangle}, so the determinant is
\m{h'\langle\Xi,\mathrm{pr}\,\Xi'\rangle}. The form \m{y'_0} is integral at~$\iota$ because the central operators at
integral points preserve integrality; it is tested by \m{(O_{Y'},\tupleIdem{r}{X-1},(\CC\Xi_{Y'})^{\perp})} because the central
operators commute with the idempotent and the Hecke polynomials, and because the sum over the classes of the
central translates of a form orthogonal to~\m{\Xi_{Y'}} lies in \m{\ker\theta_{O_{Y'}}} (the theta lift to
$\mathrm{U}(1)$ pairs a form with the kernel forms, Lemma~\ref{l:central}(b); the stabilisers are trivial), so
that the test of~$\Xi'$ against \m{\ker\theta_{O_{Y'}}} applies.
\end{proof}
\leanline{\lean{rankonebound}}

\sect{Where the results are used}{s:use}
Theorem~A is used once, in the proof of Theorem~C of the handout (integral forms are long): the use stands in
the file of the endgame on long forms, which the overview's map places in Chapter~II. There the
fit family provides, for each \m{(\ell,\iota)}, the twist at which the volume of the chain is compared with the
volume at the trivial twist; (F1) makes the index of the rows available at~\m{\theta_{\ell,\iota}}, (F2) bounds the
denominators of the Gram determinants of the single lifts and of the rank-one lines, and (F3) the heights of the
ratios of $L$-values. Lemma~\ref{l:rankone} is not used by Theorem~A; it is used in the unit of the rows
(Chapter~VII), in the height of the Hecke values, where the clause on the rank-one lines of (F2) enters through
it. Lemma~\ref{l:normproj} is not used by Theorem~A or by Theorem~B either; it is used in the unit of the rows
(Chapter~VII), in the height of the volume of a lift. The definitions of the Hecke algebra away from~$g$ and of the adjoint lifts (Definitions \ref{d:hecke},
\ref{d:lifts}) are also used by the notes on the volume of the chain.

\sect{Notes on the formal text}{s:notes}
\begin{enumerate}[label=\textup{\arabic*.},leftmargin=1.6em,itemsep=5pt]
\item \emph{Letters.} Besides the changes of the first two notes: the formal text writes {\leanfont E} for the
element~$e_0$ of the group algebra in Theorem~B, {\leanfont O} for an auxiliary pair, {\leanfont Y} for
$L$-value data (here~$Y'$), {\leanfont r} for the exponent~$r'$ of Definition~\ref{d:law}, {\leanfont p} for a
function of eigenvalues as well as for a pair, {\leanfont s} for the scalar~$s'$ of Theorem~B, {\leanfont θ}
both for the family of twists and for the function~$\theta'$ of Theorem~C, {\leanfont d} for the coefficient~$d'$,
and {\leanfont B₀} for \m{\Bzero{D}}.
\item \emph{Arguments.} The arguments were read from the text of the files; none was confirmed by Lean's own
answer, since the artifact was not compiled here. In several declarations the program lists as arguments
variables that occur only as names of fields or as bound names: a place~$w$ in the definitions of split
Frobenius data and of laws and in the ramified and rank-one pairs; the field and bound variables in the
definitions of the Hecke algebra and its relatives; and, in the definition of a fit family, the prime, the
isomorphism, the group algebra element, the auxiliary pair and the index, which are bound inside its clauses.
Lean presumably takes for a fit family only the data and the family~$\theta$. These definitions are therefore
written out in words here (Definitions \ref{d:den}(b), \ref{d:hecke}(b), \ref{d:law} and the box of Theorem~A), with
every listed variable introduced where it is bound. Conversely the program misses the arguments ``a prime, with
the isomorphism~$\iota$'' in Lemma~\ref{l:rowtwists} and Theorem~E (their statements use~$\iota$; the calls in the
last file pass it), and ``the data'' in Lemma~\ref{l:resfield}, whose statement begins with a local definition.
In the definitions of the maps \m{\auxLift{D}{\vartheta}{j}{\cdot}} and \m{\heckeOnGood{D}{V}{\cdot}} the argument of the
map is counted as an argument. All this is to be confirmed with Lean's own answer when the artifact is built.
\item \emph{Theorem~A.} The formal statement is ``for all data there is~$\theta$ that is a fit family''; the
index set is the set of the pairs of a prime and an isomorphism. The clause (F2) of the formal definition
quantifies the lift index first, then the function~$\gamma$, then the pair \m{(\ell,\iota)}, then the pair of rank
one; (F3) quantifies the lift index, the pair and the sign before the hypothesis ``ramified''. The reading of
these clauses, of the height and of the Euler factors of the ratio follows the second reader's message on the
fit family; I compared it with the formal definition and found no difference.
\item \emph{Hecke operators.} The Hecke operators away from the chosen places form a commuting family of
\emph{normal} operators (Lemma~\ref{l:eigenparts}); the dossier calls them unitary, which would be a further
condition (second reader).
\item \emph{Input~4.} The catalogue row for the bound of the odd part asked \m{t\ge1} and called the module
``the odd part'' without restriction; the second reader's verdict, followed here, corrects both: $t$ is any
natural number, and for \m{e=0} there is no oddness condition.
\item \emph{Imports copied by name.} For 211 of the declarations of other units that this part of the dossier
quotes, the reading in words was copied from other dossiers of the workspace; where such a declaration is
stated here as an Input, its statement is that of the first two notes or of the catalogue.
\item \emph{Outlines.} The proofs of Theorems B, C, D and of Lemma~\ref{l:law} are outlines of long formal proofs
(of 120, 100, 100 and 140 lines);
I compared the hypotheses used at each stage with the formal text and with the reader's step-by-step rendering,
but did not re-derive the estimates. In particular the comparison of constants at the end of Theorem~C, the
recursion of Theorem~D, the count of torsion in Lemma~\ref{l:law} and the construction of the operator in
Theorem~B (which uses the note on patching) are taken from the formal text.
\item \emph{What the final theorems use.} Lemma~\ref{l:rankone}, Lemma~\ref{l:normproj} and
Definition~\ref{d:central} are not used by Theorem~A; they are used in the unit of the rows (Chapter~VII).
\item \emph{The transfer between the spaces of the lifts.} The formal proof of Theorem~A builds the rows of the
good eigensystems from the space \m{W_{\Dm{D}}(D)} and uses them for eigensystems found on \m{W_j(D)}; it passes
from one to the other by two results of the unit on good forms (Input~6), in the contrapositive form of
part~(b). They are imports of this unit and do not appear in the concordance, which lists the declarations of
the unit only.
\item \emph{Theorem~C.} At the inertia above~\m{\gD{D}} the polynomials $o_v$ are the sums
\m{1+X+\dots+X^{\ordch{D}-1}}, and the finite set of Ribet's lemma is the set of clause~(T2) enlarged by the tame
generators \m{\tau_k^{[\pm]}}. The formal proof ends with a comparison of the resulting bound with
\m{\#r\cdot\rowBudget{D}{\ell}}, in which the first term carries the coefficient
\m{4+\tD{D}\,(2+10\fD{D}+6\hhD{D}^{22})+\hhD{D}^{22}+\hhD{D}}; I did not expand that comparison.
\end{enumerate}

\subhead{How this note was written.} From the dossier of the unit (its third chunk): the statements in words and
the structured proofs, which a program compared line by line with the source (\UnitSteps{} leaf steps in the
\UnitProofs{} proofs of the whole unit), and from the second reader's message on the fit family, which takes
precedence. I read in the Lean text: the file of the definitions of this note completely; the last file (the
proof of Theorem~A) completely; the files on the twist for rows, on the existence of rows and on rank one
completely; and the statements, with the surrounding definitions, of the four long files (Ribet's element,
reciprocity laws, prime triangles, the denominator bound), whose proofs I followed in the reader's rendering.
In particular the proofs of the three lemmas taken from the file of Ribet's element (Lemmas \ref{l:obeysel},
\ref{l:traceconj}, \ref{l:blockloss}) were written from the reader's rendering, not from the Lean text; the
referee of this chapter compared them with the Lean source. In the revision of this note I read in the Lean text
the passages of the four long proofs that the revision touches: the choice of the polynomials and of the finite
set in Theorem~C, the set of Frobenius data and the test in Theorem~B, the local conditions in
Lemma~\ref{l:law}, and the call of the transfer in Theorem~A.
The statements of this note were compared with the Lean statements argument by argument, and the order of the
quantifiers is that of the formal text. I did not open the proofs of the inputs. The references of Inputs 4
and~5 and their places are those of the second reading of this write-up and of its catalogue; I did not check
them in the literature.

\head{Appendix A. Concordance with the Lean artifact}
{\fontsize{8.6}{11.6}\selectfont\renewcommand{\leanfont}{\ttfamily\fontsize{7.7}{10}\selectfont\color{leancol}}
\setlength{\LTleft}{0pt}\setlength{\LTright}{\fill}\setlength{\tabcolsep}{0pt}
\begin{longtable}{@{}p{27mm}p{61mm}p{41mm}r@{}}
\textbf{In this note} & \textbf{Lean declaration} & \textbf{File} & \textbf{Lines}\\\midrule\endhead
\ConcordanceRows
\end{longtable}}

\subhead{The definition of a fit family and its existence, as written in the artifact.}
\quoted{IsFitFamily}
\quoted{existsisFitFamily}
In these texts {\leanfont PrimePadicIso} is the set of the pairs \m{(\ell,\iota)}, {\leanfont ι.1.1} the prime and
{\leanfont ι.2} the isomorphism; {\leanfont D.MFoldTwist} the set of the $m$-fold twists, {\leanfont levelNorm}
the product \m{\Nm(Q(\cdot))}, {\leanfont powTwist j} the $j$-th power; {\leanfont D.B₀} is \m{\Bzero{D}}, {\leanfont D.e₀}
the good idempotent, {\leanfont goodLift} the adjoint lift on good forms, {\leanfont lValueData} and {\leanfont
trivialLift} give \m{O_p(\vartheta)} and \m{\Xi_p(\vartheta)}, {\leanfont tupleIdem \_ fun \_ ↦ .X - 1} is \m{\tupleIdem{r_p}{X-1}}, {\leanfont
HeightLe} the height bound and {\leanfont ratio} is \m{\varrho_p}. The formal definition of the denominator bound
(Definition~\ref{d:den}(d)) is not reproduced: it contains a sign, for the orthogonal complement, that the fonts
of this note lack; its place is in the table above.

\head{Notation}
\begingroup\fontsize{8.6}{10.9}\selectfont\raggedright
\setlength{\columnsep}{9mm}\setlength{\parskip}{0pt}
\begin{multicols}{2}
\nt{$\ell$, $l$, $\iota$, $(\ell,\iota)$, $\Ql$, $\Zl$, $\Qlb$, $\Zlb$, $\Cl$, $\mathbb{C}_l$, $\mathbb{Z}_l$, $\mathbb{Q}_l$, $\Gl$}{primes, the isomorphism \m{\CC\to\Qlb}, the pair that indexes the fit family, the $\ell$-adic objects}
\nt{$\Qb$, $\mathsf{c}$, $g^{\mathsf c}$, $g^{[b]}$, $E$, $E^{+}$, $\GE{E}$, $h_E$}{as in the first note; the maximal real subfield; the class number}
\nt{$P$, $D$, $\ED{D}$, $\Dm{D}$, $\Dr{D}$, $\kD{D}$, $\tD{D}$, $\fD{D}$, $\Dg{i}{D}$, $\gD{D}$, $\lvl{D}$, $\hhD{D}$, $\bone{D}$, $\ND{D}$, $d$, $\wt{D}{j}{v}$, $\ordch{D}$, $x(D)$, $\operatorname{budget}(D)$, $\Dmu{D}$, $N_0$, \m{\Mom{P}{\Dmu{D}}}, $\omega_D$}{parameters, data and its numbers}
\nt{$\mathsf{G}(D)$, $e_0(D)$, $W_j(D)$, $\mathcal{M}_D(V)$, $\mathbb{T}(D)$}{the chosen group, the good idempotent, the space of the $j$-th lift, forms of principal level, the Hecke algebra away from~$g$}
\nt{$w$, $v$, $u$, $\fq$, $\cj{w}$, $\Nm(w)$, $\F_w$, $\fP$, $\fP_w$, $G_{\fP}$, $I_{\fP}$, $\tau_{\fP}$, $\pd{E}{M}$, $\inab{S}$, $S^{\complement}$}{places and the attached groups ($u$ also a family of representations)}
\nt{$\cyc$, $\cyc_M$, $\vl$, $v_\ell$, $v_l$, $v_A$, $\len$, $\nlen$, $\cO$, $\cO'$, $\cO(-e)$, $A$}{characters, valuations, lengths, coefficient rings}
\nt{$\Sel{\varrho}{S}{L}{B}$, $\Ob{L}{B}$, $\Sha^{\mathrm{odd}}$, $\delta_g$}{Selmer modules, obeying functions, the module of Input~4, indicator relations}
\nt{$\vartheta$, $\tau$, $\theta=(\theta_{\ell,\iota})$, $Q(\vartheta)$, $\Nm(Q(\vartheta))$, $\psi_{\vartheta,w}$}{twists, $m$-fold twists and the fit family; supports and their norms ($\tau$ also a family of tame generators)}
\nt{$p$, $\nu_p$, $z_p$, $c_p$, $\psi_p$, $\varrho_p(\vartheta,s)$, $L^{(S)}(z,\nu)$, $\mathrm{Ch}$}{pairs; ratios of $L$-values; $L$-functions without the Euler factors at a set; the chosen places}
\nt{$O$, $V$, $O_D(\vartheta,j)$, $O_p(\vartheta)$, $O_{Y'}$, $Y'$, $\mathcal{F}(O)$, $w_i$, $\mu_O$, $\psi_O$, $a_O$, $e_O$, $b_{Y'}$, $\theta_O$}{auxiliary pairs, $L$-value data, forms, characters, the weight, the theta lift}
\nt{$\Xi$, $\Xi'$, $\Xi_p(\vartheta)$, $\Xi_{Y'}$, $F$}{forms; the lifts of the trivial character}
\nt{$\langle y,y'\rangle$, $\langle y,y'\rangle_\sigma$, $\|y\|_\sigma$, $\mathrm{pr}_Z$, $Z^{\perp}$, $\operatorname{im}$, $\ker$, $\dim$}{Petersson products, norms, orthogonal projection and complement}
\nt{\m{\tupleIdem{r}{c}}, \m{\tupleIdem{r}{X-1}}, $r_p$, $e_0$, $e_0'$, $H_w(X)$, $X_w$, $X^{*}$, $\nu_w$, $T_{w,i}$, $T_V$}{idempotents; Hecke operators; the test value}
\nt{\m{\groupAlgAction{O}{e_0}}, \m{\IsTestForm{O}{l}{\Xi}}, \m{\DenominatorLe{\iota}{O}{e_0}{Z}{\beta}}}{Definition~\ref{d:den}}
\nt{\m{\IsSplitOffLevel{D}{w}}, \m{\heckeOperator{D}{x}}, \m{\eigenPart{D}{V}{p}}, \m{\occurringLabels{D}{V}{p}{x}}, \m{\heckeOnGood{D}{V}{T}}}{Definition~\ref{d:hecke}}
\nt{\m{\auxLift{D}{\vartheta}{j}{F}}, \m{\goodLift{D}{\vartheta}{j}}, \m{\RankOneIndex{z}}, $L_{\vartheta,j}$}{Definition~\ref{d:lifts}}
\nt{\m{\IsRamified{p}}, \m{\IsRankOne{p}{j}}, \m{\Bzero{D}}, \m{\rowBudget{D}{\ell}}, $X_D$}{Definitions \ref{d:pairs}, \ref{d:rowbudget}}
\nt{\m{\primeBelowPlus{w}}, \m{\centralRep{Y'}}, \m{\twistedCentralOp{Y'}{z}}}{Definitions \ref{d:law}, \ref{d:central}}
\nt{\m{\obeyingCocycles{u}{\varrho}{a}}, \m{\RowBlockIndex{r}{i}}, \m{\blockRep{T}{i}{g}}, \m{\crysRels{u}}, \m{\blockLoss{u}{w}{M}}}{recalled, Definition~\ref{d:recall}(1)}
\nt{\m{\cmSign{\sigma}}, \m{\IsRowType{D}{r}{t}{\HT}}, \m{\lambdaLoss{h}{b}{l}}, \m{\loss{D}{l}}}{recalled, Definition~\ref{d:recall}(2)}
\nt{\m{\AbsGal{D}}, \m{\placesAboveG{D}}, \m{\badPlaces{D}{\ell}}, \m{\invAvatar{\vartheta}{g}}, \m{\crudeUnit{\ell}{D}}}{recalled, Definition~\ref{d:recall}(3)}
\nt{\m{\IsConjSelfDual{t}{R}}, \m{\SelmerLengthLe{\psi}{R}{Q}{m}{\gamma}{z}}, \m{\IsColumn{D}{t}{\HT}{S}{R}}, \m{\IsRow{D}{t}{\HT}{R}}}{recalled, Definition~\ref{d:recall}(4)}
\nt{\m{\chiOneInv{O}}, \m{\normalisedRep{O}{R}}, \m{\twistedLabel{O}{c}{i}}, \m{\IsApt{O}{c}{\vartheta}}, \m{\tupleProj{O}{c}}, \m{\HasRow{O}{z}{R}}, \m{\IsLabelledColumn{O}{k}{\HT}{c}{S}{R}}}{recalled, Definition~\ref{d:recall}(5)}
\nt{\m{\IsAuxiliary{S}{I}}, \m{\evenClasses{\cO}{S}{\nu}{e}}, \m{\IsNewPlaces{\nu}{w}{\tau}{\phi}}, \m{\withConjInv{\tau}{k}}, \m{\IsAuxTwist{S}{Q}{\psi}{c}}}{recalled, Definition~\ref{d:recall}(6)}
\nt{\m{\IsPrimeTriangle{\nu}{w}{\tau}{\phi}{a}{t}{\mathrm{ty}}{h}{c}}}{prime triangles, Definition~\ref{d:recall}(6)}
\nt{\m{\ribetNormaliser{\gamma}{a}{y}{p}{\theta}{g}}, \m{\HasTriangleCost{Z}{\tau}{\phi}{I}{C}{b}}, \m{\IsTraceIdeal{\mathcal{J}}{u}{p}{\tau}{\phi}{k}{\varepsilon}}, \m{\IsRibetElement{\mathcal{J}}{\beta}{\pi}{c}}}{recalled, Definition~\ref{d:recall}(7)}
\nt{$R$, $R'$, $R_x$, $R'_x$, $R_i$, $T$, $T_0$, $T_x$, $u_x$, $u'_x$, $\varrho$, $\psi$, $\xi_e$, $\HT$, $\HT'$, $\HT_x$, $\HT_i$, $\HT^{\flat}$, $\HT_{\kappa,N}$, $\kappa$}{representations, lattices, characters and types, a weight ($T$ also a finite set of primes or places and a Hecke element)}
\nt{$\mathcal{S}$, $\mathcal{J}$, $\mathcal{X}$, $\varepsilon_k$, $\beta$, $\pi$, $\theta'$, $d'$, $f$, $\mathrm{id}$, $\mathrm{pr}_1$, $\mathrm{ty}$, $t_V$, $m_0$}{the objects of Ribet's lemma in Theorem~C ($\beta$ also a real bound)}
\nt{$\gamma$, $\gamma_x$, $\gamma_i$, $\phi$, $\phi_j$, $\phi_x$, $\delta$, $\delta_v$, $o_v$, $g$, $g_i$, $a$, $u'$, $s$, $s'$, $z$, $y$, $y_x$, $y_i$, $k$}{elements of Galois groups, scalars, signs, forms and adelic points, as said at each place ($\gamma$ also the function of (F2), $z$ also a law or Hecke data, $y$ also a number in Theorem~D, $k$ also a number or an adelic element)}
\nt{$x$, $x_i$, $w_x$, $\fP_x$, $W$, $U$, $J$, $K$, $I$, $Y$, $\mathfrak{p}$, $\mathcal{A}$, $Z$, $Z'$, $N'$}{Frobenius data, sets of them, an open set of Galois elements, index sets, a detecting set, spanning sets and subspaces ($x$ also an eigensystem, a pair \m{(w,i)} or a form)}
\nt{$n$, $r$, $r'$, $r_i$, $t$, $t_i$, $e$, $e'$, $\nu$, $h$, $h'$, $N$, $M$, $M_i$, $M'$, $j$, $n'$, $i$, $b$, $c$, $q$, $q'$, $d$, $m$, $a'$}{finite sets and natural numbers ($c$ also a family of labels or a cocycle, $b$ also Hecke data, a sign or a ring element, $q'$, $d'$, $e'$ also functions of eigenvalues, $m$ also the number in Lemma~4.5 of the first note)}
\nt{$H$, $\alpha$, $\eta_l$, $\sigma$, $c'$}{the multiplier and bound of a height; the budget of (F2) and the local bounds of a height; an automorphism of~$\CC$ or a complex embedding}
\nt{$B$, $E_s$, $G$, $\lambda$, $p_1$, $p'$, $\eta_a$, $y_a$, $y'_b$, $w'$, $L$, $L'$, $S$, $Y_0$, $C$, $p_0$, $b_0$, $b_e$, $X$, $X_k$, $\zeta$, $\Lambda$, $\epsilon$}{objects of the proofs: a base level, operators, scalars, linear forms and forms, linear maps, sets of places, the constants of Ribet's lemma, variables}
\nt{$\mathrm{U}(1)$, $\mathrm{U}(2z)$, $\CC$, $\mathbb{R}$, $\NN$}{unitary groups; the complex and real numbers; the natural numbers}
\nt{$s'$, $\gamma_0$, $j'$, $K$, $U'$, $T'$, $\gamma_i$, $n_i$}{in the outlines and in Input~6: the enlarged finite set of Theorem~C (also a scalar in Theorem~B); the inertial element of (W4); a second lift index; the objects of Input~1 of the second note}
\nt{$\Phi$, $\Sigma$, $\varphi$, $\varphi'$}{the set on which a law tests even classes; the set \m{\pd{\ED{D}}{\ell}\cup\placesAboveG{D}}; forms in Lemma~\ref{l:heckefacts}}
\end{multicols}
\endgroup

\head{References}
{\fontsize{9.4}{12.6}\selectfont
The references are those named by the second reading of this write-up, with the places given there; I did not
open them.\par
\smallskip
C.~Greither, \emph{Class groups of abelian fields, and the main conjecture}, Ann. Inst. Fourier 42 (1992).\par
B.~Mazur, A.~Wiles, \emph{Class fields of abelian extensions of~$\mathbb{Q}$}, Invent. Math. 76 (1984).\par}

\end{document}
